% =====================================================================
%  GENERATED; DO NOT EDIT BY HAND.  Built by
%  python3 ebook-method/moderne-filosoffer/tr_parts.py build --write
%  from tr_template.tex and the six parts in ebook-method/moderne-filosoffer/tr/*.texfrag.
%
%  Harald Høffding, MODERN PHILOSOPHERS.  English translation of
%  "Moderne Filosofer" (Copenhagen: Gyldendal, 1904), made 2026-10-14
%  from the collated transcription (transcription.tex in this folder) by
%  six translators working from GLOSSARY.md and TRANSLATOR-BRIEF.md
%  (ebook-method/moderne-filosoffer/).  The printed page numbers of the
%  1904 edition are in the margin.  Titles of books and quotations in other
%  languages are left as Høffding gives them.
% =====================================================================
\documentclass[12pt,a4paper,openany]{book}
\usepackage[utf8]{inputenc}
\usepackage[T1]{fontenc}
\usepackage{libertinus}
\usepackage{libertinust1math}
\usepackage{textalpha}
\usepackage{amsmath}
\usepackage[english]{babel}
\usepackage{setspace}
\usepackage{fancyhdr}
\setlength{\headheight}{28pt}
\addtolength{\topmargin}{-13.5pt}
\pagestyle{fancy}
\fancyhf{}
\fancyhead[LE,RO]{\thepage}
\fancyhead[RE]{\textit{Harald Høffding: Modern Philosophers}}
\fancyhead[LO]{\textit{\leftmark}}
\setstretch{1.3}
\newcommand{\whiteline}{\par\vspace{\baselineskip}\par}
\newcommand{\secrule}{\par\begin{center}\rule{2cm}{0.4pt}\end{center}\par}
% the headings (moderne-filosoffer): a group (FØRSTE GRUPPE. + its name), a chapter on one line (I. WILHELM WUNDT.)
% or with its numeral above (II. / ROBERTO ARDIGÒ.), and the lettered sections of the Guyau and Nietzsche chapters
\newcommand{\grouphead}[2]{\clearpage\begin{center}{\Large #1}\\[8pt]{\large #2}\end{center}\addcontentsline{toc}{part}{#1 #2}\bigskip}
\newcommand{\chapterhead}[2]{\bigskip\begin{center}{\large #1 #2}\end{center}\addcontentsline{toc}{chapter}{#1 #2}\markboth{#2}{#2}\medskip}
\newcommand{\chapterheadb}[2]{\bigskip\begin{center}{\large #1}\\[4pt]{\large #2}\end{center}\addcontentsline{toc}{chapter}{#1 #2}\markboth{#2}{#2}\medskip}
\newcommand{\sectionhead}[1]{\par\medskip\begin{center}#1\end{center}\addcontentsline{toc}{section}{#1}\smallskip}
\newcommand{\supplied}[1]{\textlangle#1\textrangle}
\usepackage{graphicx}
\DeclareUnicodeCharacter{0254}{\reflectbox{c}}   % ɔ: (id est), p. 103
\usepackage{hyperref}
\hypersetup{
  pdftitle={Modern Philosophers},
  pdfauthor={Harald Høffding},
  hidelinks
}
\usepackage{marginnote}
\newcommand{\opage}[1]{\marginnote{\footnotesize\textit{[#1]}}}
\newcommand{\apage}[1]{\marginnote{\footnotesize\textit{[#1]}}}

\begin{document}

\frontmatter
\begin{titlepage}
\centering
\vspace*{2cm}
{\Huge MODERN PHILOSOPHERS}\par
\vspace{1.2cm}
{\footnotesize BY}\par
\vspace{0.8cm}
{\large HARALD HØFFDING}\par
\vfill
{\large COPENHAGEN}\par
\smallskip
{\small GYLDENDAL BOOKSHOP}\par
{\footnotesize \textsc{Nordisk Forlag}}\par
\smallskip
{\footnotesize 1904}\par
\vspace{1.5cm}
{\small\textit{Translated from the Danish}}\par
\end{titlepage}
\chapter*{CONTENTS.}
\markboth{Contents}{Contents}
\begin{small}
\newcommand{\ix}[3]{\par\noindent\hspace*{#1}#2\dotfill #3}
\newcommand{\iz}[2]{\par\noindent\hspace*{#1}#2}
\hfill{\footnotesize Page}
\ix{0em}{\textit{Introduction}}{VII}
\par\medskip\centerline{FIRST GROUP.}
\ix{1em}{The Objective-Systematic Tendency}{1}
\iz{1em}{I. \textit{Wilhelm Wundt.}}
\ix{2em}{1. Course of Development}{1}
\ix{2em}{2. The Psychological Problem}{6}
\iz{3em}{a. Physiology and Psychology.}
\iz{3em}{b. The Peculiarity of Mental Life.}
\iz{3em}{c. The Elements of Mental Life.}
\ix{2em}{3. The Problem of Knowledge}{14}
\iz{3em}{a. Naive and Critical Realism.}
\iz{3em}{b. Observation, Understanding, Reason.}
\ix{2em}{4. The Problem of Existence}{19}
\iz{3em}{a. Metaphysics and Empirical Science.}
\iz{3em}{b. The Ideas.}
\iz{3em}{c. Critical Remarks.}
\ix{2em}{5. The Ethical Problem}{24}
\iz{3em}{a. History and Ethics.}
\iz{3em}{b. Society and Individual.}
\iz{3em}{c. Ethics and Metaphysics.}
\iz{1em}{II. \textit{Roberto Ardigò.}}
\ix{2em}{1. Italian Philosophy after the Renaissance}{29}
\ix{2em}{2. Ardigò's Development}{31}
\ix{2em}{3. Doctrine of Evolution and Theory of Knowledge}{35}
\ix{2em}{4. Psychology}{37}
\ix{2em}{5. Ethics}{40}
\iz{1em}{III. \textit{Francis Bradley.}}
\ix{2em}{1. Idealism in the Most Recent English Philosophy}{43}
\ix{2em}{2. Bradley and Philosophy}{44}
\ix{2em}{3. Appearance and Reality}{48}
\ix{2em}{4. Characterization and Critique}{51}
\iz{1em}{IV. \textit{Alfred Fouillée and French Philosophy in Our Day.}}
\ix{2em}{1. Introduction (Taine. Renan)}{55}
\ix{2em}{2. Fouillée's Philosophy of Evolution}{64}
\ix{2em}{3. Philosophy of Discontinuity}{68}
\ix{2em}{4. Charles Renouvier}{69}
\ix{2em}{5. Émile Boutroux}{77}
\par\medskip\centerline{SECOND GROUP.}
\ix{1em}{The Epistemological-Biological Tendency}{81}
\ix{1em}{Introduction}{81}
\iz{1em}{I. \textit{Philosophizing Natural Scientists.}}
\ix{2em}{1. James Clerk Maxwell}{82}
\ix{2em}{2. Ernst Mach}{87}
\ix{2em}{3. Heinrich Hertz}{92}
\ix{2em}{4. Wilhelm Ostwald}{95}
\iz{1em}{II. \textit{The Natural History of Problems.}}
\ix{2em}{Richard Avenarius}{98}
\par\medskip\centerline{THIRD GROUP.}
\ix{1em}{The Philosophy of Valuation}{107}
\ix{1em}{Introduction}{107}
\iz{1em}{I. \textit{Jean Marie Guyau.}}
\ix{2em}{1. Critique of English Ethics}{110}
\ix{2em}{2. Ethics}{111}
\ix{2em}{3. Aesthetics}{114}
\ix{2em}{4. Philosophy of Religion}{115}
\iz{1em}{II. \textit{Friedrich Nietzsche.}}
\ix{2em}{1. Characterization and Biography}{118}
\ix{2em}{2. Nietzsche's Writings}{126}
\ix{2em}{3. The Starting Point in the History of Culture}{131}
\ix{2em}{4. The Goal of History and the Social Dualism}{133}
\ix{2em}{5. The Slave Revolt in Morality}{138}
\ix{2em}{6. The Philosophical Grounding of the Social Dualism}{141}
\ix{2em}{7. The Last Yes --- and Zarathustra's Death}{143}
\ix{0em}{\textit{Notes}}{148}
\end{small}
\clearpage

% --- p. VII ---
\chapter*{INTRODUCTION.}
\addcontentsline{toc}{chapter}{Introduction}
\markboth{Introduction}{Introduction}
\opage{VII}In my presentation of the history of modern philosophy I stopped
at the year 1880. Various reasons led me to set this year as the boundary. In the first place, by this time
a certain provisional conclusion had been reached in the debate and interaction
between the two great tendencies of thought in the nineteenth century,
Romanticism and Positivism. Both had unfolded their consequences
and had in part found their correctives. In the characteristic standpoints of Lotze and Spencer
and in the critical philosophy now pressing forward anew
provisional resting places had been won. Alongside of these
new endeavors had indeed appeared, but they
had not yet taken shape as clear and characteristic tendencies.
In the second place, the presentation and valuation of immediately
contemporary tendencies of thought is subject to other conditions than the treatment
of philosophical phenomena that lie historically complete.
A psychological and biographical illumination does not become so
easily possible; and yet such an illumination is, after all, necessary
in the philosophical domain on account of the intimate reciprocal relation
between thinking and personality. The manner of treatment
must therefore be different from that toward tendencies that belong more decidedly
to the past. And in the third place, my own
more independent philosophical works begin at about 1880; this too
makes it less easy for me to stand objectively toward
the works of others from this time onward.

Now that I attempt a characterization of the philosophical endeavors
of the last quarter century, which I regard as
% --- p. VIII ---
\opage{VIII}the most significant, I am aware that the personal factor
will make itself felt more strongly here than in the earlier work.
It will be at work in the choice of representatives, in presentation, and in valuation.
This is the reason why I let
this book appear as an independent work, and not simply
as the third part of the history of modern philosophy. ---

If I am to single out at once, from the beginning, a peculiarity
of the philosophical attempts of the most recent time, it must be
this: that they are still harder to classify than the philosophy of earlier
times. The more thoroughly one studies philosophy, the more it
appears how little the usual rubrics, all the many
-isms, suffice. Both thought and existence have too great a manifoldness
in essence and activity for such an external ordering
to be possible. In particular it appears again and again how
a philosophical movement is both a work of thought and a sign of the times.
Philosophy is both treatment of problems and symptom,
and in our day it appears in a pronounced form in both
these ways. In part existence shows itself to us constantly from several sides,
in greater depth and variety; in part one becomes more and more
aware of the cooperation and accompaniment of subjective factors
in all thinking. In both respects --- one might call them
the objective and the subjective side of philosophy --- the philosophizing of
our time appears with a stronger accent than that of earlier
times. One stresses more strongly than before the necessity of observation,
analysis, critique, objective coherence, and one stresses
more strongly than before the subjective choice of points of view and starting points
and of concluding thoughts. And yet (this is perhaps
the most characteristic thing) one does not really believe that there
is an irreconcilable opposition present here; rather there is a quiet
conviction that the objective and the subjective movement of thought
will in the end coincide.

The philosophizing personalities are led, in their inquiry, now by more objective,
now by more subjective reasons. Wilhelm
Wundt and Ernst Mach have passed from natural science to
philosophy, and even natural scientists such as Clerk Maxwell and Heinrich
Hertz, who wish to remain within the bounds of their subject, feel
nevertheless a need to render account to themselves of the first presuppositions
% --- p. IX ---
\opage{IX}on which they build, and thereby come into contact with philosophy.
Roberto Ardigò is led from Catholic theology to critical and positivistic
philosophizing, without a sharp break, but in such a way that a long
work of thought imperceptibly leads him over to a peculiar
positivistic standpoint. Friedrich Nietzche begins as a philologist
and historian of culture, but through the problem of culture, as
it presented itself already within a narrower historical frame,
a passionate work of thought is awakened in him, which leads him
through violent criticism and polemic to an ideal construction,
in which fancy and feeling work together with thought. In men
such as Francis Bradley, Richard Avenarius and Jean Marie Guyau
the philosophical impulse seems to have awakened more directly, though
it too appears in them under individually different forms. ---
We all philosophize, really, though it is --- on account of
the various inner and outer circumstances --- only a few in whom
a coherent work of thought comes to development. ``Among
a hundred who ponder, one becomes a thinker.''

A division of the most important philosophical phenomena of the
last quarter century may be had by distinguishing between
three groups or tendencies.

\textit{The systematic tendency} is represented by a group of
thinkers who aim above all at giving an illumination of the problem of existence,
and thus work at the development of a coherent
world-view. Such men are Wundt, Ardigò,
Bradley and Fouillée. --- In another tendency, the \textit{epistemological-biological,}
the problem of knowledge stands in the first rank,
and the endeavor is directed to finding the simplest ways by
which the demands of the life of thought may be satisfied, in that the life of thought itself
is regarded as a special kind of life, following the general
laws of life. To this belong men such as the natural scientists Maxwell,
Hertz and Ernst Mach; this tendency receives its most characteristic execution
in Richard Avenarius's attempt to give a Natural History of
Problems. --- A third tendency keeps essentially
to the problem of valuation. This \textit{philosophy of valuation} takes up
the fundamental ethical and religious problems and seeks to illuminate or
solve them from new points of view. In the authors of this tendency,
Guyau and Nietzsche, the subjective factor predominates. ---

% --- p. X ---
\opage{X}Within all three groups psychological observation and
analysis play an essential role, so that the fourfold set of problems
which I found in ``Den nyere Filosofis Historie'' by the historical
route, and which I have sought in the essay ``Filosofiske Problemer''
to determine and ground more closely, is confirmed here too. ---

Through the characterization of the three groups of modern
thinkers named, I hope to be able to give an orientation in the thought-world of our time.

\mainmatter

% --- p. 1 ---
\grouphead{FIRST GROUP.}{THE OBJECTIVE-SYSTEMATIC TENDENCY.}

\opage{1}\chapterhead{I.}{WILHELM WUNDT.}

1. Wundt stands as a great example of the colossal working capacity of the Germans. There is no living philosopher who surpasses
him in fullness of knowledge, in the ability to bring his manifold
learning together under general points of view, and in tireless
striving for exactness and clarity in the particular as well
as in the large. He does not have Lotze's delicacy and inwardness;
he stands more broadly and solidly on the ground of reality and
does not venture so far out into metaphysics as Lotze. In comparison
with Spencer he is more of an idealist. From the broad foundation of experience
he strives upward toward the peak of thought which the philosophy of Romanticism
believed it had scaled, but which according to Wundt,
who is here inspired by Kant's critical spirit, is merely
a highest point in the field of vision, which we use for orientation.

Wundt was born near Mannheim on the 16th of August,
1832. He studied medicine, and it was through his physiological
studies that he came to philosophy. His \textit{Beiträge zur
Theorie der Sinneswahrnehmung} (1859---62) investigate above all
the influence of the eye movements on the perception of space. At the congress of natural scientists
in Speier in 1861 he gave a lecture on \textit{the physiological
time,} i.e. on the time that elapses in the apprehension and
reproduction of a sense impression. In \textit{Vorlesungen über die Menschen-
und Thierseele} (1863, 3rd ed. 1897) the
% --- p. 2 ---
\opage{2}ethnological method is applied alongside the experimental-physiological. After
having served for a number of years as extraordinary professor
of physiology at Heidelberg, he became in 1874 Albert Lange's successor
as professor of inductive philosophy at Zürich, but took over
already the following year a professorship of philosophy at Leipzig,
where he is still active. Wundt's life is the course of his studies.
His personality is not very prominent in his works;
to one who comes into contact with him in life, it appears
with mild warmth and open cordiality and with quiet intellectual
enthusiasm.

Shortly before Wundt went to Leipzig, he published his \textit{Physiologische
Psychologie} (1874), a principal work of experimental
psychology, next after Fechner's ``Elemente der Psychophysik.''
It presents above all the methods and the physiological foundation
of psychology, and only those parts of psychology that can in particular
be illuminated along this way are treated at length. Description
and analysis recede into the shade before the experiment. --- A more
evenly detailed presentation of psychology Wundt has later
given in his \textit{Grundriss der Psychologie} (1896). --- At Leipzig
he established an experimental-psychological laboratory, the first
university institution of that kind, and began the publication of the journal
\textit{Philosophische Studien} (20 volumes, from 1883---1903).

On the significance that experimental-psychological studies had
for Wundt in the transition from physiology to philosophy, he
has expressed himself in an interesting way in the essay \textit{Ueber
psychische Causalität und das Princip des psychophysischen
Parallelismus} (Philos. Studien 1894. p. 121---124). They convinced
him of the independence of psychology in relation to
physiology and metaphysics, which had already been asserted by the English
school and by Kant. ``If I were asked,'' he says, ``in what
the psychological value of experimental observation has
consisted for me and still consists, I would answer that it
has called forth in me, and ever more firmly established, an entirely new view
of the nature and connection of the mental processes. When I
first came face to face with psychological problems, I shared
the common prejudice, which lies so close at hand for the physiologist, that
the genesis of sense perceptions is merely a work of the physiological
properties of the
% --- p. 3 ---
\opage{3}sense organs. I learned first through the operation of the sense of sight
[especially as regards the perception of space] to understand the act
of creative synthesis, which gradually became a guide for me
to wresting from the development of the higher functions of understanding and imagination
too a psychological comprehension, such as the old psychology
had not been able to help me to. When I then went on to
the course of ideas in time, there opened for me
a new insight into the development of the functions of will [namely through the
significance that preparation and tension have for the shortening of the physiological
time], the development of the external functions of will
from the internal, of the compound from the simple, --- and at the same time an insight
into the close connection among all the mental functions,
which are otherwise separated by artificial abstractions and names
(idea, feeling and will), in other words into the
indivisibility and inner uniformity of mental life at all stages.''

Wundt had come by this way to the concept of synthesis
which Kant in his day had won through the analysis of the functions of understanding and of the perception of space
and time. By ``creative
synthesis'' he understands a bringing together and combining, whose
product acquires properties which none of the elements taken separately
had. Not seldom Wundt overlooks that this concept does not
give any explanation, however well it suits the description
of what is peculiar to mental processes. He stresses
in particular the moment of activity in synthesis, and is thereby led
gradually (in the later editions of his psychology) to designate
the will as the central thing in mental life.

Psychological studies were not, however, the only way by
which Wundt was led over from natural science to philosophy.
As a natural scientist he became aware that all science
of nature rests on certain presuppositions. The question arises
then involuntarily how these presuppositions, which ground
everything for us, can themselves be grounded. In this way epistemology
arises. In his work \textit{Die physikalischen Axiome
und ihre Beziehung zum Causalprincip.} Ein Capitel aus einer
Philosophie der Naturwissenschaften (Erlangen 1866) Wundt discussed
this question. He sought here to derive the fundamental principles of physics
from the proposition that everything has its ground, in
% --- p. 4 ---
\opage{4}connection with the second proposition, that all natural processes are
motions. This second proposition he grounds in turn by showing
that motion is the simplest change, in that a thing that
moves merely changes its position in relation to other things,
without needing to lose any of its properties. --- Wundt
is here not yet clear about the epistemological difficulties
in this conception, which in comparison with his
later epistemology has a somewhat dogmatic character.

In two ways, then, Wundt was led into philosophy, and
in his later writings he has further deepened and extended his concept
of philosophy.

The discussion of the problem of knowledge led him to work out
his large work on logic \textit{(Logik,} eine Untersuchung
der Principien der Erkenntniss und der Methoden wissenschaftlicher
Forschung. 1880-{}---83. 2nd ed. in 3 volumes 1893---95).
What gives this work its great value is above all the rich
material from the special sciences that is worked into it.
Wundt regards science rather as a historically given
fact in the life of the human spirit than as an endeavor
whose possibility and justification are to be investigated. This is
a trait characteristic of Wundt's whole philosophy. His logical
work thereby takes on rather the character of an encyclopedia than
of a discussion of questions of principle. In formal and in principle
respects the logical works of \textsc{Sigwart}, \textsc{Schuppe} and \textsc{Benno Erdmann}
are perhaps to be preferred; but none of them has
the real fullness of Wundt's.

Wundt's presentation of \textit{ethics} (1886. 3rd ed. 1902) has
a character that is analogous to the character of his logic. He
regards the ethical, like science, as a great fact,
which individuals are to acknowledge and continue, and he does not go
into the difficulties which the demand for this acknowledgment
contains, when a definite grounding of it is required. Ethics in Wundt rests on folk psychology: it shows
how the ethical fundamental principles spring from the historical
and social context within which every human life
is led. --- On the place of ethics in Wundt's whole philosophy I shall
speak later. ---

% --- p. 5 ---
\opage{5}Wundt's principal philosophical work is \textit{System der Philosophie}
(1889, 2nd ed. 1897). It is also the work in which his
art of presentation is at its greatest. He gives here his final conception
of philosophy as a whole and of its problems.
--- To this work is joined \textit{Einleitung in die Philosophie} (1901),
which, however, brings various nuances into the conception as compared with
the earlier writings. The largest place in this work is taken by
the characterization of the different tendencies in the history of philosophy.
---

What leads to philosophy is, according to Wundt, an intellectual
need, which appears more and more in its difference from
the religious need, without however needing to abolish the latter.
Philosophy means: a world-view and view of life that satisfies
both the demands of the understanding and those of feeling. In modern times
philosophy is no longer, as with the Greeks, alone in meeting
that intellectual need. The individual sciences have
developed and have divided experience among them. They can
be divided into three groups. Mathematics is the doctrine of number, space and
motion. Natural science is the doctrine of the natural processes
and of the objects of nature. The science of spirit is divided into psychology,
philology and history. The difference between these three groups consists
not merely in their different objects, but above all in
the different points of view they adopt. The development of all
these sciences has powerfully contributed to the growth of
intellectual interest. Philosophy cannot, as romantic speculation
believed, push them aside; but its task becomes
to combine the knowledge won by them in a manner free of contradiction,
so that none of the special points of view comes
to rule one-sidedly, and so that harmony is brought about
between the intellectual interests and the ethical needs.
Historically there can, according to Wundt, be demonstrated a course of development of
our concepts from ancient philosophy through the special sciences
back to philosophy.

Wundt divides philosophy into theory of knowledge and doctrine of principles.
Under the \textit{theory of knowledge} fall the doctrine of thinking
(formal logic), and the history and theory of knowledge (real
theory of knowledge). The \textit{doctrine of principles} has a general part, which
% --- p. 6 ---
\opage{6}Wundt calls metaphysics, and of which he remarks that
it is the conclusion of philosophy, not its beginning, --- and special
parts, namely philosophy of nature and philosophy of spirit.

According to this there should exist only two philosophical problems,
the problem of knowledge and the problem of existence. --- Ethics receives
its place in the doctrine of principles as one of the special
branches of the philosophy of spirit (coordinate with aesthetics and philosophy of religion), but represents
no independent treatment of problems. Yet there is
here a wavering in Wundt. In his presentation of the philosophical
tendencies (in ``Einleitung in die Philosophie'') he distinguishes expressly
among three problems, setting the ethical
problem alongside the epistemological and the metaphysical;
the epistemological problem has man as a
knowing being as its object, the ethical problem man as a willing being,
and the metaphysical problem concerns both sides
of his essence. And this also agrees with his conception of
philosophy as that which is to bring about agreement between
the intellectual interests and the ethical needs; for this
task presupposes the independence of the ethical problem.
--- Psychology has its place among the special sciences, alongside
philology and history, and is not reckoned as philosophy.
But within Wundt's own psychology there appears nevertheless clearly
a difference between the special, experimental investigations and
the points of view for mental life that become of significance for the ``philosophy of spirit,''
in particular for its general part, which is to give
``a coherent fundamental conception of spiritual being and
becoming.'' He admits furthermore that psychology receives an
exceptional position among the individual sciences in that it
enters into a more direct relation to the philosophical theory of knowledge
than any other science, since every act of knowledge
is an empirically given mental process, which as such comes
under the forum of psychology. From this it follows that it will be justified
to speak of a special psychological problem, --- and we
can thus apply the four philosophical problems mentioned earlier
to the division of our presentation of Wundt's philosophy.

2. It must be left to the history of experimental psychology
to present Wundt's significance for the methodology
% --- p. 7 ---
\opage{7}and organization of this discipline. Wundt stands here as the most significant inquirer
next after the founder of the field himself, Fechner. But whereas
Fechner (as Wundt has emphasized in his memorial address
on him) really had no pronounced psychological interest,
since he cared only for what concerned the threshold
of consciousness and thereby presumably could throw light on the relation
between the spiritual and the material, Wundt's psychological
interest is more comprehensive, and it is moreover at an entire
series of different points that psychology becomes for him
of significance for philosophy as a whole. I shall here dwell
on three such points, namely the relation between psychology
and physiology, the peculiarity of mental life and the elements of mental life.


a. The relation between psychology and physiology necessarily acquires
philosophical significance, since it becomes decisive for the determination
of the relation between the spiritual and the material.
Besides in his Psychology and in his System, Wundt has also treated
this question in his essay, already mentioned above,
``Ueber psychische Causalität und das Princip des psychophysischen
Parallelismus.''

In opposition to Fechner, Wundt does not regard psychophysical
parallelism as a complete solution of the problem.
The difference between the sciences rests, after all, more on differences
in points of view than on differences in objects, and one and the same
object can therefore occur in very different sciences.
Physiology and psychology regard the same object from different
points of view. But the fact that different points of view are necessary
does not prove that there is a difference between independent
objects present. And when we follow step by step the development of the
spiritual from lower to higher stages, we are compelled to
the assumption that this series of stages is prepared in the unconscious
nature, so that nature stands as the process of self-development of spirit.
Wundt thus finally holds a decided idealistic interpretation.
``Parallelism'' is for him only an auxiliary hypothesis, not a
definitive standpoint. But as a provisional hypothesis it is necessary,
because we cannot assume anything other than that there must be
uniformity between cause and effect, so that an action
% --- p. 8 ---
\opage{8}of the spiritual on the material, or conversely, would stand
as a miracle. He also lays weight in this connection
on the law of the conservation of physical energy, which requires physical
equivalents for every increase or decrease of physical
energy, whereby the usual doctrine of interaction is excluded.

Repeatedly Wundt has had to assure us that he steadily
holds fast to this view, even if he apparently expresses
himself in contradiction to it. He declares that, where
it is not a matter of questions of principle, he uses the
popular mode of expression with the same right as that with which a
Copernican lets the sun rise and set. I believe, however,
that the difficulty one finds here in Wundt does not merely
stem from the mode of expression. For there are two points where
he expressly makes exceptions to parallelism. In
the first place, the connection between the elements of consciousness
is not supposed to be capable of illumination from the physiological point of view; only the
individual elements, not their connection, have physiological
correlates. And in the second place, the value that
attaches to the mental phenomena is likewise not supposed to have such a correlate.
--- The first limitation of parallelism rests on an
unjustified application of the distinction between form and matter,
connection and element. The concept of element always designates merely
an approximation to a merely given, and it is therefore only in very
rough outline and purely relatively possible to distinguish between the connection
and the elements. Moreover the connection is just as much
a factual given as the elements are. The second limitation
rests on an unjustified distinction between the value and that
to which it is attached. That something is felt as valuable is a psychical
process, like any other. Or, if one will, value is just as
much a psychical element or a psychical quality as color
or tone is. It seems thus to be inconsistent
to make these exceptions. That Wundt has come to make
them is certainly due to his zeal to assert the independence of psychology.
In earlier times --- he holds --- this independence was especially threatened by encroachments from the spiritualist
side; but in our day it is the materialist encroachment that is especially
to be feared\textsuperscript{1}).

% --- p. 9 ---
\opage{9}A particular difficulty for a correct understanding of Wundt's theory
has been caused by the fact that he indeed assumes that the physical energy in the world
is constant, but speaks of a growth of energy in the spiritual
domain. Yet it is seen from the principal passages for this part of
his doctrine that by growth of spiritual energy he means partly the
arising of new psychical qualities, partly the genesis of new
values\textsuperscript{2}). It seems to me most correct here to speak
of concentration and organization of energy instead of
its growth. Psychical value presupposes such a concentration,
whereas it does not presuppose an increase of energy in
the whole. And on the other hand the principle of the conservation of
psychical energy is precisely necessary for understanding such
morbid states as are due to one-sidedness and divisions in
consciousness.

b. Through all stages and forms of expression of mental life there appear,
according to Wundt, properties that place it in opposition
to the material life of nature, although this is its external, sensibly
intuitable side. In his various presentations
Wundt enumerates these peculiarities in somewhat different ways.
But the most important traits are always the three that
will be discussed in what follows.

\textit{α.} Every mental content is a process, an activity, in
constant interaction with other processes and determined by antecedent
processes. The soul is no ``thing,'' no ``substance,''
but activity, actuality. This concept of the soul Wundt calls
the actual, in opposition to the substantial concept of the soul,
which both spiritualism and materialism, each in its own way,
hold\textsuperscript{3}). The essence of the soul reveals itself in the very continuous
connection of all mental occurrence. It is this
that makes it impossible to conceive the soul in analogy with a material
atom, as spiritualism has always been inclined to do.
At all points where we can observe mental life, we find
always continuity, the more exactly our investigations are carried out.
It will then be justified to assume that it holds also beyond
the domain where psychological observation is possible, in that
the unconscious is used as an auxiliary concept.

\textit{β.} Another peculiarity of mental life is its capacity
% --- p. 10 ---
\opage{10}by combining given elements to produce a qualitatively
new content. It is this peculiarity that may be designated
by the concept of creative synthesis. It manifests itself in every
sense perception, most clearly in the perception of space, in that this arises
through the cooperation of sensations of sight, touch and movement,
and in the timbre of tones, which arises through
the cooperation of overtones and undertones. It was first discovered
in the higher mental processes, in the genesis of images of fancy,
of concepts and of ideas. But it is a peculiarity
that appears in the most elementary psychical processes
as well as in the highest. From this side too the continuity
in mental life shows itself.

In the expression creative synthesis Wundt lays quite special
weight on the word creative. He does not bring out (as already
touched upon) the problem and the limit of understanding that
lie in this. For the qualitative new that comes in --- e.g.
when the form of space, according to the genetic theory, arises by synthesis
of heterogeneous sensations, or when a brilliant idea
springs forth on the basis of motives which it brings into entirely new
connection --- does not find its explanation by being called
the work of a creative synthesis. This expression may rightly be
used to describe, and it is well suited to stress the
new in the product in comparison with the factors; but it
leaves a large problem remaining and must therefore not be used
as a conclusion once and for all.

\textit{γ.} Very closely joined to the peculiarity just mentioned
is a third distinguishing mark of mental life. Nowhere does
mental activity manifest itself, according to Wundt, more clearly than in the manner
in which that which stands as object for consciousness is articulated,
distinguished into individual constituents by a process that may
be designated as relating analysis (beziehende Analyse).
This articulation takes place, as a rule, not in such a way that
the parts thereafter appear as independently existing units, but
in such a way that they continue to stand in relation to the whole and have
their significance through their definite place within it. As
an illuminating example Wundt uses the manner in which, within
the whole field of vision (Blickfeld), we are able to let a single point
% --- p. 11 ---
\opage{11}(Blickpunkt) stand out in particular, by letting the impression that
corresponds to it fall on the place of the most distinct vision on the retina.
The distinguishing attention that manifests itself in the
relating analysis Wundt calls apperception. By its
help the parts of an image of a whole step forth successively before
us, without losing connection with the whole. Thus
for the speaker the thought stands as a whole, before the individual parts are articulated;
and for the artist the work stands as a whole, before the
individual parts are distinguished and executed. It is a progressing clarity
and distinctness that is thereby attained.

c. Wundt first became aware of the active side in the life of consciousness
through the investigation of sensation and idea.
This investigation led him to the concept of apperception, which
I have already provisionally mentioned --- one of the most difficult concepts
in Wundt's philosophy.

Apperception is not merely an activity that manifests itself in the domain
of ideas; it manifests itself also in all feeling, in that
pleasure and displeasure are determined by the relation of the content of ideas to
the psychical activity, --- and it manifests itself in what in a narrower
sense is called the will. When Wundt as an illuminating example
uses the difference between field of vision (Blickfeld) and point of fixation
(Blickpunkt), it must nevertheless be noted that apperception
according to his theory is not merely distinguishing, as when a single
point is fixated within an entire context, but also combining,
as in all formation of concepts and judgments. It stands
in opposition to the predominantly passive and involuntary processes
--- thus among others to the mere association of ideas.
Wundt speaks against the psychological theory that would trace
everything in consciousness back to external association between independent
ideas, the so-called association-psychology.
He maintains that in all association an ``apperception'' is at work,
so that the distinction between association and apperception really
rests on ``an abstraction, to which reality can more or
less approximate.'' But it is not easy to bring all
of Wundt's statements into agreement with this. For Wundt speaks of
simple processes of association, in which the will is supposed to have
no influence at all; they are characterized ``durch die
% --- p. 12 ---
\opage{12}völlige Einflusslosigkeit des Willens auf die Art ihres Eintritts.''
And when he assumes that there is in the brain a special ``center of apperception,''
there seems to lie in this too that apperception
must be an independent process in relation to association
and other more passive processes\textsuperscript{4}). But the relation becomes
still more complicated in that within apperception, in turn,
a distinction is made between passive and active apperception. Active apperception
is characterized by being determined by an antecedent
idea and knowing more than one possibility, whereas
passive apperception is determined by a sense impression and
goes in one definite direction.

I have not been able to convince myself that there was any
necessity for introducing this concept into psychology. The concept of
attention contains in reality everything that experience
offers. As involuntary attention (corresponding
roughly to Wundt's passive apperception) it takes part
in every sensation (e.g. when a single point in the field of vision
is fixated) and in every association of ideas (in that it always
depends for an essential part on the prevailing feeling and interest
which idea is called up). As voluntary attention
(corresponding roughly to Wundt's active apperception)
it appears when, prior to the act of attention,
there is an expectation of or a striving to apprehend or
form an idea of something definite. It is then the task of psychology
to trace all special degrees and forms of attention,
from the most elementary to the very highest.

The essential thing in Wundt's theory of apperception consists in
its stressing the significance of the inner, central conditions
for psychical activity in contrast to the outer and peripheral. In particular
he stresses that the whole antecedent history of the development of the life of consciousness
determines the manner in which we self-actively
work over and apprehend sense impressions and ideas. As I
would prefer to put it: the degree and direction of attention
is determined by the need and striving of life at the given stage of development,
and this need and striving is in turn determined by the prehistory
of life.

The difficulties in Wundt's theory of apperception probably
% --- p. 13 ---
\opage{13}have to do with the circumstance that he has
been led to form the concept of apperception partly in the
experimental way, namely by experiments concerning prepared and
conscious attention, partly after historical models (especially
Leibniz and Kant), in whom pronouncedly conscious functions
likewise lay at the basis of the doctrine of apperception. Thereby the relation
between the clearly conscious and semi-conscious, voluntary
and involuntary functions of mental life has come to remain in a certain
obscurity, even after Wundt later on had extended the concept of apperception
to hold for the whole of mental life. The concept did
not come to undergo the changes that the greater extent necessarily
demanded. It is Wundt's striving, which appears clearly
through the series of works and editions, to give the concept of
activity, as expression of an activity and striving determined from within,
a prominent place in psychology. The concept of will
comes more and more to stand for him as the typical,
central concept, in analogy with which everything else in mental life
is to be thought. In the second edition of his Logik he designates his
psychological conception as voluntarism, a term which, by
his own statement, he has borrowed from \textsc{Friedrich Paulsen}, in whom
it is placed in opposition to the intellectualism of the older
psychology\textsuperscript{5}). Yet Wundt's psychology is not, even in its latest
forms (in the \textit{Grundriss,} in the third edition of \textit{Menschen- und Thierseele}
and in the fifth edition of \textit{Physiologische Psychologie)} presented
as a psychology of will, which follows the will in its development
from stage to stage and sees the other elements of consciousness
in their relation to it. Such attempts, attempts of no small
interest, have been made by \textsc{Friedrich Jodl}, by \textsc{Alfred Fouillée}
and by G. F. \textsc{Stout}. Indeed, Wundt does not even reckon the will among
the elements of the life of consciousness. He treats the phenomena of will
as the most compound and special forms of the life of consciousness,
and counts only sensations and feelings as psychical ``elements.''
This is understood only by the fact that Wundt began his
psychological studies with sensation and the intellectual
functions connected with it, and worked his way forward to seeing
the principled significance of psychical activity, without its having
succeeded for him in making this the all-conditioning principal element.
% --- p. 14 ---
\opage{14}The relation between the elementary and the compound psychical
phenomena always comes to stand obscurely with him\textsuperscript{6}).

The difficulties which thus present themselves in the study
of Wundt's psychological works bear witness, by this explanation,
precisely to his uninterrupted and tireless labor through 40
years, starting from the point at which he had come over from natural inquiry
to the domain of philosophy. On this road
he has, by his versatility and fullness of knowledge, thrown light on a
multitude of phenomena which a more rectilinear and sharply delimited
presentation perhaps would not have been able to.

3. As regards Wundt's treatment of the problem of knowledge,
I keep in this presentation mainly to his
\textit{System der Philosophie,} in which it seems to me that the
decisive points of view appear with the greatest clarity and exactness.


\textit{a.} Already in ordinary knowledge, which is not yet
scientifically formed, the most important processes of thought
appear distinctly. Knowledge is partly descriptive, partly narrative:
it aims partly at stating properties, partly at establishing
events. To description there corresponds, in strict thinking,
the demonstration of identity, to narrative the proof of
a logical relation of dependence. In the one appears the constant
or the static, in the other the changing or the dynamic. The relation of dependence
appears with full clarity in inference,
in the manner in which a proposition is derived from other propositions.
Therefore thought seeks to trace all relations in time and
space back to logical relations of dependence, so that they become of
the same kind as the relation that obtains in an inference between premises
and conclusion. One does not then remain standing at the
mere ascertaining of properties, but varies them in order to find
their mutual relations of dependence. This endeavor is
natural, since logical dependence is the only kind of dependence
that is immediately one with the lawfulness of thought itself;
the dependence here does not become an external one; the freedom of thought
and necessity are here immediately one. All real relations of dependence
are therefore sought to be reduced to logical relations of dependence.
This endeavor finds its expression in the
% --- p. 15 ---
\opage{15}principle that everything has its ground (the principle of sufficient reason). This principle
connects our various acts of thought with one another in the most intimate
manner.

Knowledge always begins with the conviction of the real validity
of ideas --- as \textit{naive realism.} No distinction is yet made
between knowledge itself and its object; such
a distinction is due to a later reflection. Philosophers
often overlook that reflection always designates a later stage,
not the original stage. Sufficient grounds are required
to give up the reality of the immediately given. Such grounds
arise above all from the opposition between the changing and the
constant. What has validity can only be that which
persists despite all change, --- and this is found by work of thought,
not by immediate apprehension. When the work of thought forms
concepts of objects that can replace the changing content of immediate
apprehension, naive realism passes successively over into
\textit{critical realism.} This is only a continuation of the process
that already begins within naive realism, when this
distinguishes between observation, memory and fancy. --- On the relation
between naive and critical realism Wundt has written an
excellent essay, found in the 12th and 13th volumes of ``Philosophische
Studien.''

On the road from naive to critical realism three stages are
to be distinguished: those of observation, of the knowledge of the understanding and of
the knowledge of reason. These three stages correspond respectively to the
standpoint of practical life, to that of the special sciences and to that of philosophy.

\textit{b.} Within \textit{observation} time and space appear as two
forms in which the ordering of the given takes place, --- as the constant
in opposition to the changing qualities,
--- and as expression of a steady work of thought. The laws of space and time
can be fixed in concepts, and thought therefore has a natural
tendency to hold to them and to regard the qualities
as something subjective in comparison with them.

\textit{The knowledge of the understanding} begins when the strict demand
for constant persistence leads away from immediate intuition
and replaces this with words and symbols. Only that which can
% --- p. 16 ---
\opage{16}be determined in concepts is here called object. We come here
necessarily into the domain of hypotheses, but it is logical
grounds that compel to this, in that thought strives to find
a coherence free of contradiction among all parts of the content of experience.
Everything here goes back to the laws of thought. But pure
thinking we have at no point, any more than pure experience. Thereby
apriorism is rejected, as well as empiricism. --- In
mathematics, natural science and psychology alike it is the knowledge of the understanding
that is at work. --- A good example is the natural-scientific
concept of matter or of material substance.
This concept is formed by tracing all material properties back
to relations in space, --- to motions and positions.
And this reduction is motivated precisely by the fact that only thereby
does it become possible in a simple manner and by virtue of the laws of thought itself
to derive one material change from another.
The empirical qualities are used in natural science only as
aids for inferring back to the relations of the objects in space
and time. It is not the task of a theory of matter to develop
ideas that resemble the phenomena of empirical bodies,
indeed not even to form intuitable ideas at all,
but to establish concepts by whose help empirical
phenomena can be derived from previously given phenomena. What
holds for the concept of matter holds also for concepts such as
inertia and energy. It is constantly the demand for connection free of contradiction
in accordance with the principle of sufficient reason
that lies at the basis of these formations of concepts. --- In the
psychological domain there is, according to Wundt, no reason to
undertake a construction of concepts in analogy with those that lead
to the concept of matter in natural science. For here observation shows
clearly that the most constant thing that is given in us is our
will, and that it is a never resting activity, a constant
occurrence; there then becomes neither possibility of nor reason for
forming a concept that corresponds to what position in space is
in the mode of operation of the material phenomena. (Cf. above on
Wundt's ``actual'' concept of the soul). The constructions and hypotheses
that we need in the psychological domain
concern partly the relation between the psychical elements that
% --- p. 17 ---
\opage{17}cooperate in the inner processes, partly the relation between the psychical
and the material. --- Although natural science traces all changes
back to changes in position and motion, it does not deny
that the things in the world have inner properties which
are not expressed by their external relations; it is merely not its
business to discuss them. But without such inner properties it would
be impossible to understand how life, especially mental life, has
been able to arise in the world. And everywhere that a series of events
appears as a whole, within which each member
stands at its definite place as cooperating toward a decisive
result, it is justified to apply the concept of end, even
if one does not think that the end exists in an idea
prior to the result. It is the manner in which the individual
member shows itself to be determined by its relation to and significance
for the whole that makes the teleological conception justified, without
its at any point excluding the purely mechanical. Only by
such a conception does it become, according to Wundt, justified to attribute
significance to the concept of development beyond the organic
domain.

\textit{The knowledge of reason} leads beyond experience, in that it seeks
an unlimited connection, whereas in experience all connection
is limited. The word reason Wundt takes in a narrower sense,
of the tendency to unity that leads to the formation of
totalities on the basis of the fragmentarily given. Reason
comes into being when consciousness awakens to the fact that the activity of thought,
in accordance with its own law, always goes further, in that every limitation
in grounding must be capable of being overstepped or at least thought
as overstepped. There then arises the Idea of a totality of all relations of dependence.
The Ideas of reason arise through a continuation
of the process that led the understanding to form its concepts.
Whereas the understanding merely sought to connect the given in a manner free of contradiction,
reason seeks --- by virtue of the same
laws of thought that guided the understanding --- to bring about a coherent
system.

This task thus presents itself to Wundt with the force of a necessity of
thought. In vain will empiricism and skepticism
deny the presence of the problems that here present themselves. Wundt
% --- p. 18 ---
\opage{18}admits that one cannot by way of logical proof refute
the empirical and skeptical standpoints; but he holds
that they are refuted by the very existence of thinking: for it is self-contradictory
to use thought to connect particulars with one another
and yet to reject the task of bringing the connections
thus won into mutual coherence. The principle of
sufficient reason leads ever further. The scientific significance of empiricism and skepticism
rests on their holding speculation
within certain bounds and seeing to it that it does not cause confusion
by being mixed up with the knowledge of the understanding.

The knowledge of reason can go beyond experience or become
``transcendent'' in two ways. It can continue the formations of series
that are begun within experience, in the same direction.
Thus in the series of numbers, in the dimensions of space, in the series of time,
in Spinoza's two attributes (spirit and matter). Here we think
as steadily continued continuity what in experience
is always given only as fragment. This overstepping
of experience Wundt calls \textit{real transcendence.} But the knowledge of reason
can also supplement experience by assuming
other series than those that appear in experience. Whereas
with real transcendence there is only a question of quantitative
infinity, here a qualitative infinity is introduced, in that
one assumes more and other sides or properties of existence
than the given ones. An example from the history of philosophy
is offered by Spinoza's assumption that existence has infinitely
many ``attributes,'' not merely the two that experience shows us (spirit
and matter). This kind of overstepping of experience Wundt calls
\textit{imaginary transcendence,} because it assumes new series of qualities
in existence, just as the imaginary numbers can be interpreted
as expression of new directions in relation to the direction that is
designated by the series of positive and negative real numbers,
--- ``sideways units,'' as a famous mathematician has called
them. A formation of hypotheses of this kind too is justified,
when it satisfies reason's striving for unity, and when it
does not violate the science of experience.

Only by these two ways can, according to Wundt, a general
world-view be attained. All metaphysics that is to have
% --- p. 19 ---
\opage{19}significance arises either by continuation or by supplementation of empirical knowledge.
By these two ways we are also led over from
the problem of knowledge to the problem of existence.

4. \textit{a.} Metaphysics begins really already within the
special sciences, namely in so far as these set up definitive
principles or hypotheses whose content cannot be shown empirically,
although they are indispensable for the connection of the content of experience.
Perhaps they cannot even be constructed completely
in analogy with objects that are given in experience. They
make experience comprehensible, but themselves lie beyond experience.
This holds e.g. of the natural-scientific hypotheses about
the constitution and fundamental properties of matter. Wundt lays
weight on retaining the expression metaphysics for such assumptions,
so that they shall not be confused with facts.
Every definitive hypothesis is metaphysical, and every metaphysics
is hypothetical.

The motive for metaphysics is the need for coherence and thereby for understanding of the content of experience,
--- in the end the need
for a coherent world-view. Each single domain of experience,
and as a consequence each single science, has a
certain tendency to regard itself as absolute, as a key
to the whole of existence. Therefore a constant critique is needed here\textsuperscript{7}).

Wundt's relation to positivism, criticism and romanticism
appears here in a characteristic manner. He approaches
positivism through the weight he lays on experience as
the standard. With great caution he goes step by step
toward the limiting questions of knowledge, and the manner in which
he acknowledges both the problems that present themselves here
and the difficulty of their solution shows him as an adherent
or continuer of the critical philosophy. But caution
rules in him only while he is still approaching
the limit. When he stands, or believes he stands, at the limit itself
he is not so wary, and he does not always think through
the final consequences of his concluding thoughts. He
here easily becomes heavy and dogmatic. The principle of unity captivates him,
and critical reflection recedes. One notices nothing
of the surf that naturally arises on the coasts of thought.
% --- p. 20 ---
\opage{20}On this point there is more to learn from the English philosopher
Francis Bradley, with whom we shall deal later. He
supplements Wundt in this respect, while he lacks
Wundt's broad empirical and experimental basis. --- Hitherto
we have, however, essentially come to know the empirical and the critical side
of Wundt; in the presentation of his metaphysical
standpoint the romantic element in his outlook will come to
light. ---

b. There are, as the history of metaphysical thinking
shows, three groups of Ideas, i.e. of concepts that by real
or imaginary transcendence bring a conclusion to the world-view.
Wundt calls them the cosmological, the psychological
and the ontological Ideas.

\textit{α.} The Idea of the physical world as a whole arises all
the more naturally within natural science, as it is really
only within closed wholes that one can demonstrate
definite laws. Every investigation presupposes a certain isolation
of that which is to be investigated; one must as far as possible think
of it as without influence from outside. Step by step one can
then perhaps extend the consideration. This holds not only when
the talk is of the physical world as a whole, but also when
it concerns the smallest parts of the physical world. Time and
space we can, by mere continuation, think extended beyond
experience; here there is only real transcendence present. But
every hypothesis about the physical world as a whole or about
its smallest parts becomes imaginarily transcendent, since at
every limit that is reached there is the possibility that something qualitatively
new might hold on the other side. Such
hypotheses too have, however, their justification, when they do not come into
conflict with empirical knowledge\textsuperscript{8}).

If one builds one's metaphysics on the basis of the cosmological
Ideas alone. --- that is, on concepts such as space,
time, matter, mechanics, --- this metaphysics takes on the stamp of
materialism.

\textit{β.} Wundt remarks that almost all strife in the metaphysical
domain concerns the psychological Ideas. We already know
Wundt's position on the psychological problem, in so far as
% --- p. 21 ---
\opage{21}this is placed on the ground of experience. Activity and striving for unity
are for him the characteristic marks of mental life, and the concept of soul-substance
contains for him, even where it is set up with spiritualistic
motivation, a concealed materialism. To him it is also
--- we may add here --- too individualistic, since it
isolates the individual mental beings from one another. The difference
between soul and body exists only in our conception. What
experience shows us is a spiritual organization, which in itself
is one with a bodily organization. But experience does not lead
to any conclusive concept here, and therefore every concept of soul,
the actual as well as the substantial, finally becomes imaginarily
transcendent. But when one wishes to conclude one's Ideas on
this question, one must according to Wundt not forget that the
individual mental life in our experience always appears as a member of
a community, of a spiritual whole, within which its motives,
its tendencies and its content alone become intelligible.
Wundt therefore lays great weight on folk psychology. As the
highest psychological Idea there appears then finally the thought of an
all-embracing ground of unity (ein Uebergeistiges), which lies at the basis
of all mental beings and of their mutual connection.
This Idea becomes imaginarily transcendent, since for us
all psychological determinations can be applied only to individual
beings, for which reason we lack in our experience expressions
for that which is to express the totality of all idea and will.

If one builds one's metaphysics on the psychological
Ideas alone, it takes on the stamp of idealism.

\textit{γ.} If we do not wish to follow the more or less one-sided ways
of materialism and idealism, but to strive for a more comprehensive
concept of existence, then we must according to Wundt combine
the cosmological and the psychological Ideas. Thereby we come
over to what he calls the ontological Ideas.

This combination can take place in the way that the cosmological
analysis is supplemented by the psychological. The cosmological
investigation shows us existence as a connection
of elements, of whose ultimate nature it can say nothing definite;
the psychological investigation has at least shown us the will
as our real essence. If we now combine both, we get
% --- p. 22 ---
\opage{22}the Idea of existence as a whole of striving and willing
beings.

One cannot, Wundt holds, like Spinoza remain standing
at the Idea of a being whose constitution is not more closely determined.
The question must always be asked whether the principle of unity in existence
does not agree with some one of our given concepts,
--- whether it lies nearer to the material or to the spiritual. Either
one must --- he maintains --- think the world as a material
unity or as a spiritual unity: a third is not given! ---
And here the choice is not doubtful for him. The only activity
that is immediately given to us is and remains our willing.
An infinite total will (Gesammtwille) becomes then the final, conclusive
Idea. This Idea, the Idea of God, is imaginarily transcendent,
and its content is indeterminable. By its help it becomes
nevertheless possible to apprehend the cosmic mechanism as the
outer husk for spiritual activity and striving, and our own
psychophysical existence stands before us as a world in miniature,
a microcosm. Thereby the demand of reason for unity is at once
satisfied, and we can see our human ideals
as consequences that spring from the ground of the world itself. The empty
and comfortless world-view which, supported by mere knowledge of the understanding,
sees the essence of things exhausted in their outer
relations and orderings, is abandoned.

\textit{c.} The definitive position which Wundt takes toward the limiting questions
naturally occasions some critical remarks, partly
with regard to the method, partly with regard to the motives for
the standpoint he ends by taking.

As regards the method, he himself holds that he is undertaking exclusively
such a continuation and supplementation of empirical knowledge
as the thought's need for unity demands. He rejects, on the contrary,
expressly the possibility that one should here be able to rely
on analogy: no analogy would suffice! --- But when he
is to ground the idealistic turn which he finally gives his
philosophy, he relies, surely, decidedly on the analogy between
microcosm and macrocosm. He asks which of our given
concepts the principle of unity in existence \textit{agrees with}
most nearly (übereinsstimmt). There can now here, he holds further, only
% --- p. 23 ---
\opage{23}be two possibilities: ``Either the world must be thought by us as a
material unity or as a spiritual unity, if it is at all
to be thought as a unity; --- a third is not given.''
(System.\textsuperscript{1} p. 411). It is correct that we know from experience only
spiritual and material phenomena; but we therefore have
(as Spinoza already saw) no right to assume that the
essence of existence is exhausted by them, and that a third is not given.
Wundt seems on this point to forget his own theory of
imaginary transcendence, which points to the fact that existence might
be more manifold in its essence than the metaphysicians suppose.
In any case it is certain that he relies on an analogy,
and that he even chooses between the two analogies that are, in his opinion, the only
possible ones. What determines his choice is the regard
to which of the two kinds of phenomena is most immediately
known to us. (Cf. System.\textsuperscript{1} p. 434: --- dem gleichend, das
wir in uns selber erleben). His line of thought here decidedly reminds
of Leibniz's and Lotze's, only that these with full consciousness rely
on an analogy. There is lacking then a closer investigation of the justification of
analogy in thinking, and of the various applications
of analogy that science knows\textsuperscript{9}).

As regards the motives, it is clear that Wundt's concluding
thought is not due merely to theoretical need. He declares
that the world-view to which we would be referred,
if we would not go the way he goes, is ``empty and comfortless.''
And purely from the standpoint of experience we stand, he says further,
with our ethical ideals, however valuable and indispensable they
may be, at an abyss over which no
bridge leads.


There is lacking here a series of investigations of a psychological and
ethical nature for a closer illumination of the need by which
the ontological Idea of unity should thus become practically necessary
and justified. In the first edition of Wundt's ``System''
there occurred the proposition that philosophy cannot indeed turn faith into
knowledge, but that it can prove the necessity of faith. This
proposition has fallen away in the second edition, perhaps because Wundt
has seen that the psychological necessity of faith cannot be proved.
--- There is lacking, moreover, a discussion of the difficulties which
% --- p. 24 ---
\opage{24}every monism, and especially a monism grounded on ethical motives,
has to contend with in the face of the disharmonies that experience presents\textsuperscript{10}). ---


At the conclusion of Wundt's treatment of the ethical problem
we shall return to some of the points touched on here.

5. On the somewhat obscure position of \textit{the ethical problem} in Wundt's
philosophy I have already spoken above. It finds in part
its explanation in the character of Wundt's ethics, in that this in part
tends to become one with folk psychology, in part issues in
a metaphysical or religious contemplation.

\textit{a.} Folk psychology is, according to Wundt, the vestibule of ethics. He
has in his old age begun a large and broadly planned work
which is to treat the most important subjects of folk psychology: language,
myth and custom (Sitte); only the first part is as yet available, and
a coherent discussion of Wundt's folk psychology
can therefore not yet be undertaken. But his main thought is that the
individual consciousness through language, religion and common habits of life
and habits of action is bound up with the life of the people, indeed with that of
all humanity. The individual will finds itself
as an element in a total will (Gesammtwille), by which
it is determined both in the motives by which it is led, and in
the ends toward which it strives. Culture and history
constitute a true common life and are not mere results of the coincidence
of countless individual endeavors. With great wrong does
the individualism from which the whole modern age suffers
regard the individual will as the only reality. There exists no original
isolated individual human being. Individualization proceeds
successively, from a state of social unity and without
any complete detachment from the total will.

The existence of society is the mightiest of all historical
facts. By it the individual is borne, even where he
believes he moves most independently; by it are determined sympathy
and piety, the foundation of social feeling. And participation
in the total spirit (Gesammtgeist) is greatest in the great spirits
of humanity, the leading spirits; they are able to draw
such nourishment from the total spirit that they can in turn assign it
new tasks, new directions.

% --- p. 25 ---
\opage{25}It would be Abderitism, if the value of the history
of humanity were to depend on whether the welfare of the individual or of
the particular group of men is promoted. Individuals and peoples
are perishable, and they are subject to passions, prejudices
and weaknesses. But the spirit of history is imperishable
and is always right. Historical development follows laws which
the individual and the people, standing in the midst of development,
cannot survey. Of the greatest importance here is the
fact that the effects of human actions always
more or less extend beyond the conscious motives
and intentions of individuals. In this is expressed a metamorphosis of ends
(die Heterogonie der Zwecke kalder Wundt det), which makes new
subjective motives possible, since the unforeseen effects can
call forth new feelings and impulses. The new formation of motives out
of given effects is the most important law of development that holds
for the ethical consciousness\textsuperscript{11}). It brings it about that we cannot become
conscious of the highest ends of development; the direction
in which they lie we can indeed surmise, but the conviction
that all stages serve a constantly progressing development
can in the end be only a faith, no knowledge. ---

The difficulty that lies in this for a scientific
ethics Wundt does not sufficiently stress. It seems clear, however,
that the possibility of new formations must make all ethics more
empirical than Wundt with his appeal to the spirit of history will
admit. We stand here again at a point where he reminds
of the philosophy of Romanticism. And he avows openly
that his ethics in certain fundamental thoughts stands close to speculative
idealism. Already Hegel had acknowledged a real moral
power in the total will. Hegel's error was, according to Wundt, that he merely
regarded the total will as objective moral power, whereas the
individual will was merely to be the unconscious bearer and
executor. In his doctrine of the relation between individual and society
Wundt seeks to remedy this one-sidedness.

\textit{b.} The individual is indeed borne by society, but works back upon it
by his own direction of thought and will. The individual consciousness
is creative, the social consciousness preserving. The new
comes from individuals, but society makes it fruitful
% --- p. 26 ---
\opage{26}for a
later development and thereby works for the continuity of
spiritual life. Not all are productive here to the same degree.
It is only the leading spirits who decisively determine the direction
of the total will, and most of all this holds of an ethical
genius such as the spirit of history perhaps produces only once
in centuries or in millennia, and who carries the ethical
life further by awakening hitherto slumbering impulses to life.

When Wundt thus, in opposition to Hegel, places productivity
in the individual persons, he comes into contradiction
with his own strong emphasis on the ``total will.'' Consistently
he would have to attribute to individualism a far greater significance
than he does, for the development of powerful individualities
must be a principal ethical end. In this respect
he could have learned much from the English school, on which he looks
down because of its empiricism. And not only as starting points,
but also as end points, as ends, he would have had
consistently to place the individual persons higher than he does.
``The individual,'' he says, ``be he ever so richly endowed by fortune
and ever so perfect, --- is yet a drop in the ocean of life.
What can his happiness and his pain mean for
the world?'' --- I cannot answer this question. But I
am unable to connect any meaning with words such as end
or value, when they are not illuminated at all by a relation to the
conditions of life to which beings with the capacity to feel pleasure and pain
are subject. Wundt's ethics ends in a mystical dualism,
in that the ends lie in the ``total will,'' the means in the ``individual wills.''
And for the tragic conflicts that can arise from the collision
between the individual will and a historically formed total will
he has no place in his ethics\textsuperscript{12}).

Characteristic of this side of Wundt's ethics is also
the proposition that justice is not an individual but a
public virtue, because it presupposes power to establish right and
to impose duty. He overlooks (perhaps under the influence of German
bureaucratism) that every individual through his position in the family,
society and state has a certain power, and that in his conviction,
in his judgments about others and in his conduct toward them
he has the opportunity to show justice. Even the bondsman has it in
% --- p. 27 ---
\opage{27}relation to the master. The thought of power has made the philosopher blind to
the significance of the many independent starting points for
judging and acting.

And yet several utterances show that Wundt has an eye for the
defects of the present social order. He complains that
the thief is judged more severely than the spendthrift, the usurer
and the gambler. He finds that the existing relation
between possession and labor produces two opposite occasions
for immorality: possession without a calling breeds craving for pleasure,
possessionlessness without a calling breeds envy. These evils
can only be remedied by a new legal order. Our present
condition of society suffers from the opposition between an antiquated
conception of right and new elements of culture that cannot be brought
under the old concepts.

\textit{c.} Wundt asserts the independence of ethics in relation to speculation
and metaphysics, precisely for the reason that it furnishes one of the
most important contributions to the foundation of a general world-view.
And as we have seen, ethics also finally dissolves
into speculation or faith. Indeed, ethics needs, according to
Wundt, more than other domains a metaphysical supplementation.
By virtue of this conception (which stands in contradiction with the assumption
of the independence of ethics in relation to metaphysics) ethics
for Wundt passes over into philosophy of religion. When the ideals
go beyond what can be attained by human striving, they acquire
a religious character. Philosophy is here able only to give
indeterminate hints; the positive religions give, on the other hand,
concrete symbols. The main religious thought is ``the demand
that all spiritual works (geistige Schöpfungen) possess
an absolute, or imperishable value.'' The more religion
raises itself above the standpoint of nature religion, the more it agrees
with science; a God who works miracles is a
God of nature, not the God of ethical religion. The development of Christianity
did not avoid relapses into nature religion. But the final task
of Christianity must --- in accordance with the statements of its founder
himself --- be the overcoming of all such elements
in religious faith as hinder the ethical content of the religious Ideas.
Christ will, precisely when one conceives him not
% --- p. 28 ---
\opage{28}as God, but as man, continue to stand as an ethical model,
and at the same time as a witness to the infinite, unfathomable
ground and end of the world, identical with the ethical ideal. ---

Wundt's philosophizing, as I have now sought to characterize
it, may be said to be typical of the thinking of our time.
Where it is in need of correctives, it itself indicates the points of view
and methods that must be applied to carry it further.
Especially is this the case where he wishes to ``supplement'' the science of experience
with metaphysical and religious elements. However energetically
he seeks to hold fast to and carry through a purely objective
standpoint, subjective factors nevertheless make themselves felt, which
call for a more definite investigation. Wundt's objectivism, which
at certain points turns into mysticism, is both his strength and
his limitation.

% --- p. 29 ---
\chapterheadb{II.}{ROBERTO ARDIGÒ.}

\opage{29}1. In my presentation of the history of modern philosophy I have had occasion to speak of Italian philosophy only in connection with the Renaissance. It was only in the Renaissance that Italian thought appeared with an originality and a boldness that could give it significance for the general development of philosophy. The spirit that in antiquity had animated the Pythagoreans and the Eleatics, and that had inspired Lucretius to his mighty didactic poem, came to life again in Pomponazzi and Telesio, in Bruno and Galileo. But the Italian movement of the Renaissance was struck to the ground after Bruno's pyre and Galileo's forced perjury. In the following centuries there are only a few names of philosophical interest to mention, such as \textsc{Giambattista Vico} († 1744), the precursor of modern sociology. At the end of the eighteenth century French philosophy acquired great influence, and around the psychologist and sociologist \textsc{Romagnosi} († 1835) a group of enthusiastic disciples gathered. But toward the middle of the nineteenth century the philosophical movement took on a different character, as religious, national, and philosophical enthusiasm entered into a close connection. The aspirations for freedom and unity in Italy stirred powerfully among a throng of young clergymen, who dreamed of a great harmony between religion and thought and between Church and State, and who believed in the possibility that the Church would put itself at the head of the
% --- p. 30 ---
\opage{30}national resurgence. Just as in the time of the Renaissance the enthusiasm for Italy's freedom lay behind Machiavelli's political philosophy, so it now animated the speculations of \textsc{Rosmini} († 1855) and \textsc{Gioberti} († 1852). Their philosophy was a kind of Platonism, borne by faith in an eternal truth exalted above all experience. For them philosophy was partly an introduction to religion, a kind of doctrine of the Logos, partly a means for patriotism. The difference between the two modern Platonists rests chiefly on this, that Gioberti assumes an immediate intuition of ideal truth, whereas Rosmini maintains (in a manner similar to Schelling in his later doctrine) that thinking leads only to a system of possibilities, and to that extent stands nearer to critical philosophy. \textsc{Mamiani} († 1885) developed Italian idealism in a direction that gave experience greater influence, and in this mediating form it ruled for a time at the Italian universities.


The rhythm that so often shows itself in the development of philosophical ideas expressed itself in a characteristic way in the role that a positivist tendency came to play in Italy in the last generation. Various causes worked together in this change in philosophical interests. Italy had attained its unity and freedom, so that there was room for special scientific interests; ideal enthusiasm could be replaced by real work. At the same time there came a strong influence from recent French and English philosophy through the writings of Comte and Mill; \textsc{Villari} applied the fundamental thoughts of positivism to the conception of history, \textsc{Angiulli} especially to the field of psychology and pedagogy. To this was added the influence of modern natural science. And on the other hand the Catholic Church now set itself more and more sharply against every philosophy that did not strictly follow in the footsteps of medieval thought. The disciples of Rosmini and Gioberti within the clergy were persecuted. It was even declared that the task now was to lead civilization back to the Catholic conception from which it had strayed in the last three centuries! In 1864 Pope Pius IX issued a ``Syllabus,'' a list of the errors of the age, among which are counted freedom of conscience, naturalism, and rationalism, and quite
% --- p. 31 ---
\opage{31}particularly the assumption that the method and principles of scholastic theology are not in accord with the demands of our time and the results of science. In an encyclical of 1879 Pope Leo XIII later declared that the philosophy of Thomas Aquinas' was a divine aid, which ought to be laid as the foundation by all teachers and used for the refutation of modern errors. Later he specially condemned critical philosophy in a circular letter to the French bishops (1899)\textsuperscript{13}).

Thus not only were the motives for the intellectual harmony that the Italian philosophers in the middle of the century had hoped for and striven toward fallen away, but the opposition between the Church and science had been sharpened, partly by the development of thought, partly by the Church's attempt to turn back the clock.

Under these conditions \textsc{Roberto Ardigò} developed, in the solitude of the cloister, from a childlike believing and patriotically enthusiastic Catholic into an energetic positivist. His development and his standpoint have interest not merely as a symptom of the times, but are of importance for the treatment of problems through his power of thought and his psychological talent, as well as through the peculiar change that positivism undergoes in his hands.

2. Ardigò was born near Cremona in the year 1828. His father, who was a farmer, moved to Mantua so that his son could study. His mother was a pious Catholic, whom Ardigò often mentions in his writings with heartfelt gratitude. It was owing to her influence that he became a priest. After the death of his parents Bishop Martini took care of him and made him a canon of the cathedral in Mantua. In his solitude he studied scholastic philosophy and modern science with zeal, deeply convinced that the ``modern errors'' must be capable of refutation. But slowly there arose --- more clearly to those around him than to himself --- an entirely new structure of thought in his inner being, and it then dawned on him at last, at a definite moment, that he no longer subscribed to the Church's dualistic doctrine, but believed in a great continuity in all things. In one of his writings \textit{(La morale dei positivisti} II, 3, 2, in a chapter that deals with the possibility of morality without religion) he expresses himself in an interesting way about this breakthrough.

% --- p. 32 ---
\opage{32}``By the example and teaching of my mother, a simple and poor peasant woman, religious faith and religious life arose and were nourished in my soul. Even to this day I cannot think of the exalted simplicity of my mother's religious feeling without the greatest enthusiasm and the tenderest emotion being awakened in me thereby. The faithful image of this feeling still exists in my sister, who owes her development neither to school nor to cultivated circles, and I feel, in the face of the way in which it appears in her, the greatest reverence, as for a sanctuary. This childlike religiosity of mine was later greatly strengthened by my life together with Mgr. Martini, which extended over more than twenty years. He had taken me to himself after the death of my parents; he gave me the bread I lacked, and made the scientific career possible for me ... He, who regarded science and sincerity of mind as the foundation and soul of religion, loved me because he saw me enthusiastic for study and found in me a character without guile. He hoped that I should one day be of use to the Church against the ignorant, superstitious, bigoted, and hypocritical kind of religion that he abhorred .... I devoted myself (in addition to the studies of the natural sciences and of philosophy, which I have never given up) with my whole soul to theology, especially the dogmatic and apologetic. I collected a library of Church Fathers and theologians, spending my best years of youth on the study of them, especially of St. Thomas. At last I wrote and published a book on confession against the Protestants. But the outcome of my studies turned out entirely opposite to what I was working toward and expecting. Little by little the doubt, which had already arisen on all sides in my earlier years, but which I had combated with unceasing reflection and study and for a long time believed I had overcome, came at last to stand without an adversary and appeared one fine day, to my greatest astonishment, as a decided conviction and indisputable certainty. Strange! Right up to that day I had devoted myself to the endeavor to stand firm in my old religious faith, and yet in my inner being, without my knowledge, the positivist system had unfolded completely beneath the system of religious
% --- p. 33 ---
\opage{33}ideas, which were the fruits of so much labor and of so long a time. The new system I found, to my greatest amazement, already fully finished and immovably fixed in my mind at the very moment in which I --- sitting on a stone under a bush in the little garden I had laid out by my canon's residence --- noticed how the last reflection broke asunder the last thread that still held me bound to faith. It then suddenly seemed to me that I had never in all my life believed, and that I had never done anything but strive to develop the purely scientific tendency in me. And this sprang, I believe, from the zeal with which I had from first to last sought, as far as possible, to know all the grounds that were opposed to religion, in order to be able to believe with a good conscience and to protect faith against all attacks.'' --- Ardigò adds that the step he now took was bound up with pain and melancholy on account of the memory of his mother and on account of his relation to his benefactor. But the will was strengthened in this struggle, after which religion came to stand as a poetic memory. And he found that moral ideality was not weakened; on the contrary, it stood from then on even more firmly for him that the true goods are those that are acquired by following the law of conscience in useful activity and by devoting oneself to the work of thought.

There now came some straitened years for Ardigò. He had to earn his living by teaching in schools, and only after several years did a liberal minister of education appoint him professor of philosophy at Padua (1881). His tendency was opposed not only to theology, but also to the tendency that under Mamiani's influence prevailed at the Italian universities. His teaching has exerted great influence and aroused great enthusiasm, which found public expression in a Festschrift on the occasion of his seventieth birthday (1898).

Ardigò's main problem is closely connected with his personal development. Already in a speech he gave in 1869 on Pietro Pomponazzi, in whom he sees a predecessor, he depicts thought as a force that comes into being without being noticed, and that, when it has matured, breaks forth in an irresistible manner.
% --- p. 34 ---
\opage{34}That was what he himself had experienced, and it was to him an example of how all development proceeds. In his inaugural address at Padua (1881) he alludes to his own course of development as that which had set him a problem that recurs on a larger scale, just as the motion of a molecule may remind one of the rotation of a celestial body. All development consists in a transition from the indeterminate to the determinate, yet in such a way that continuity is preserved: the whole, which appears in a more indeterminate manner at the first stage, constantly lies at the basis as the supporting foundation for the differences that successively appear. The first scientific task he set himself was that of following the natural development of human thought (la formazione naturale del pensiero). This plan was not carried out, but the material he had collected is to be found used here and there in his various writings. But the starting point became decisive for his standpoint.

He maintains that he reached his result independently of French and English positivism, which he only later came to know. When his conception of the concept of development recalls Spencer's, there nevertheless remains the difference that whereas Spencer builds essentially on the analogy with biological evolution, Ardigò builds essentially on the analogy with the development of thought, this most wonderful of all natural formations (la piû miragliosa delle formazione naturale). Ardigò calls himself a positivist. But what is essential in positivism is for him the starting point in experience, not the conclusion. The positivist is in no hurry to conclude; he does not wish to form an idea that can be used as a party banner, but goes forward step by step, as the truth opens itself to him. The point is to keep the horizon open. This line of thought appears particularly in Ardigò in the face of the problem of the unity of existence, in opposition partly to materialism, partly to the inclination to remain standing at a multitude of scattered facts. In his most recent work \textit{(L'unità della coscienza} 1898) he expresses himself particularly in this direction, as regards the psychological problem. One of his disciples, \textsc{Marchesini}, has in a special work, which he has called ``The Crisis of Positivism'' \textit{(La crisi del positivismo.}
% --- p. 35 ---
\opage{35}Torino 1898), discussed from this point of view the relation between Ardigò's philosophy and positivism. A certain similarity between Ardigò's standpoint and Wundt's emerges at this point. Already
Villari and Angiulli, moreover, stood in a similar critical position toward the older forms of positivism.

3. The fundamental features that Ardigò had found in the development of the life of thought he rediscovers in all natural development. The development of thought is for him an example of a law of the world. In his work \textit{La formazione naturale nel fatto del sistema solare} (1877) he has set forth this law more fully through an analysis of the well-known hypothesis of Kant and Laplace, which he regards as an example of a scientific explanation. Properly the work mentioned was to have been only a single chapter of the projected work on the development of human ideas.

According to the hypothesis mentioned, the present state of the solar system is due to a process of distinction (\textit{distinzione}), in that smaller parts or wholes have formed within the great, indistinct mass. But the whole is not thereby abolished; the whole --- the Indistinct (l'Indistinto) --- continues to exist, and only thereby does it become intelligible how there can be interaction between the parts that have been distinguished out (the globes): these belong together now just as much as before they were distinguished out. The existence of the Indistinct is the ground of solidarity. --- If one asks how the differences could arise out of the Indistinct, Ardigò answers that they were already present in the initial state as possibilities or in a latent state (forza latente, o virtuale); the special form is reached through successive unfolding. He concedes that this explanation is merely an inference backward from experience. We cannot know beforehand which special forms and parts will develop out of the indistinct initial state. Only observation and experiment can show them. Neither mathematics nor metaphysics can construct nature. And the reason for this is that the indistinct whole does not specialize ``of itself,'' but requires external conditions, unforeseeable by us and \textit{to that extent} ``accidental,'' in order to obtain its special forms and parts. But what is thus successively unfolded is nevertheless only the equivalent of what already existed as
% --- p. 36 ---
\opage{36}coexisting in the indistinct state. In the world-substance the work performed by the preceding causes is stored up, and thereby the possibility of future effects is conditioned\textsuperscript{14}). The continuity that thus appears between l'Indistinto and the formed parts (distinti) is what we express by the word \textit{Nature.} (La natura è la continuità di una cosa con tutti le altre. \textit{Form. nat.} p. 205).

Ardigò sees clearly, however, that one must ask in what those possibilities or ``latent forces'' properly consist. Possibility is for him only a reality of a different character from the given, and it consists in an activity: l'essere è attività \textit{(L'unità della coscienza} p. 479). Continuity is a constant energy. Therefore every indistinct whole, every Indistinto, points back to a still more comprehensive whole, out of which it has distinguished itself. The difference between indistinto and distinto, as well as that between possibility and reality, is therefore always only relative. A conclusion is not possible here. We enter upon an infinite series. But the last word of science concerns the fundamental relation between distinto and indistinto and states that all differences, wherever they present themselves, spring from a whole and are comprised by a whole. ---

Ardigò's doctrine of evolution and his theory of knowledge are closely connected. Every explanation is a distinguishing, an analysis. What continues to stand indistinct stands as unexplained. Thought is constantly driven further, although it has an inclination to stop at a distinto finito; --- and it is the very infinity of nature that thus drives thought forward. Indeed, it is the infinite nature that is the very energy in the proper law of logical thinking\textsuperscript{15}).

That the law of thought or of explanation is thus the same as the law of nature cannot be surprising, since thought itself is nature or a natural formation, like all other things. But when thought is a fact of experience like any other, and is itself only an example of natural formation, it is impossible to explain the whole of nature by deriving it from thought, as metaphysics and theology would have it. Thought itself is understood, after all, only through the general law of development from
% --- p. 37 ---
\opage{37}indistinto to distinti. We are directed to explain the one fact by the other, without being able to obtain any absolute conclusion. Nature is an immense circle whose center is everywhere and whose circumference is nowhere. ---

Ardigò has let the theory of knowledge be absorbed to too high a degree by the doctrine of evolution. Considered psychologically or biologically, the development of thought is indeed only a single example of the general laws of development. But the problem of knowledge presents itself as soon as one asks how we ground the validity of the general laws that we believe we find for thought as well as for other phenomena. This problem does not fall away through a mere reference to the law of development as common to all phenomena. For there is this peculiarity about thought, that it is by and through it that we apprehend and express all development and all law, both its own and that of other things. And the question is then what validity this apprehension and form of expression has\textsuperscript{16}).

4. Ardigò has labored with zeal for the independence of \textit{psychology} as a science of experience, and in this his talent for psychological description and analysis and his knowledge of natural science stood him in good stead. His psychological interest and gifts go in two directions. On the one hand he seeks to find a greater connection in mental life than appears to preliminary observation; he strives to find continuity between conscious states among themselves and between conscious and unconscious states. On the other hand he seeks to demonstrate finer nuances and differences than ordinary observation knows; he asserts, for example, that apparently uniform mental states in reality consist of rhythmic motions. Some of his most eminent disciples, such as the physiologist Giulio Fano and the criminologist Enrico Ferri, have in the Festschrift for the old philosopher dwelt especially on these peculiarities of his teaching. These are indeed the two directions in which all inquiry goes and must go; but Ardigò has --- as also his doctrine of the Indistinto and Distinti and of the relation between them shows --- had a special eye for them and sought to unite them.

His chief psychological works are \textit{La psicologia come}
% --- p. 38 ---
\opage{38}\textit{scienza positiva} (1870) and the work that he himself has designated as his philosophical testament: \textit{L'unità della coscienza} (1898).

For Ardigò the task of psychology is the study of our inner states, whose common peculiarity is expressed by the concept of soul. It is a subjective fact from which psychology proceeds; but it must be illuminated along an objective, physiological path. This double illumination shows us the difference of the phenomena of the soul from all other phenomena. However great the importance that physiological investigations have for psychology, physiology can nevertheless never take the place of psychology. The relation is this, that the psychical and physiological phenomena are expressions of one and the same ``substance,'' the psychophysical substance, or, as Ardigò prefers to call it, the psychophysical reality (realtà psicofisica). In his later writings (where his general doctrine of evolution is set forth) the psychophysical reality is one with l'indistinto\textsuperscript{17}). It is a natural totality that precedes and underlies the differences that we designate as soul and matter. The concept of the Indistinto expresses here, as everywhere, unity and solidarity. Both materialism and spiritualism rest on mere abstractions. Physiologically viewed, mental activity is connected with the physico-chemical processes. Psychologically viewed, the material, just as much as the mental, is given to us only as sensation and idea; both motion and thought we know only as psychical acts. Instead of asking how matter can pass over into soul, it would be more correct to ask how our originally more indeterminate ideas can gradually be differentiated, so that some of them come to stand as expression of an I, others as expression of a not-I. The popular, materialistic conception arises, according to Ardigò, from the fact that one does not realize that properties such as extension and motion are owing to psychical functions just as much as color and smell and the other so-called secondary properties.

The idea of the psychophysical reality gives no explanation, however --- this Ardigò emphatically insists. Its significance is only to assert a connection that we in our abstractions are inclined to dissolve. He compares it with
% --- p. 39 ---
\opage{39}the idea of gravitation, which likewise gives no explanation. It must be left to the future whether it can find a real explanation. The positivist is in no hurry to conclude. Further than a provisional concept we cannot come. The inner nature of the relation we cannot determine. The riddle that remains does not, however, stand alone for Ardigò. How a thought can be the equivalent of a motion is for him no greater riddle than how a billiard ball can set another in motion by a blow. Everywhere in nature we know the relations of simultaneity and succession only by experience, not through insight into the inner essence of the terms. --- By this last remark, however, Ardigò overlooks the fact that science, wherever it is possible, strives to transpose the external connection between the different phenomena into a connection of homogeneous terms, and that thereby a greater degree of continuity is attained. The greater the resistance this endeavor meets, the more enigmatic must the phenomenon be said to be at which we stop. ---

In his philosophical testament, the work on the unity of consciousness, Ardigò expresses his admiration for Kant's brilliant intimation in the doctrine of the unity of consciousness. True, he erred in assuming too great an opposition between the matter of knowledge and its form. But he grasped the fundamental scientific thought of psychology. There goes on in the life of consciousness an uninterrupted process of composition and combination, whereby all inherited or acquired dispositions and all new elements work together in determinate directions. There is from first to last a solidarity among all psychophysical functions, a constant tendency to unite in one stream. This tendency to ``mental confluence'' (confluenza mentale) underlies, among other things, the association of ideas. Often it is discovered only when one takes account of unconscious starting points or intermediate links, and it is especially this that makes the aid of physiology so indispensable to psychology. The unity that is the mark of the life of consciousness cannot be explained as a mere product of the cooperation of manifold elements, for we discover these elements only by a distinction, which always presupposes a previously given whole.
% --- p. 40 ---
\opage{40}And it is from the observation of the life of consciousness that we derive our concepts of unity and plurality at all. Everywhere that we speak of unity and manifoldness in the world, we have abstracted these concepts and the idea of their connection from the life of consciousness. Microcosm and macrocosm illuminate each other reciprocally.

5. Ardigò's \textit{Ethics} is borne by the same fundamental thoughts that characterize the other parts of his philosophy. The individual develops within society, and stands in the same relation to it as the distinguished element (distinto) always stands to the whole (indistinto). And society develops --- like the solar system and like thought --- by a natural process (formazione naturale), which it is the task of sociology to describe, and which ethics presupposes. It is the task of ethics (as nomologia, in contrast to sociology as nomografia and nomogonia) to discriminate between those elements in social life that have lost their significance and those that can be made fruitful in the same or a new way. Thereby it paves the way for the practical morality of the future.

Through the original and constant connection with social life there will develop in the individual an antiegoistic feeling. (Ardigò prefers this word to the word altruism). Human ideas receive their first content from the social environment. Its judgments and interests become provisionally valid for the individual. And all ideas have from the first an impulsive character, a tendency to pass immediately over into action; it is due to a later differentiation (distinzione) when thought and action are separated. Both, the involuntarily received social content of ideas and the original impulsive character of ideas, Ardigò comprehends in the concept of social ideality (idealità sociale). By ideality he understands the capacity to be determined by thoughts that point beyond the moment. Social ideality is developed already in the family as an enduring society that tends and protects the germs of the coming generation. Family feeling becomes a hearth fire from which warmth can radiate into larger circles. Social ideality has many degrees and forms. It is practiced by imitation and repetition. --- In this whole doctrine Ardigò develops a
% --- p. 41 ---
\opage{41}line of thought that is found in antiquity especially among the Stoics, in modern times in Adam Smith, and at the present day in Tarde, Leslie Stephen, and Baldwin. The latter's concept of social heredity would fit well into Ardigò's philosophy.

Social ideality is not developed to an equal degree in all, although the capacity for unegoistic action must be said to lie in human nature. Family love bears witness to this, as does involuntary sympathy, --- the craving for honor, which is properly meaningless in an egoist, --- and indignation at outrages and violations of right. At its high point the ethical disposition becomes a kind of holy rage (furore santo), which leads to self-sacrifice without thought of one's own advantage, in the faith that out of the tragic downfall of the human an eternal and divine will rise (eterno divino che sorge delle ruine tragiche dell' umano, an expression Ardigò borrows from a kindred Italian author). The truth in the theological idea of grace is that there is an involuntary impulse, a passionate yearning, which leads man to disinterested actions.

Such heroic action is possible without religion, just as ethics in general is independent of religion, when by religion one understands not a relation to the Infinite (after Max Müller's definition), but a relation to the supernatural. The scientific concept of the Infinite stands in sharp opposition to the supernatural: it denotes the continuous function whose limited and particular form the finite is. The Infinite is the law by which the individual's own essence is conditioned; thereby the deep satisfaction the individual feels when his yearning toward the Infinite is satisfied is understood. The positivist thus finds the Infinite in himself when he becomes conscious of the law of his essence. Here again the relation between element (distinto) and whole (indistinto) repeats itself. The idea of the supernatural, by contrast, arises when the law of things is made into a thing different from the activity of things themselves, which then can be set in relation only in an external manner to what goes on in nature.

Yet the concept of the supernatural denotes only the theoretical side of religion, and does not make up its whole essence. The
% --- p. 42 ---
\opage{42}essential thing is the relation of dependence in which man, with regard to his whole weal and woe, feels himself to stand to the supernatural. The Holy was from the first the Fearful; it is fear that creates gods. In any case the feeling toward the Holy is not, as has been asserted, an entirely simple and unanalyzable feeling\textsuperscript{18}). Gradually the element of fear is separated out in the transition from religion to scientific consciousness, which everywhere puts law in place of the idea of a mysterious being as that which determines man's fate. Religion always finds a certain social ideality and takes it up into itself; therein lies the fact that morality is independent of religion. Many transitional forms are possible here; but in the end it is a matter of an either-or. Artificial maintenance of religion leads only to the people's standing unprepared when its time is truly past. ---

It is a rather elementary psychology of religion on which Ardigò builds, and he has not undertaken an extensive discussion of the religious problem. His standpoint is essentially determined by his intellectual interest, and presumably also by an effect of contrast, which the sharp opposition between his earlier and his later view of life would naturally entail. It must nevertheless arouse wonder here how religion, after his reversal, could stand for him as a ``poetic memory,'' when he regards fear as the core of all religion --- from the lowest to the highest forms. The energetic thinker has surely not here made use of his own real experiences.

% --- p. 43 ---
\chapterheadb{III.}{FRANCIS BRADLEY.}

\opage{43}1. If one attends to the spirit and tendency of English philosophizing in the most recent period, one cannot avoid coming to the result that the classical English school --- that which begins with Locke, and whose latest representative is Spencer, --- has ceased to exist. It can be said to have fulfilled its mission, which consisted in asserting the right of experience, in demanding that problems be posed on the ground of reality, and in preparing practical, reformatory endeavor. It began with Locke by laying the individual's experience as foundation, and it ends with Spencer by laying the race's experience as foundation. What has here been inculcated has been handed on to other philosophical tendencies and has worked to correct and remold these. It is no longer the distinguishing mark of the single tendency\textsuperscript{19}). But the English school had on its side certain deficiencies, which stood out glaringly the more it came into interaction with other tendencies. To this belongs the mechanical atomistic conception, which it carried over from natural science to the science of mind, and which led it to regard mental life as a product of independent psychical elements and society as an external union between independent individuals. In contrast to this, in the most recent period, partly as a reaction, partly through the influence of German thought, the problem of the whole has come forward more strongly. The question is asked about the ultimate connection between the elements, which must be presupposed,
% --- p. 44 ---
\opage{44}in order that they may come to form a whole. There was, to be sure, also earlier in the nineteenth century a movement that went against the ``insular'' philosophy; names like Coleridge, Carlyle, and Hamilton bear witness to this. But a methodical study of German thought is first undertaken in earnest in the last quarter of the nineteenth century. It is Oxford University that leads the way here. There had always been earlier resistance here to the national English philosophy (to Hobbes, to Locke, to Hume); yet it had been more on theological than on philosophical grounds; it is only now that a truly philosophical opposition asserts itself.

As the one who first aroused a larger spiritual movement of this kind must be named \textsc{Thomas Hill Green} (b. 1836, fellow of Balliol College 1860, Professor of Moral Philosophy 1878, d. 1882). Strongly influenced by Wordsworth and Carlyle, he developed a peculiar religious idealism, whose philosophical foundation was owed to the study of Kant and Hegel. He had an extraordinarily great influence on the student youth at Oxford, especially through the ideality and enthusiasm of his personality. He worked not only in a scientific but also in a socio-reformatory and a free religious direction. Almost the whole younger generation of social thinkers at Oxford owes its awakening to him.
His chief works are \textit{Introduction to Hume} (1874), a penetrating critique of the principles of the older English school, and \textit{Prolegomena to Ethics} (1883), a peculiar combination of epistemology and ethics, through which a non-empirical grounding of the ethical is sought.

2. England's most considerable thinker in the most recent period is undoubtedly \textsc{Francis Herbert Bradley}, who was born in 1846 and lives as a fellow of Merton College, Oxford. During his course of study he was influenced by Green and by the writings of Hegel and Lotze. His inclination to solitude and his reserve were strengthened by ill health. He is a Spinozistic nature; only he lacks Spinoza's realistic eye for psychological and social phenomena. With great energy of thought he immerses himself in a single idea, which again and again sets his reflection in motion and brings him now closer to skepticism, now to mysticism. Even sharp
% --- p. 45 ---
\opage{45}critics acknowledge the stimulation and the discipline of thought that the study of his chief work especially produces. Yet his acuteness in detail often becomes hair-splitting, and in his special analyses the lack of empirical material is often felt. In energetic and indefatigable reflection he stands perhaps highest among the thinkers of our time, in any case as regards the outermost problems of the boundary. He forms a counterpart to Wundt. Whereas the latter approaches the boundaries of thought cautiously, but concludes somewhat hastily and dogmatically at the boundary, Bradley hurries too hastily to the outermost boundary, but then moves there with watchful critique and discusses the problems from all sides.

Bradley's first considerable work is his \textit{Ethical Studies} (1876). He here takes the field against the atomism of English psychology. Consciousness, he maintains, cannot be described as a mere collection of elements; for it would be impossible to understand how such a collection could become aware of itself (aware of itself). This problem had (what Bradley has not noticed) been posed by Stuart Mill himself in one of the later editions of the ``Examination of Sir William Hamiltons philosophy.'' It led Bradley to the study of German philosophy, which has particularly emphasized the unity and inner coherence of consciousness. He judges English thought severely, blaming it for one-sidedness and dogmatism and for lack of boldness, especially as regards the religious problem. ``We live on an island, and our national way of thinking will, if we do not widen it, acquire an island-like (insular) character''. --- This judgment is itself one-sided, however; Bradley forgets that critical philosophy was properly already founded by Locke, and that English philosophy of experience, with its demand for a showing of the origin of ideas, has been a mighty weapon against dogmatism.

Man's ethical striving, according to Bradley, aims at realizing his self, at developing into a whole, a harmonious and completed whole. For it is the essence of the self to be a totality, --- at once completed and rich. In a similar way theoretical striving aims at thinking existence as a coherent and completed whole (a consistent whole). When
% --- p. 46 ---
\opage{46}we cannot become a whole ourselves, we must make ourselves part of a more comprehensive whole, just as, where a smaller whole presents a contradiction, we must form the thought of a larger whole. There is thus agreement between our practical and our theoretical nature. And we have in our own nature a standard for what is to be called higher and lower: it depends on the degree of self-realization, thus on the harmony and the independence that are in our thought and in our life. Pain at contradiction man could not feel if he were not himself a whole and had an inkling that he was. Contradiction is called forth partly by lack of harmony within, partly by lack of agreement with external conditions.

The standard for theoretical and practical perfection that Bradley here set up gives the fundamental thought in his whole philosophy. It contains the germ of an entire view of life and world-view, as he later (1893) developed it in his chief work \textit{Appearance and Reality,} after he had in his \textit{Principles of Logic} (1883) given an exposition of the principles of knowledge\textsuperscript{20}).

The work on ``Phenomenon and Reality'' Bradley characterizes as a metaphysical treatise, and he explains this expression by saying that an investigation is intended of the principles that underlie the recognition of reality as opposed to phenomenon or appearance (appearance), or the standard by whose help we can distinguish between higher and lower degrees of reality. A system he will not give.

Bradley imagines various objections against the attempt he here wishes to make. --- One considers the attempt perhaps hopeless. But beforehand this cannot be known --- except precisely by virtue of a metaphysical insight, i.e. an insight into what standard must be applied in order to distinguish between reality and appearance. --- Or one holds that the result will be worthless. But even if the result is imperfect, it will have its value if it serves to illuminate what reality is. Indeed, even if we ended in complete skepticism, there would thereby be provided a salutary counterweight to the dogmatic tendencies --- theological orthodoxy on the one side, the
% --- p. 47 ---
\opage{47}vulgar materialism on the other side, which otherwise divide minds between them. --- After advancing these remarks Bradley adds: ``There is yet one reason, which as far as I am concerned perhaps has the greatest weight. I believe that we all more or less feel ourselves drawn beyond the domain of the usual facts. Each in his own way we think we stand in contact and communion with something that lies beyond the visible world. In different ways we find something higher that both bears us and humbles us, both chastises and inspires us. And for some natures the intellectual striving for an understanding of existence is the most important way along which they experience the Divinity. No one who has not felt this, however differently he may express it, has troubled himself much about metaphysics. But where it is strongly felt, it is its own justification.'' To a closer understanding of this utterance, which opens a view into his inner life, and to avoid misunderstanding, Bradley adds:
``I have been compelled to speak of philosophy as a satisfaction of what one might call the mystical side of our nature, a satisfaction that for some individualities cannot be attained in any other way. It might perhaps look as if I thought that metaphysicians are initiated into something that is far higher than what the great multitude can possess. Such a doctrine, however, would rest on a lamentable error, namely on the superstition that mere understanding is the highest side of our nature, and on the false idea that the intellectual work which is applied to higher objects is therefore higher work ... No calling whatever, no endeavor whatever is a private way to the Divinity, and the way that leads through speculation on the ultimate truths does not stand above other ways. There is no sin which philosophy can so little justify as spiritual pride, however inclined the philosopher may be to it.''

In these utterances Bradley has contributed not only to his own characterization, but also to the understanding of the innermost core of all philosophical striving, when this is on the grand scale. The inward craving for understanding and the constant experience of the limit of thought, --- the passionate exertion
% --- p. 48 ---
\opage{48}of the energy of thought and the constant feeling that it is not we who think, but that ``it thinks in us,'' --- the intensive devotion to the life of thought, as if it were the only valuable life, and the constant experience that there is also another craving for life in us, which demands satisfaction, --- all this can be read out of the words of the English thinker. These words are worth noting in face of the handicraft character that scientific work in many fields seems to be in the process of assuming.

Having considered Bradley's general relation to philosophy, I now go on to bring out the fundamental thoughts in his chief work.

3. It is a matter of a critical investigation of the ideas by whose help one has sought to understand existence.

Such concepts as matter, space, time, energy, with which natural science works, are excellently suited to determine the mutual relations of limited phenomena, but lead to contradictions when they are to express the true essence of existence. They are relational concepts, which characterize things in comparison and connection with one another, but say nothing about the individual things themselves that stand in these relations. They lead into infinite series, in that one can always ask about the relation between the terms themselves and the relations in which they stand, and in that the investigation of the terms themselves leads to establishing relations within each of them. Such concepts are intellectual constructions, which can be expedient and necessary within the individual science, but do not illuminate the innermost essence of existence. They are working thoughts (working ideas), whose significance is technical, not purely theoretical. Therefore natural science cannot be metaphysics; but therefore there also cannot be any conflict between natural science and metaphysics. If one founds a metaphysics exclusively on the fundamental concepts of natural science, it receives the stamp of materialism.

If one now takes refuge in the fundamental concept of the science of mind, the concept of soul, one has the advantage that in this field we find a more intimate connection of unity and manifoldness than in the field where the natural-scientific
% --- p. 49 ---
\opage{49}fundamental concepts find application. Unity and manifoldness do not stand in mental life in the external relation to each other that they do in physical nature. Therefore psychological experience is the highest experience we have. But neither is it suited to express absolute reality. Within the self experience and analysis show us oppositions, differences, and relations, and the self as a whole undergoes changes. It is therefore impossible to comprehend its essence in a single intuition, to form a completed concept of it. Therefore we can as little build a metaphysics on psychology alone as on natural science alone. Idealism can as little as materialism express the full truth. Psychology is a special science, and every special science finds only half-truths and operates with expedient fictions. The concept of soul is an abstraction, just as much as the concept of body. Reality can be neither ``soul'' nor ``body.'' What we have given is always only events, which present two sides to consideration. ---

But the presupposition of this whole line of thought, which leads to rejecting both idealism and materialism, is nevertheless that we have a standard for what reality and truth are. Such a standard lies in the \textit{concept of experience.} In this concept two things lie in close mutual connection: an extent and a harmony. Experience presupposes a given manifoldness and an inner agreement, a consistent and harmonious relation between the manifold elements. Perfect experience would consist in an all-embracing content combined with full consistency and harmony into a whole. Every experience we actually have contains an approximation in the direction of this ideal.

The standard for reality and for value is the same. In our practical ideal every side of our nature must be satisfied in harmony with the other sides. Every unsatisfied impulse is an unfulfilled thought; in every pain a disharmony expresses itself and there lies a spur to work one's way out of this disharmony.

This standard shows us how far we stand in connection
% --- p. 50 ---
\opage{50}with reality. We cannot form any idea of that which completely satisfies the standard. ``We cannot construe the one absorbing experience to ourselves''. The Highest must be all-embracing and absolutely harmonious. But for us there is a constant opposition between extent (extension) and concord (harmony) or self-agreement (self-consistency). And the reason is that our extent is too small and our harmony too incomplete. Inner disharmonies are due to limitation and can be removed only through a more comprehensive content, which removes the dependence on external conditions and thereby makes it possible to attain inner coherence. It is the disturbances from outside that cause the inner disorder. Therefore only an infinite being can be completely harmonious. The two components in our standard for reality (and for value) thus hang closely together. In the theoretical field this shows itself in that we can remove the contradictions that experience presents only by widening our experience.
The relation of time stands especially as an obstacle to attaining harmonious totality. It thus stands in conflict with the criterion for reality, and therefore it can have only phenomenal significance. Of progress and regress there can be talk only in the case of limited and imperfect realities; the Absolute can have no history, although it contains historical processes without number; it has no seasons. Nothing perfect, nothing truly real changes.

Our thought therefore constantly strives after something that is More than thought, our personality after something that is More than personality, our morality after something that is higher than all morality. There is no contradiction in a being's longing for a fulfillment in which it itself is lost. The river flows to the sea, and the self loses itself in love. The Higher must always be more comprehensive than the Lower, must include or be this and More besides. ---

Philosophy leads, according to Bradley, to a sound skepticism, for which science is a slight thing in comparison with the richness of existence. Even the highest truth we can find is conditioned by something we do not know. We do not know what other kinds of experience may exist besides our own. What we can
% --- p. 51 ---
\opage{51}attain is only a justification of our right to call something higher or lower than something else.

Religion fares in the face of the ultimate riddles as philosophy fares. It too is compelled to express the Highest by ideas that are taken from experience. And it does not even undertake, as philosophy does, an exact investigation of the nature and validity of such ideas. Seen from this side, as knowledge, philosophy stands above religion. But in another respect --- as an endeavor to let the recognition of the reality of the Good permeate every side of our being (the attempt to express the complete reality of goodness through every aspect of our being) --- religion stands highest. ---

Is this result now not unsatisfactory? To this question Bradley replies with the counter-question: Who says that we are to have full satisfaction of all our craving? Incompleteness, unrest, and unsatisfied ideality are, after all, the lot of the finite! And he adds that every attempt to imagine a full satisfaction always makes a selection among our needs, a selection whose justification cannot be demonstrated.


4. Bradley's philosophy is a thinking-through of the problem that Wundt and Ardigò approached, but did not take up for express treatment, --- that which Kant had posed by his doctrine of the ``Ideas'' as concepts of totality, and which has its origin in the very nature of thought as a comprehending activity (synthesis). Bradley himself thinks he stands in great debt to Hegel, and he is, like several English philosophers in the most recent period, impressed by Hegel's bold dialectic. But in reality, if he were finally to be classified, he would far rather be called a Kantian than a Hegelian. On individual points he recalls William Hamilton. With a firm hand he draws, at the very boundary of thought, the consequence of the laws of thought itself. It is especially two laws on which he relies: the law of relation (the relational way of thought) and the criterion of reality (complete experience as standard of reality). They lead him both to the same result, --- to the impossibility of stopping at any point, although we
% --- p. 52 ---
\opage{52}can step by step arrive at a closer determination of reality.

A certain skepticism characterizes Bradley, in that he lays more weight on the impossibility of a conclusion than on the closer determinations of reality that are possible. Thereby he comes to ascribe too little positive significance for the world-view to specialized and empirical knowledge. The standpoints with which work is done in the special fields of experience he dismisses as ``useless fictions,'' ``mere practical compromises,'' without taking into account that they would be of no ``use'' if they did not in some sense lead us nearer to reality. He forgets Goethe's words, that nature has neither kernel nor shell.

He is also a skeptic in the sense that there are several individual problems that he declares insoluble. We cannot derive the manifoldness of existence, the many finite centers from which experiences are had, all the fragmentary, at which we must stop, from one principle, even though the fact of manifoldness does not conflict with there being a comprehensive unity. Why there should be ``appearances'' (appearances) and not merely ``reality,'' we cannot say, but this does not prevent us from holding fast to the concept of reality. In particular the inorganic is a boundary for our knowledge. The organic satisfies with great approximation our concept of true reality as harmony and closedness, although in the organic field too there are found partiality and disharmony; but the inorganic stands far further from this concept of reality. Bradley is inclined to think that the assumption of an absolute inorganic in nature is due only to our ignorance. He concedes that the special science for practical reasons must distinguish between organic and inorganic, but does not dare to ascribe absolute significance to this distinction. Here his distance from Hegel appears clearly; for romantic philosophy of nature worked boldly with the thought of the whole of nature as one great organism. --- The problem of the relation between soul and body he also regards as insoluble. It is impossible to form a concept of how these two forms of existence are related to each other. But in this he finds a
% --- p. 53 ---
\opage{53}confirmation of his general theory of reality; for the difficulty arises from the fact that by setting soul and body in opposition to each other we make two abstractions into real beings and thus treat ``appearances'' as if they were realities.


Yet skepticism would not be the right word for Bradley's standpoint. He does not stop at a chasm between the work and the goal, the phenomenon and reality. At every stage the Highest is present, and every stage has its truth; there are many degrees and stages, but they are all indispensable. We cannot find any province in the world so slight that the Absolute should not dwell in it. --- He is most nearly to be called a mystic, and he is decidedly so when his thought comes to rest, and when he polemicizes against the concept of time and against the significance of activity. He passes here over into an immovable contemplation, a resting gaze, a beholding sub specie æterni. It has gone with him here as it went with Spinoza. For Spinoza's ``substance'' is properly the standard for reality beheld as a perfected being, the criterion of reality as an existing ideal.

I cannot see that we can here get beyond the constant possibility of new processes, new activity. And why should it not be the Highest to develop through time? Especially since it will always prove impossible to eliminate the concept of time from a knowledge of reality. From the constant incongruence between thought and reality and from the constant necessity of new work of thought --- for the individual and for the race --- one might perhaps be justified in drawing the ``metaphysical'' consequence that existence itself is not completed, is not finished. And it \textit{cannot} be finished when thought, which after all is a part of existence, is not finished. This conception leads to a more intimate relation between thought and reality than that in which Bradley ends. For when time and activity are not forms that belong merely to the phenomena, ``the appearance,'' thought will precisely as striving and working be able to feel itself one with the innermost of existence\textsuperscript{21}). --- But I will not here go further into what I have developed in the
% --- p. 54 ---
\opage{54}epistemological part of my Religionsfilosofi and in my ``Filosofiske Problemer.''

In England itself Bradley has been criticized partly from an empirical and critical standpoint (thus by \textsc{James Ward} in ``Mind'' 1894), partly from a standpoint that emphasizes the economic and technical character of our fundamental concepts and grounds an idealistic world-view on personal need (thus in the work published by eight Oxford academics, \textit{Personal Idealism.} London 1902). This latter tendency considers itself a continuation of the line of thought that has prevailed in Oxford in the last thirty years. It is most in the right in its critique where it disputes Bradley's negative attitude toward truths that hold within determinate conditions, and maintains that we can have valid knowledge even if we cannot attain an exhaustive understanding of existence. On the other hand it does not seem to me that the English critique has sufficiently recognized the real profundity and energy of thought in Bradley's work.

% --- p. 55 ---
\chapterheadb{IV.}{ALFRED FOUILLÉE AND FRENCH\\ PHILOSOPHY IN OUR DAY.}

\opage{55}1. In the history of French philosophy Auguste Comte
(† 1857) stands as the greatest figure of the nineteenth century. The
influence he has had far and wide found expression in the unveiling
of his statue in May 1902. Positivism is the most distinctive
and significant tendency of thought that France has produced in
the last hundred years. For the greater part of the century it had as its opponent a popular spiritualism,
represented by Victor Cousin and his disciples
and favored by the governments. Yet from about the middle
of the century Taine and Renan exerted a great influence on
the younger generation, not only in philosophical but also, and quite as
much, in literary respects. Beside them, but as a
more solitary figure, stands Charles Renouvier, representative
of a strict critical philosophy. --- There is on this point a
gap in my account of the history of modern philosophy, in that I
stopped at Comte as far as France is concerned. This gap
I do not now see myself in a position to fill either. I can only
give a brief characterization of Taine and Renan as an introduction
to the French form of the philosophy of evolution, represented
by Fouillée, and later a similar characterization of Renouvier in
connection with the role that the principle of discontinuity comes
% --- p. 56 ---
\opage{56}to play in the most recent French thought. A monograph
(or several) on these three men as philosophers would be of
no little interest in throwing light on the course of
intellectual development from the middle of the nineteenth century on. The Royal Danish Academy of Sciences and Letters
has attempted to call forth a monograph
on Renouvier, whose extensive authorship presents
peculiar difficulties for survey and characterization;
but no answers to the task set were received.

\textsc{Hippolyte Taine} (1828---1893) is best known as
a historian of literature and art and as a critic in the ``realistic'' spirit.
Our literature possesses a valuable characterization of him from
this point of view in \textsc{Georg Brandes's} essay \textit{Den franske
Æstetik i vore Dage} (1870). As a critic Taine strove first of all
to understand works of art and their producers
by investigating the external conditions (le milieu) under which they
developed, the situation (le moment) in which they worked,
the race from which they had sprung, the dominant faculty
that expressed itself in their work. This last element was essentially
conditioned by the first three.

From his earliest youth Taine worked under pressure and
resistance, but with unabated energy and enthusiasm. Letters
from his young days, recently published, give a beautiful picture of
him from this time. Later, as professor of the history of art
at the École des beaux-arts, he had occasion to reach
larger circles and came to work under more favorable conditions,
an occasion he used with the iron diligence that was
his hallmark throughout his life. Characteristic of Taine as
speaker, writer and inquirer was above all his ability to produce
a picture of the whole through the accumulation of single features.
At his lectures one could --- by my personal
experience --- at the outset feel repelled by the dry,
objective description that went from part to part, from property to
property. But there was such a wealth in the material that stood
at his disposal, and such an energy in the way in which
it was brought forward, that the total picture before the end of the hour stood
clear and living before the imagination. I remember, for instance, a lecture
on Greek sculpture. It was not accompanied by
% --- p. 57 ---
\opage{57}pictures, but Taine sat during the long description he gave
and drew in the air, and at last it was as if a
statue stood before him on the table. His art of presentation went
in the opposite direction from that of Julius Lange, whose strength consisted in
analyzing a picture of the whole that was called forth by an illustration
or by immediately awakening the imagination of the hearers. Taine went from part
to whole, Lange from whole to part.

A purely philosophical work Taine produced in \textit{De l'intelligence}
(1870), a psychology in the spirit of the English school. It betrays
the strong influence of J. S. Mill, Bain and Spencer. Of special interest
in this brilliantly written book is, in the first place, the explanation
of the development of knowledge through a struggle for existence
among the individual psychical elements. There arise
continually elements with the stamp of reality, a kind of normal
hallucinations; they clash with one another, and the victor is
a sensation, a sense perception, which can therefore be defined as
une hallucination vraie. In this clashing the original tendencies to
movement of the elements of knowledge then take part.
Every sensation or idea is connected from the first
with an impulse to movement. Gradually
this motor tendency is worn away, and only then do purely theoretical sense perception
and idea come into being. Finally it is also characteristic
of Taine's psychology that he makes such extensive use
of morbid states to illuminate the nature of the life of consciousness.
--- This last trait has since remained characteristic of
French psychology, perhaps because French nervous pathology,
especially in \textsc{Charcot} and his school, supplied such great material. \textsc{Ribot},
whose first writing, \textit{La psychologie anglaise contemporaine,} appeared
in 1870, has in his later writings gone especially this way; likewise
\textsc{Alfred Binet} and \textsc{Pierre Janet}. Modern French
psychology has thereby received its distinctive stamp, just as English
psychology (in Ward, James and Stout) has its in analysis and
German psychology (in the disciples of Fechner and Wundt) its in the experiment.


Earlier (especially in the writing \textit{Les philosophes français du 19e
siècle)} Taine had sharply criticized the dominant spiritualism
(Cousin's school) and pointed to Comte. A positivist he was
% --- p. 58 ---
\opage{58}yet only in the wider sense of that word. Positivism was for
him, as for Ardigò, a starting point, but not a conclusion.
In part he demanded that the analyses of concepts and of states be carried
further than the positivists of the strict school found necessary,
in part he held fast to the problem of a conclusive world-view
and to that extent believed in a metaphysics. With
an eye to this problem he ends his work De l'intelligence thus:
``Je vois les bornes de mon esprit; je ne vois pas les
bornes de l'esprit humain''. In Taine's famous works on the history of literature and
art psychology plays a great role, often
too great a role, in that he is inclined to proceed deductively
and to derive an artistic phenomenon exclusively from a
single element (a dominant faculty or a milieu). That hinders
him at times, e.g. in his treatment of Shakespeare,
from catching the finer nuances and the more inward expressions of the soul
that reveal themselves only after prolonged
observation and more patient analysis. These works all date from before 1870.
It had now been his intention to continue his purely psychological
work and to supplement the book on knowledge with
a book on the will. But the terrible misfortunes of his native country
led him into historical studies. He wanted to understand the
fate of his country from its past and through understanding to win
new hope. From this arose the great work \textit{Les origines de la
France contemporaine} (1876---1894). Of philosophical interest are
especially the first volume (L'ancien régime), in which, however, on essential
points he follows Tocqueville, and the last two volumes (Le régime
nouveau), which no one who wishes to understand the nineteenth century may
leave unread. They are important contributions to sociology. Of his
treatment of the French Revolution (in the volumes
between) it has rightly been said that it laid the foundation of the scientific
view of this event, as opposed to
the declamatory and agitational one. Against it, on the other hand, it has been objected
that he adheres too much to the general causes
and underestimates the influence of the special and momentary
political situation. His hatred of the phrase has moreover
carried him so far that he misjudges the enthusiasm that the great
% --- p. 59 ---
\opage{59}hopes awakened. One does not understand, from his account,
that this was the time when the Marseillaise came into being.

Taine's view of life was determined by his temperament.
He was by his own statement un homme naturellement triste.
The coup d'état of 1851 and the defeat of 1870 had to increase his despondency
in the face of external conditions. His best life he lived in
work and in the world of thought. He felt a need to immerse
himself in the great, necessary coherence of existence. Marcus
Aurelius and Spinoza were his favorite authors. In religious
respects he sympathized with Protestantism, and after
his visit to England he hoped for a free religious movement.
But he remained a Stoic to the last and still in his
last days read Marcus Aurelius ``as a kind of liturgy''\textsuperscript{22}).

\textsc{Ernest Renan} (1823---1892) presents in his course of development
an unusually long series of stages, and the biographical
art that is to be able to find the inner coherence
here has not yet been applied. His childhood and his youth
he has himself described in that one of his books which stands
highest in the inwardness and depth of its mood \textit{(Souvenirs d'enfance
et de jeunesse).} It was, by his own assertion, historical
criticism that led him to leave the seminary and
the Catholic faith. Meanwhile there soon developed in him
an enthusiastic faith in science and in the significance its
results would have for the great multitude. This philosophical
and democratic optimism he expressed in a writing on ``The Future
of Science'' \textit{(L'avenir de la science),} which he composed
in the years 1848 and 1849, but which he published only much later,
after he had lost the hope expressed in it. It was
the coup d'état and the reaction that instilled in him a pessimism that
he never overcame. Later came the defeat and the rising of the Commune:
as a scholar he became especially famous through
his works on the history of languages and of religions, which it is not here
the place to discuss more closely. Under the Third Republic he obtained
a high scientific position, as Director of the Collège de France,
and all the while he was putting the last touch to his works,
his reflection occupied itself in a light and jesting manner with the
great problems that had engaged him so strongly in his young
% --- p. 60 ---
\opage{60}days. Large circles listened to the skeptical causeur,
of whom it was known --- at least at third and fourth hand ---
that he was a remarkable man. In his utterances
and writings of this period he was now full of feeling, now
saucy, --- now sublime, now blasé; and for many this
constant oscillation, --- an oscillation whose center was not easy to
discover, if there was one, --- was an expression of true freedom of spirit
and wealth of spirit. Work of thought was now for him himself
rather a game, by which he distracted himself from his more rigorous
studies. And he had become apprehensive of every definite
position on the questions at the boundaries, because he himself had taken
so many different positions. Thereby he cheated
himself and others of a valuable result of his long striving
and inquiry. It has rightly been said of him: Il fut
dupe de la peur d'étre dupe.

Renan's significance as a thinker in the later years of his
life lies more in his standing as a symptom of a certain tendency
of the time than in his having done real work
in the service of the problems. The secret was perhaps
this, that his heart clung to the religion of his childhood, that
into which he had lived himself while he was ``brought up by women
and priests'' (an upbringing in which he himself thought
the explanation of his merits and his faults had to be sought). What
animated him later on was an echo of it --- comparable
to the sound of the bell sunk in the sea, which
even in still weather is said to be heard on the coast of Brittany. A
general idealistic faith is indeed possible; but it suffices only
so long as we still live under the dominion of the old habits.
That will not last. We live --- says Renan in several places
--- on the shadow of a shadow; what will men live on after us?
To a firm and manly position in face of this question
Renan did not come: it was only in certain moods that he felt its
sting.

There are two works by Renan that are of special interest for
his philosophical characterization in his later years.

\textit{Dialogues et fragments philosophiques} was composed in the spring of
1871, when Renan during the rising of the Commune stayed in
% --- p. 61 ---
\opage{61}Versailles, but was first published in 1876. They are conversations that
are conducted between ``three philosophers of the school whose basic principles
are the cult of the ideal, the denial of the supernatural, and the investigation
of reality''. They discuss the question of what the
ideal end of the development of the world is. For an end there must
be; the universe is not an empty surge of waves that ends in nothing.
The end of development cannot be sought in the great mass of
beings, who have a minimum of enjoyment and of cultivation, but
a maximum of labor. The end must be the dominion of reason,
and only a small circle can come to share in it. The
significance of the great mass can only be that it is the soil
from which the geniuses grow. The geniuses --- and the intelligent,
who understand how to value them --- are the final end of history.
It is regrettable that the ignorance of the great mass is
the condition of this end being reached; but nature does not trouble herself
about such things. It would lead to degeneration, to a lowering of
the level, if one were to go the ways of democracy; along these ways
no gods are produced. Democracy and science are antipodes.
Development is marked by instinct being replaced by reflection,
religion and art by science.

Renan sees in this something that the Middle Ages already had an eye for in their
own way. He found a truth in the medieval
order, that there were men who prayed for those who
had no time to pray. The men of science are to be
the priests of the new age. And it is evidently an echo of
the priest in Renan that comes forward here and --- in any case
for a moment --- lets him land in this solution of the problem of culture,
which incidentally --- as we shall see later --- was expressed approximately
at the same time by a German author with more point and
punch, though with greater inconsistency. For Renan it seemed
as though the knowing, who are the end of development, might very
well love the people, so that their dominion might redound to the
benefit of the latter. Nietzsche tries to deny this.

Some years before his death Renan gave a kind of philosophical testament
in his \textit{Examen de conscience philosophique} (1888). (Printed
in ``Feuilles détachées''). A statement characteristic of Renan
occurs immediately at the beginning of this essay. ``The
% --- p. 62 ---
\opage{62}first duty of the honest man'', it says, ``is not to exert
influence on his own opinions, but to let reality
be reflected in him as in the photographer's darkroom and to stand
as a spectator at the inner struggles that thoughts carry on in the depth of
consciousness ... We ought to be passive in face of the inner changes
on our intellectual retina.'' --- It is of course
true that when we are to settle our intellectual account,
we must not intervene, but see what has really been thought in
us. But in the coming into being of the individual items of the account
a striving after definite ends takes part, an involuntary
yearning, just as our eye at every moment is directed so that that
which attracts attention can be taken in by the central
part of the retina. It will therefore be an illusion when
one believes oneself to be a purely passive spectator of what goes on
on the mental retina. In every scientific hypothesis
as well as in every view of life there is something of a venture.
The struggle of the fundamental thoughts and fundamental values we by no means follow
as disinterested spectators. And if we try it, it by no means
turns out to the advantage of thinking.

Two thoughts underlie Renan's philosophical testament.
The one is that both when we go into the great and
when we go into the small we come upon the infinite. However firmly
we must, within our world (or our part of the world),
hold to experience, we have yet so little right to suppose
that its results have absolute validity. From the standpoint of infinity
nothing is impossible, and the unsurveyable future makes
many difficulties vanish. Let us then deny nothing,
affirm nothing, but let us hope! --- The second thought supplements
the first. In the midst of all the enigmatic and uncertain that surrounds
us, there appear the four exalted authorities, love, religion,
poetry and virtue, which the egoist denies, and which yet
lead the world forward. Through them we hear the voice of the universe or,
if one will, of God. It is the harmony of the heavenly spheres,
the tones of the infinite. In them expresses itself the deep striving
(le nisus profond) that works its way forward through the development of the world
against the resistance of the recalcitrant matter (perhaps especially recalcitrant
in our part of the world). It will win its way in clarity
% --- p. 63 ---
\opage{63}and in power to a degree of which we cannot form any idea. ---

The philosophical testament is less playful and skeptical
than the earlier utterances. The theory of the great men
as the end of the world does not occur in it. But whether he has
retracted it, one cannot see. Nor can one see
whether he still holds fast to what a few years before he had uttered
at the end of his essay on Amiel (which is also
printed in Feuilles détachées). Amiel had reproached him for the
playful and ironic manner in which he spoke of the most serious
things. Renan answers that irony must be the last word of philosophy.
What the end of the development of the world is we do not know, not even
whether there is an end. Perhaps the whole is une
mauvaise farce. And let us then at any price not be
duped! By oscillating between the two possibilities --- that life
is earnest, or that it is a bad joke --- one at least avoids
being \textit{wholly} deceived, whereas by definitively choosing
the one or the other one exposes oneself to suffering a complete
disappointment. It is especially for the sake of others that one must be
cautious. For one's own sake one might embark on great
ventures; but one has no right to lead others astray into
exposing themselves to complete shipwreck. Therefore the point is to
be ad utrumqve paratus, and that one attains by oscillating
--- oscillating between faith and doubt, between optimism and irony.

\textsc{Charles Renouvier} has said to this, with justice, that by
such oscillations one easily loses one's title in both directions,
and he has pointed to the possibility that \textit{the choice itself}
might be the essential\textsuperscript{23}). It might be, after all, that the personal
energy that chooses and holds fast to its position and takes
the consequences of it stood nearer to the real reality in the world
than the two possibilities Renan depicts.

A psychological explanation of Renan's standpoint in his
later days one perhaps gets from his previous history and his
race. The contrast between the fullness of childhood faith and the apparent
emptiness of criticism had gripped deeply into his life, and
it seemed to him as though the only positive thing in him was really
only an echo, whose tones he could not prevent from
% --- p. 64 ---
\opage{64}dying away. For Ardigò too religion in his later period stood
as a poetic dream; but he built himself, with all the energy of thought
and will, a new, firm position. Renan's Catholicism has
no doubt also been of a different kind from Ardigò's. As a young theologian
he had become accustomed to receiving his experiences of life in a
passive manner, from without, not through energetic inner work, not
through active participation in life. Therefore it also came more easily to him
later to stand as a mere spectator at the struggle of thoughts,
as well as at the struggle of life. He let moods follow one another
in kaleidoscopic fashion; no inner process of crystallization took place
in him. And his Gallic temperament had
no capacity to let the opposed moods unite into a
whole with a definite timbre.

2. If one were to name an inquirer who in the philosophical
domain can be said to continue Taine's work, \textsc{Alfred Fouillée} (b. 1838)
would first of all have to be named. Taine's strong
emphasis on lawful connection is found again in
Fouillée, who at the same time stands in definite opposition to Renan's significant
philosophical dilettantism in his later years. But besides this
Fouillée's philosophical writings are testimony to the idealistic
tendency which philosophical thinking on the whole took
toward the end of the century.

Fouillée has served as professor of philosophy at Bordeaux and
in Paris, but later, for reasons of health, withdrew to
Mentone.

His first writings concern the history of Greek philosophy.
Through the study of Plato he worked his way forward to his
philosophical conception. \textit{(La philosophie de Platon.} 1869). For
Plato the world of Ideas and the world of experience stood in sharp opposition
to one another, so that only the former was the true reality.
Modern naturalism, on the other hand, asserts the world of experience,
that is: the natural connection given in experience,
as the only reality. When Fouillée from the study of Plato
returned to the problems of the present, it had to become
his task ``to lead Plato's Ideas back from heaven to earth
and thereby to reconcile idealism and naturalism''. Thus
he himself characterizes his endeavor. \textit{(Le mouvement idéaliste et}
% --- p. 65 ---
\opage{65}\textit{la reaction contre la science positive.} 1896. p. XXI). The possibility of
this reconciliation he builds on the fact that thought
(idée) can lead to action. He emphasizes in particular that the thought
of freedom can awaken striving and release force, whether
we are ``free'' in the indeterministic sense or not. \textit{(La liberté
et le determinisme.} 1872). Here we already stand at the fundamental thought
in Fouillée's psychology as well as in his whole world-view.
This fundamental thought he himself has expressed by the word
thought-force, \textit{idée-force.} This concept makes it, in his view,
possible to bring about an approximation between the different
philosophical tendencies. He first used the expression
in the Revue philosophique 1879. With the thoughts in consciousness
tendencies to movement in the brain are immediately united. It
is at once a physiological and a psychological fact that we here
have before us, so that there reveals itself here an ideal that expresses
itself in the very determinism of nature. Already prior to conscious
thought it works in the organic tendency of development and expansion
as well as in instinctive sympathy. There
stirs a will in the whole of nature, both in outer motion and
in inner sensation. The spiritual and the material, consciousness
and life, the individual and the social, freedom and
solidarity work here as one. \textit{(La science sociale contemporaine.}
1880).

In the most significant of his works, \textit{La psychologie des
idées-forces} (1893), Fouillée has developed his fundamental thought from
the purely psychological side. It is the best presentation of
the psychology of voluntarism. Fouillée directly defines psychology
as the study of the will. He sees clearly that our willing is not an object
of immediate observation; but instead of drawing from this
the conclusion that some draw, that then the will is nothing,
he draws --- in my view\textsuperscript{24}) with full right --- the
opposite conclusion, that our willing is one with ourselves. Psychology
has according to Fouillée hitherto suffered from intellectualism. It has
not been seen that the mental phenomena are always expressions of
need or impulse (appétition), accompanied by pleasure or pain,
according as they are favored or hindered. It is here a matter of indifference
whether we regard the phenomena from the physiological or
% --- p. 66 ---
\opage{66}from the psychological side. When we set these two sides in opposition
to one another, we must remember that in them we really only
have two abstractions (deux extraits d'une réalité unique et
totale). Every discrimination (discernement), even the most elementary,
presupposes already a choice, a preferring (préférence,
choix pratique rudimentaire). Discrimination and preferring in the
simplest case coincide, e.g. in the animal's discrimination between
the feeling of pleasure in eating and the feeling of displeasure in hunger. Only that
is sensed which is of significance in the struggle for existence, and it
is the will (this word taken in the widest sense) that drives
sense perception forward and conditions its more differentiated forms.
What holds for sense perception holds also for recognition
and memory. What is recognized and remembered is first
and foremost what has practical interest. Step by step
interest grows --- in interaction with experience --- out
beyond the momentary and the near at hand. Even the abstract logical
principles have formed themselves along this way. Every thought
or idea denotes a more or less conscious direction of our
striving and sensing life. There expresses itself here a great continuity
among all mental processes. Every state of consciousness
can be called thought (idée), insofar as it consists in a discrimination,
--- force (force), insofar as it consists in a preferring. ---

Fouillée's ethics is closely connected with his psychology.
He emphasizes in particular that I cannot become conscious of myself without
having consciousness of other beings: I apprehend them in
analogy with myself and apprehend myself in relation to them. With
this relativity in our apprehension a certain solidarity
and an altruistic tendency is already given. It is for me already purely intellectually
impossible to make myself the only being. The limitation
of intellectual arrogance entails a limitation
of practical egoism. Here we have again an idea that
is at the same time force, in that it is one with an involuntary striving,
which --- when it becomes conscious --- appears as an imperative.
The ideal lies in the prolongation of the direction in which our striving involuntarily
moves, and forms itself (by means of the connection
between relativity and solidarity) as the idea of a realm of freedom,
equality and justice. --- This conception
% --- p. 67 ---
\opage{67}of ethics Fouillée has not hitherto carried out in a connected
presentation, but has set forth in various places in his writings. It
underlies a brilliant critical work, of great interest
for the history of modern ethics. \textit{(Critique des systêmes
de morale contemporains.} 1883). ---

In his general world-view Fouillée lays particular
weight on there being no ground for interpreting existence merely
by the most elementary phenomena or the most abstract
points of view. Materialism is guilty of this when it makes
motion everything; motion is only the point of view of reality
that rests on the sensations of sight and touch.
Idealism is guilty of it when it makes thought everything;
thought is just as much an abstraction as motion is.
The individual sciences have in general a tendency each
for itself to make their points of view the sole expression of
reality. But every special point of view is one-sided. A
world-view must conceive reality as a totality
within which the special points of view each have
their justification and significance. And since psychical existence
is the only immediately given reality we know, it is
justified to interpret existence as a whole in analogy with it.
Since Kant's time thought has far too often been conceived as that
which separates us from reality, instead of being conceived as
that which unites us with reality and is itself a kind of
reality. And metaphysics is directed to build on the
only immediate experience we have. The experience that pain
is shunned is more immediate than that impact moves. --- His
metaphysical ideas Fouillée has developed especially in \textit{L'avenir de la
metaphysique} (1889).

All metaphysics is hypothetical and builds on analogy. It
is a supreme expansion, sprung from the need of personal life
for unity with the universal life. This universal
life is conceived as a great whole, a society of striving forces.
It is thus not merely psychological but also sociological
analogies that underlie the metaphysical ideas. The word
God --- which
always more or less expresses a representation drawn from human
relations --- denotes the innermost ground
% --- p. 68 ---
\opage{68}and striving in the universal society of existence. But beyond a
hypothetical scheme we cannot reach here. The final synthesis
that our thought attempts cannot be represented by us in the same positive
manner as a special connection between phenomena
of which we have a definite and concluded knowledge. We need
not with the Kantians introduce new dogmas in the name of
morality; just as little are the final ideas an object of irony
or mockery, as in Renan's dilettantism. In the very need to
act according to our ideals in spite of the uncertainty of their victory
there lies hidden an idealistic presupposition: c'est une spéculative
en pensée et en acte sur le sens du monde et de la vie!
The formulation of this presupposition in definite propositions is
always hypothetical. There is a constant struggle for existence
among the metaphysical systems. The question is which of them
thrives best in the prevailing scientific atmosphere. The possibilities
of a metaphysical conception will be more and more limited
as knowledge advances. The conception that will triumph
will be the one that lets both analysis and synthesis
come into their own in the most complete manner. But Isis's
veil will never be lifted, and there will always stand room for play
open to the imagination and to religious symbolism. --- This
last line of thought is directed not merely against Renan, but also
against Albert Lange and against a thinker who incidentally stands very close
to Fouillée, and with whom we shall come to occupy ourselves later,
namely Guyau.

3. Taine and Fouillée, as well as Wundt and Ardigò, are decided
philosophers of continuity. Their endeavor is to
find as much connection as possible within experience,
and to supplement experience in such a way that the principle of continuity is, as far as
possible, maintained. In contrast to this kind of attempt
there stands a series of endeavors that, partly on grounds of experience,
partly on ethical grounds, assert the significance of discontinuity. Such
endeavors are characteristic of the idealistic tendency
at the transition to the twentieth century, whereas the
idealism that prevailed at the beginning of the nineteenth century
was essentially a philosophy of continuity. In French
philosophical literature the philosophy of discontinuity appears in a
% --- p. 69 ---
\opage{69}particularly interesting and energetic manner. Three different motives
can be distinguished as guiding the philosophy of discontinuity.
Experience presents differences of quality that neither
speculation nor the doctrine of evolution has succeeded in reducing.
Comte's positivism had expressly recognized the gulf
that separated the different domains of nature from one another. Every
new science meant for Comte a peculiar, irreducible group
of phenomena. --- To this is added that even within each single
group of phenomena the principle of causation can find its confirmation only imperfectly.
One therefore draws Hume forth
again and sets his empiricism against the attempts of Kant and of evolutionism
to overcome it. --- Finally one points to the consciousness
of initiative, of the ability in one's thinking and acting
to put something new into the world, and one stresses strongly
the moral significance of this ability. --- In the two thinkers we
shall put forward as representatives of the philosophy of discontinuity
these three motives work together in somewhat different ways.

4. \textsc{Charles Renouvier} is the Nestor of philosophy in our time.
He was born in 1818, and even in 1903, thus 85 years old, he
has published a well thought-out work. He received his scientific
education as a pupil at the Polytechnic School and
was in his youth an eager Saint-Simonian and republican. Throughout
his life he has lived as a private man, but has exercised
no little influence through his books and through a
periodical that he published for a time. An energetic and knowledgeable
thinker and a serious character appears to us through
his writings; yet he is, like so many French philosophers,
strongest as a critic. He appeared both in the name of logic and of
morality against every kind of doctrine of continuity, whether it
rested on theological mysticism, on metaphysical speculation or
on natural-scientific formation of hypotheses. His activity as
critic and thinker begins long before the period that the present
work treats. There is here, as already remarked,
a gap in my account of the history of modern philosophy
as regards modern France, a gap I do not now either
see myself in a position to fill. I pass over here his earlier
works (of which the most important is \textit{Essais de Critique}
% --- p. 70 ---
\opage{70}\textit{générale.} 1854---1869) and keep to his writings of the
last quarter-century. For orientation regarding his
standpoint it will be important to study \textit{Classification
systématique des Doctrines philosophiques} (1885), \textit{Les dilemmes
de la métaphysique} (1901) and \textit{Le personalisme} (1903),
besides various essays in the yearbook published by Renouvier's friend and
collaborator Pillon, \textit{L'année philosophique.}

Renouvier has himself given an account of his course of development,
in that he lays weight on every system --- as the
work of a human being who lives and thinks under definite
inner and outer conditions --- being conditioned by personal factors
that have an influence on the point of view chosen. (See the essay
\textit{Comment je suis arrivé à cette conclusion} in the second part
of the great work on the classification of philosophical theories).

In his youth he had, as already mentioned, a period
in which he was an eager socialist of Saint-Simon's school, and at the same time
his mathematical studies led him to occupy himself with
the concept of infinity. By both ways he had come
upon a philosophy of continuity: as a socialist he emphasized solidarity
at the expense of individuality, and as a thinker he let
all finite differences fade away through infinite transitions
and continually continued processes. But in both respects
doubt was awakened, nourished by eager study of Descartes and
Kant. He now saw that every attempt to construct a system
that is to hold for existence as a whole ends in contradictions.
Kant's doctrine of the antinomies worked especially strongly
upon him, although he could not agree with Kant that neither
antitheses nor theses could be refuted; in his conviction
the theses, which assert the finitude of the world, were right as against
the antitheses, which assert its infinity. The contradictions
in which speculation finally entangles itself haunt not only the purely
philosophical attempts, but also every theological attempt to
unite God's immutability with his creative activity and
his infinity and foreknowledge with human freedom. Renouvier's
thinking is borne not merely by purely theoretical interest,
but also by the endeavor to pave the way for a rational
reform of the religious representations. --- Gradually
% --- p. 71 ---
\opage{71}Renouvier was led to a peculiar form of critical philosophy, which in
France is called le néocriticisme, and which I shall now characterize
in its main features.

The first principles of our world-view are determined,
according to Renouvier, by a choice between the oppositions that thinking
sets before us. --- Thus a choice must be made between
infinity and finitude. The concept of infinity contains
a contradiction if by it one means anything other and more than
a possibility for thought of always going further (e.g. of
adding ever new parts to my imaginative picture of space). As
actually given the world must be finite, although our experience
cannot embrace it. Kant's theses are right against his antitheses.
--- By this choice continuity is already broken; for
with the concept of infinity falls the possibility of infinitely many
transitions between the phenomena. From this it follows again that the principle of causation
can have only limited validity. Renouvier
here joins Hume and asserts with him the indemonstrability
of the principle of causation. New beginnings may occur, causes
that are not effects. This becomes of significance for
morality, in that freedom of the will thereby becomes possible. Philosophy
can therefore not become construction; it can only become a doctrine
of what happens before and after the pure beginnings, and
of how freedom is used in the right way. Freedom itself
is a first truth in the world of knowledge; it is a presupposition
that the determinist himself really makes use of
when he chooses to make necessity the first principle, ---
only that he thereby entangles himself in a contradiction that the philosopher of
freedom avoids\textsuperscript{25}). --- It is thus a pronouncedly voluntaristic
(or ``arbitrary'') epistemology that Renouvier embraces, ---
not only because the principles are freely chosen, but also because that
which is chosen is the principle of freedom, the principle of discontinuity. The
contradictions and difficulties that Renouvier encounters or finds
in the history of thought do not lead him to cease
thinking, nor yet to play with thoughts, but to hold fast to
the assumption that both logically and ethically stands before him as
the right one, with an act of will. He is Renan's decided antipode.
--- His fight against continuity Renouvier carries on not only
% --- p. 72 ---
\opage{72}as regards the individual, but also as regards the development of the race.
There is, he maintains, no all-embracing historical law,
but many laws, each valid for its phase in history, in that
each single phase is introduced by a new beginning, an act of freedom,
that could have turned out otherwise than it
actually did. In the appearance of great men he sees such new
beginnings, and he contests the view that regards
the arising of such personalities as the product of the slow
development of long periods. The creative in them is thereby
not understood. He sees a great danger in the historical way of regarding things,
which regards everything that actually happens as the
only possible; thereby one comes to justify everything.
He combats Romanticism, which has especially introduced this historical
conception, which rehabilitates all the past without regard to the
possibilities of a wholly different direction of development that are present at every
point. In his work \textit{Uchronie} (L'Utopie dans
l'histoire) (1901) he has tried to show that ancient
culture by no means needed to perish to be replaced by superstition
and barbarism, and he has depicted a reformation of the Roman Empire
and of the ancient religion that lay within the limits of the possibilities
and would have led us around the Middle Ages.
By this ``European history of civilization, as it
has not been, but as it could have been'', Renouvier wants to
strengthen the consciousness of freedom and of responsibility and to counteract the reactionary
tendencies that rest upon a cult of history.

Through his critique of the concept of infinity and of the principle of continuity
Renouvier came to see the significance of the law of relation or
the principle of relativity. The concept of relation became
for him the foundation of our knowledge, the fundamental concept that
manifests itself in all the categories with which knowledge
works. He joins here William Hamilton, who
had first pointed this out. All knowledge is reached through
finding out the relations in which things stand, and every object
we know therefore also stands before us as a system of relations.
The principle of relativity is a method (la méthode des relations),
which springs from the very essence of our consciousness. Consciousness
itself is a relation, but different from all other
% --- p. 73 ---
\opage{73}relations in that it is a relation between my self as subject
and my self as object (relation de soi à soi), thus a relation
between identical terms. Thereby it is the foundation of all other
relations, which become clear only by being referred to this basic relation.
--- From the law of relation it follows that all knowledge concerns only
phenomena, in that everything we are to know must appear
to consciousness and be determined by its laws. We know
things only as objects of representation. If one wants to go beyond
the phenomena, that can happen only by religious postulates. Yet
not even the postulates can be freed from the law of relation; God
can be as little infinite as the world can. But along
the way of the postulate we do get beyond the limits of experience.
--- In his earlier period Renouvier was inclined to transfer
his republicanism to theology and thought that pluralism (in
religious form polytheism) was a consequence of relativism.
Later, however, he found so much connection in the world that
monotheism stood before him as necessary, although he could not
take God to be an infinite being. It followed --- he thought
now --- precisely from the law of relation that one must end at a first
cause, thus at a creation, as a point where thought
stops (le point d'arrêt de la pensée). Otherwise we enter into
an infinite process, which gives us no solution of the riddles.
Renouvier holds that many thinkers confuse the constant displacement
of the problems (le reculement indéfini des problèmes)
with their real solution. But when one thus ends in
the idea of a world-cause that is itself finite, then
--- since the definitive end point for thought is now supposed to have been
reached --- every brooding over the origin of this world-cause is
fruitless. Here absolute ignorance begins.

The last choice, which determines the results of our knowledge,
stands between the concepts thing and person, between a world-view
that subordinates the personal and subjective
to the impersonal and objective, and a world-view
that conceives everything in analogy with the personal. It has cost
a long and great labor for human thought to reach
an objective way of regarding things from the chaotic personifications of the mythologies.
The point was to carry through the thought of the
% --- p. 74 ---
\opage{74}lawful, objectively given connection, which forms the presupposition
and object of the work of personal beings. But
the concept of the object (the thing, nature, substance, the ``Idea'') then easily gains
such an overwhelming power over thought that it forgets
that what is real is given only as object \textit{for it} as phenomenon,
and does not see that an intelligible conception of the world is obtained only
by laying our immediate experiences at the foundation of the interpretation
of the objective connection. By these two ways, the
epistemological demonstration of the logical priority of the subject
and the metaphysical interpretation of nature founded on analogy,
Renouvier arrives at a doctrine of monads, which he has presented in
his two most recent works, \textit{La nouvelle Monadologie} (1901) and
\textit{Le personalisme} (1903). It is his conviction that, after
the scientific understanding of the connection of nature has been secured
through strict and methodical science, the time has come for
a new application of the concept of personality in the world-view,
an application that can now take place with greater critique and with
clearer consciousness than in the mythological fantasies and in the theological
speculations. It is here particularly the concept of will he
relies upon. While he agrees with Hume that the causal relation's
objective validity cannot be demonstrated, so that we
must content ourselves with demonstrating at as many points in nature
as possible an unchangeable sequence, he separates himself from
Hume by asserting the indissoluble connection of the concept of cause
with the nature of our consciousness as one of the categories with which
it by its essence must work. And he separates himself still
more from Hume by teaching that the concept of cause has finally its
foundation in our consciousness of our willing. No relation between
cause and effect is intelligible to us (pénétrable) when
it is not interpreted as something analogous to our activity of will. The concept
of will makes the concept of force intelligible to us, but cannot itself
be defined by being traced back to something more primitive.
At this concept thought must stop --- by virtue of the law of relation.


Under these presuppositions it becomes easy for Renouvier
to conceive the origin and development of mental life in the world. The germs
of mental life must --- by virtue of the Monadology --- have
% --- p. 75 ---
\opage{75}been present from the first, so that only favorable
conditions are needed for their unfolding. To explain to himself the origin
of evil he assumes a very early disharmony, a
kind of fall into sin in the world of monads, which the Godhead, itself
limited and having limited itself, has not been able to prevent. In
his theological representations he lays weight on the necessity
of conceiving the Godhead in analogy with man and
therefore as finite and limited. Only then can an ethical world-view
be maintained, whereas the doctrine of infinity makes
all oppositions fade away, the ethical ones too.

While several thinkers (Bruno, Böhme, Spinoza), faced with the
expansion of the world-view to which the Copernican doctrine
gave the impetus, found refuge in the idea of infinity and embraced
it with enthusiasm, Renouvier insists that, if one
wishes to hold fast to the unconditional significance of the ethical ideas,
one must make man the center of the world, just as well as
Aristotle and the Bible did: l'anthropocentrisme est le
point de vue moral de l'univers. The external world is in size
and physical force far superior to us. But our earth could
perhaps nevertheless be the seat of something in spiritual respects far
more significant than the rest of the world, and when it is a matter of spiritual
power, it would then rather be we who with our superiority
(supériorité d'agents intelligents) could crush the universe,
than the reverse. Perfection consists always in limitation, and
it is therefore a fancy, leading away from true greatness,
when one loses oneself in the unboundedness of the material world.
--- This allusion to a well-known utterance of Pascal
is characteristic of Renouvier's standpoint, and it appears
undeniably in him with greater consistency than the
corresponding sentence in Pascal. In general it is characteristic
of Renouvier that even where he appears more as mythologist
or theologian than as philosopher, he forms his views
with sharp logic and makes no compromise. He is at the same time clear
that he moves in the domain of bold conjectures and hypotheses.
``I do not dogmatize'', he says, ``but
seek to comprehend. What I really have in view in
my bold guesses is to illuminate the sublime possibilities
% --- p. 76 ---
\opage{76}. . . . Where it is a matter of tracking down truths that lie beyond
experience, indeed that are first to explain experience itself,
the hypothesis, without other control than logic and morality,
is the philosopher's only way out. So long as he remains in agreement
with the laws of pure reason, he has the right to assert
the probability of the propositions on the origin and nature
of our world that accord with the law of his understanding and with his
feeling about life''. ---

Only a single, but certainly very incisive critical
remark have I to make before I leave the philosophical
patriarch of our time. He bases the assumption of a first cause,
and with it his whole peculiar philosophy of religion, on the law of relation
(loi de relativité). This law requires, he holds, that we
set a limit to the constant regress from effect
to cause. But the term at which we stop then comes to denote
a limit to the law of relation. It is supposed to stand in relation
to that which is derived from it, but not itself on the other
side be derived from anything else, as is nevertheless the case
with all the earlier terms in the series. The law of relation is thus abolished
when we have come back to the first term. This is
a decided self-contradiction when one, like Renouvier, conceives
the law of relation as a law necessary for our thought, deeply grounded
in the very nature of consciousness. The law of relation requires that we
at every point connect term to term, and where we cease
to do so, our understanding too ceases. The law of relation leads to that extent
rather to a doctrine of infinity, in that at no
point may we suppose that we have reached an absolute starting point.
Renouvier is right in saying that the reference to
an infinite process gives no explanation. But the point
is that where we face a series that by virtue of its
law can be continued indefinitely, we can reach no other understanding
than that which lies in the law itself\textsuperscript{26}). There is here a mythological
tendency in Renouvier that gets the upper hand over his
criticism. Instead of immersing himself in the inner connection
between the phenomena he wants to land in a first beginning
and gives us then only the choice between assuming a ``thing''
or a ``person'' as this first beginning. By virtue of
% --- p. 77 ---
\opage{77}the law of relation itself the critical philosophy must reject this dilemma.
The philosophy of continuity will, precisely when the need for understanding
is stressed, always retain an advantage over the philosophy of discontinuity.


5. While Renouvier works with sharply cut postulates
and with energy proclaims the great breaks in continuity,
\textsc{Emile Boutroux} (b. 1845, professor at the Sorbonne)
seeks more a purely theoretical grounding of the philosophy of discontinuity
and turns his attention especially to the small
differences. In Boutroux's philosophizing is traced the continuation of
a tendency that in an earlier period of French philosophy was represented
by Maine de Biran (who is discussed in the second volume of
Den nyere Filosofis Historie).

In two significant treatises \textit{(De la contingence des
lois de la nature} 1875, and \textit{De l'idée de la loi naturelle dans
la science et dans la philosophie contemporaines} 1895) Boutroux
discusses the principle of causation in the different forms in which
the different sciences are able to apply it. His main purpose
is to demonstrate how much is lacking in its carrying through,
when it is conceived in the strict form in which the
mechanical conception of nature applies it. Especially at the transition
from one domain of experience to another we face relations of discontinuity
--- thus at the transition from logic and mathematics
to mechanics, from mechanics to physics, and so on through
chemistry, physiology, psychology, sociology. The more concrete
sciences' fundamental laws cannot be reduced to those of the more abstract;
every new domain of experience requires new, special principles
that are not contained in those of the preceding. Boutroux is
here in remarkable agreement with Comte, who had already
strongly stressed the discontinuity of principles at the transition
from one science to another and had expressly warned
against attempting to derive the principles of a more concrete science
from those of a more abstract. Boutroux himself has been aware
how he here in a certain sense meets the founder of
positivism. In a spirited characterization of Comte,
which served as an introduction to a discussion in the French philosophical
society, he has moreover brought out in particular the
% --- p. 78 ---
\opage{78}elements in Comte's philosophy that point beyond the strictly empirical
standpoint it was Comte's endeavor to occupy.
\textit{(Bulletin de la société française de philosophie.} Janvier 1903).

Peculiar to the concrete sciences, according to
Boutroux, is that they concern wholes that are thought to consist of elements,
without it being possible, however, to demonstrate a relation of equivalence
between such a whole and the elements of which it is supposed
to consist. There is, e.g., no relation of equivalence between
a human being as a whole and the elements that can be demonstrated
in his nature. Therefore we are here always referred to experience;
only it can teach us what wholes are formed by
the interaction of elements. Here it then appears clearly that
the principle of causation is an abstract law that does not exhaust the whole
concrete reality. In actual development there appear
unforeseeable, and to that extent contingent changes (variations
contingentes). The history of things contains for us the key
to their essence, but cannot be derived from their nature. The
dynamic point of view gains, the more we go from the abstract to
the concrete sciences, the preponderance over the static, and qualities
rule instead of quantities. The laws of nature indicate
only the constant relations in which the changes take place,
the habits, so to speak, that things display. Change
is the principle, permanence is only a result. The constant
in nature, the so-called laws of nature, is the riverbed in which the stream of
events flows, --- and which the stream itself has deepened,
although it now follows it. If deeper-going variations
occurred in the things themselves, this riverbed might
change its form and its direction; the objective, logical relations themselves
would then undergo change.

When the changes through which the development from the elementary
to the higher forms of existence takes place cannot
be brought into a relation of equivalence to the preceding states, chance
to that extent rules. But this can be interpreted as the
external, empirical sign of freedom. Freedom itself cannot be immediately
experienced, but it expresses itself in human action
each time a habit is overcome. In the psychological domain
the laws concern only what the soul leaves to habit to accomplish.
% --- p. 79 ---
\opage{79}Habit has here, as everywhere in nature, the significance that through
the mechanization of the most necessary functions
force can be set free for new work.

It is thus for Boutroux the free development through
new formations that is the real working and effective thing in
existence, the real existence. The constant forms and
laws are only so-called results. If, when he wrote
his spirited treatises, he had known Hugo de Vries's
``Mutationsteori'', he could have claimed it in support
of his view, in that the said theory assumes a leaping
varying and emerging of new types. With Lamarck's theory
he evidently sympathizes, in that it assumed an inner need
and striving in living beings, leading them to adapt
themselves to conditions and thus to develop fixed forms.

But in his work on the concept of the law of nature in modern
science Boutroux lays weight also on another point of view
about the constant in nature. Epistemologically regarded
what we call the laws of nature are the sum of the methods we
have found for making things accessible to the understanding
(assimiler les choses à notre intelligence) and for using them
to attain our purposes. And what thus holds of the concept
of law of nature holds of all the concepts we form in order to understand
things, in particular the distinctions we make in order to see
more clearly. When we thus distinguish so sharply between thought and
motion, and perhaps assume these two domains to be separated
by an abyss, we have no right to ascribe to this distinction
absolute significance. Our concepts and conceptual differences
express rather our way of regarding things than the things'
own way of being. --- By these last remarks
Boutroux approaches the economic epistemology
we shall encounter in the following part. Had
he carried through the epistemological way of regarding things
he here hints at, he would perhaps have become convinced
that the ``chance'' that he interprets as the expression of ``freedom''
really only denotes a negative limit to inquiry,
which is perhaps not definitive, but sets new tasks. The philosophy of discontinuity
becomes dogmatic when it takes the
% --- p. 80 ---
\opage{80}differences or leaps at which we provisionally stop to be absolute.
It must remain constantly a problem how the
new variations have come to be. That they occur for our observation
suddenly they have in common with many events,
--- perhaps even, when one looks closely, with all events
in the world.

% --- p. 81 ---
\opage{81}\grouphead{SECOND GROUP.}{THE EPISTEMOLOGICAL-BIOLOGICAL\\ TENDENCY.}

Epistemological considerations were by no means foreign to the group of thinkers with whom we have been
concerned in what precedes. In the end, however, what mattered to them was the question whether any of
the special points of view which the individual sciences each apply in their own domain are suited to serve
as the principle of a coherent and provisionally conclusive world-view. But the difficulties which the
answering of this question met with could not but awaken the need for a renewed investigation of the
scientific points of view and principles, their origin and the conditions of their validity. The character
of the contribution which natural science can make to a general world-view was taken up for renewed
discussion --- not least within the circle of natural scientists themselves. The so-called mechanical
conception of nature, which was founded by Galileo and Descartes and which traced everything back to
pressure, impact and attraction, had been regarded with some inclination as the highest truth, as more
than a hypothesis. New phenomena now came forward (notably through the electromagnetic theory of light)
which were difficult to bring under the mechanical conception of nature, and moreover it was also impressed
upon us that a theory must not be held to be absolutely correct because it is confirmed by
% --- p. 82 ---
\opage{82}experience, so long as the possibility of another theory, whose consequences could likewise be
confirmed by experience, is not excluded.

The emphasis laid upon development works in the same direction, upon the constant extension of the
limits of experience and a possible change in the inner structure of our spirit. The dynamic thereby gains
the upper hand over the static, and our theory of knowledge will be able to hold good only for a certain
stage of development.

When, finally, the element of will is so strongly emphasized, particularly in the very knowledge of the
principles, it is natural to discuss the question what goal it really is that we set ourselves in our
knowledge, and what means we have at our disposal for approaching it. ---

By these various paths we are led over to the group of thinkers whom we have designated as the
epistemological-biological. To this group I reckon, first, some famous natural scientists who, partly
through their special researches, partly through the study of the history of their science, have been led
to reflect upon the presuppositions of knowledge, and, further, a philosopher who has undertaken a more
thoroughgoing investigation of the arising and vanishing of problems in the course of life and development.

\chapterhead{I.}{PHILOSOPHIZING NATURAL SCIENTISTS.}

1. The distinguished physicist \textsc{James Clerk Maxwell} (b. in Edinburgh 1831, d. as Professor of
Physics at Cambridge 1879), famous above all as the man who developed the electromagnetic theory of light
in mathematical form, acquires significance for the history of philosophy through the interest he took in
making clear to himself the first principles of his science. It was William Hamilton, whose lectures he
heard as a student in Edinburgh, who laid the foundation of his philosophical interest; he received there
(according to his biographers) an impulse which never lost its effect. Hamilton, for his part, also felt
great interest in the young, acute and eager-to-learn student. Maxwell's conception of the essence and
significance of thinking even later on, when he could build it upon his own profound
% --- p. 83 ---
\opage{83}inquiry, recalls Hamilton's philosophy on decisive points. The works of Maxwell which are of
particular importance are found in the second volume of his \textit{Scientific Papers,} which was collected
and published after his death.

It has been asserted, says Maxwell, that metaphysical speculation ought to belong only to the past, and
that physics ought to have eradicated it. Yet there is no danger that the discussion of the fundamental
concepts of existence will be brought to a close in our time. The exercise of the speculative faculty will
remain as attractive to every fresh spirit as it was in the time of Thales! --- But what matters is to
develop thoughts with which we can work. The progress of the exact sciences rests upon the discovery and
development of suitable (appropriate) and accurate ideas, by whose help we can give a mental
re-presentation (representation) of the facts, --- a re-presentation which on the one hand is sufficiently
general to be able to represent (stand for) any special case, and on the other hand sufficiently accurate
to form the basis for mathematical inferences. This is the road that has been followed, from Euclid to
Faraday.

When it is asked how we come in the first place to set up and apply a principle in a definite domain,
Maxwell points to the significance of analogy for science. By \textit{physical analogy} is understood the
partial similarity between the laws of two domains of experience, which makes it possible for the one
domain to illuminate the other. All application of mathematics in science is accordingly grounded on an
analogy between the laws of physical magnitudes and the laws of numbers, and the endeavor of the exact
sciences aims at tracing the problems of nature back to the determination of magnitudes by means of
operations with numbers\textsuperscript{27}). And of a similar character is really the whole so-called
mechanical conception of nature. It rests on an analogy between the laws of the qualitative changes in
nature and the laws of motion. Thereby it becomes possible to describe the changes of bodies as if they
consisted merely in motions of small parts. Such a description of a phenomenon, which makes it possible to
see it as an example of a principle that can also be
% --- p. 84 ---
\opage{84}applied to other phenomena, is an explanation, and experience shows that with every such
explanation there follows a peculiar mental satisfaction. So long as the analogy holds, one need not form
definite ideas of the nature of those small parts, nor even attribute to them extension, shape, hardness
or softness. What must be presupposed about the ultimate parts of matter is only constancy and regularity
of their motions. It is the constancy and continuity of motion that underlies the concept of matter.

On this last point Maxwell joins thoughts which had earlier been advanced by Faraday and W. Thomson (Lord
Kelvin). \textit{Only through motion,} he maintains, \textit{does one understand rest and equilibrium;}
therefore kinematics (the doctrine of motion, abstracted from the nature of that which moves) must precede
statics (the doctrine of relations of equilibrium). Kinematics differs from geometry merely in that time
is expressly introduced as a measurable quantity. Yet at bottom geometry itself is derived in relation to
the doctrine of motion, or is a part of it; for geometry really deals with the process by which figures in
space are produced. A line is not originally a stroke on the blackboard, which may just as well be called
A B as B A, but is the path of a motion from A to B. The idea of motion underlies the idea of form. That
this is so is also shown by the psychological doctrine of attention: this is at any given time limited to
a single undivided idea, but moves forward through a continuous series of ideas. One keeps as a rule too
much to the finished signs and figures and overlooks the genetic process of sensation and thinking.

The principles of knowledge thus arise from analogy and are applied to change before rest. A third point
has already been indicated: the truth of principles consists in their validity, and their validity in
their \textit{applicability,} their capacity to set inquiry in motion. The significance of a principle
consists in its leading us to put quite definite questions. Thus the significance of the energy proposition
consists in this, that all physical phenomena can be regarded as examples of transformations of
% --- p. 85 ---
\opage{85}energy, so that with every new phenomenon one can ask: what transformation of energy takes place
here? from where and to where? what are the conditions of the transformation? Thus the assumption of action
at a distance, even if it --- as Maxwell himself believes --- should be incorrect, has had the great
scientific significance that, in order to answer the question how bodies can act upon one another at a
distance, people have investigated with zeal the properties of the intervening medium. --- One is
reminded here of William Hamilton's proposition that a living untruth is better than a dead truth. ---
Maxwell even applies this point of view to the very doctrine of atoms, the presupposition of the
mechanical conception of nature. After showing that none of the causes which we call natural can explain
the existence of atoms and in particular their uniformity, he declares: ``Is it then so that our
scientific speculations have already now reached down below the visible manifestation of things, where
arising and passing away prevail, and stand at the entrance to that world of order and perfection which to
this very day is as it was created, perfect in number and measure and weight? We are perhaps mistaken. No
one has yet seen or felt an atom, and our atomic hypothesis will perhaps be supplanted by a new theory of
the constitution of matter. Yet the idea of innumerable single things that are alike and unchangeable has
not arisen in human consciousness without bearing fruit.'' Thus: even if the doctrine of atoms is not
``correct,'' it may very well have been useful for a long period of time. In an analogous way Maxwell
expresses himself about the inference that, since a dissipation of physical energy is constantly taking
place, there must have been an initial state which cannot have come about in a natural way. It is not
said, he remarks, that this hypothesis is the last word of science, --- and moreover dissipated energy
really means only energy for which we cannot imagine any use. The case here, he adds, is like that when we
speak of disturbances (disturbations), which prevent a law of nature from emerging clearly: they are in
reality circumstances which we do not know, or have overlooked, or must reserve for later investigation;
``disturbance'' is a fancy, not a fact of nature;
% --- p. 86 ---
\opage{86}in the working of nature there is no disturbance. --- One seems to hear Spinoza speak.

Yet there are other places where Maxwell is not so methodical and consistent. As with not a few natural
scientists, the theological ideas sometimes lie so ready to hand with him that they step forward when he
believes he is standing at an absolute limit, or perhaps rather that they mislead him into believing that
he is standing at such a limit. Just as Descartes derived the highest laws of nature from the will of God,
--- just as Newton declared that the coherence of the solar system could not have a natural origin, ---
and just as Linnaeus and Cuvier declared the same of the organic types, and Lord Kelvin in the most recent
time of the very origin of life\textsuperscript{28}), so Maxwell teaches in some places (which cannot be
reconciled with the passage cited above) that the uniform and unchangeable atoms, which quite bear the
stamp of manufactured articles, must have been created; we are here at the limit of all science, he adds,
for science cannot discuss creation out of nothing. The coming into being of an atom is an event that
cannot have taken place within the natural order under which we now live. But such an event must have taken
place, since it would be absurd to think of absolutely uniform elements existing from eternity. --- In
what this absurdity is really supposed to lie, Maxwell does not say, and his doctrine of atoms has here
evidently a more dogmatic character than in other utterances, which agree better with the intellectual
energy of this fine inquirer.

On another point too Maxwell found a limit to science.

Although, he says, the soul has long been called immaterial, the presupposition nevertheless lay at the
bottom, right up to the most recent time, that one would indeed not be able to come upon the soul by
anatomical or chemical means (by dissection or by weighing) in an investigation of the brain, but would,
by going back far enough, come upon a material motion which had no material cause but was directly
brought about by a soul. This assumption conflicts, according to Maxwell, with the law of the conservation
of energy, and he asserts that one would not trust
% --- p. 87 ---
\opage{87}an experiment which showed a considerable difference between the work performed by a living being
and the energy received and consumed. But how, then, can consciousness have come into being? It is of no
use to attribute psychical properties to the atoms, for consciousness cannot be explained as a product of
psychical atoms. When we come to the depths of personality, we are beyond the limit of science. The
objective facts cease. The very first proposition that holds of the subject: ``I am!'' cannot be taken by
two persons in quite the same sense and can therefore not become an object of science at all.

On this last utterance it must be remarked that when one can discover the difference between two persons'
use of the same thought, one must be able to observe and compare their thoughts, and thereby an empirical
psychology has already been introduced. Here too Maxwell stops too early, which is all the more remarkable
as he himself in an earlier cited utterance appeals to psychological observation.

2. In \textsc{Ernst Mach} (b. 1838) a very early awakened philosophical interest worked together with
inquiry in natural science, and especially with the study of the history of natural science, to put the
problem of knowledge in a new and peculiar way. This course of development shows itself outwardly in his
teaching career, in that for some years he was Professor of Physics (in Prague) and later became
Professor of Philosophy (in Vienna). His works concern in particular the history of natural science
\textit{(Die Geschichte und die Wurzel des Satzes von der Erhaltung der Arbeit.} 1872. ---
\textit{Die Mechanik in ihrer Entwickelung.} 2. Aufl. 1889. --- \textit{Die Principien der Wärmelehre.}
Historisch-kritisch erläutert. 1896). In \textit{Die Analyse der Empfindungen} (1886, 4. Aufl. 1903) he has
given a purely philosophical work.

Mach was no more than 15 years old when he got hold of Kant's ``Prolegomena'' in his father's library.
This writing made, according to what he has later stated, a tremendous impression upon him, and he
regarded it as a piece of good fortune that he had become acquainted so early with its line of thought. A
complete Kantian he did not continue to be; thus he found the ``Ding an sich'' superfluous. Moreover, the
later study of Spinoza, Herbart
% --- p. 88 ---
\opage{88}and Fechner had its effect on him. What mattered to him, as he has expressed it, was to take up
\textit{a standpoint which he need not give up when he passed from physics to psychology.} All sciences
must after all in the end constitute a whole. And Mach saw clearly that the doctrine of atoms, as it is
generally conceived, does not make possible the preservation of the same standpoint in the transition from
physics to psychology. One is more fortunately placed when from the very start one conceives all bodies as
thought-symbols for complexes of sensations. The world then does not consist of enigmatic beings in
interaction with an enigmatic ego; colors, tones, relations of space and time, etc., are only the ultimate
elements, whose connection it is our task to explore, and which we combine and delimit in the manner most
appropriate in practical and theoretical respects. There is only one great coherent mass of sensations,
most strongly connected in that which we call our ego, which thereby stands in opposition to that which we
call the world. But ego, world, spirit, matter and similar concepts are only thought-symbols which denote
limits and differences that it has proved more or less appropriate to hold fast.

This conception Mach already expressed in his lectures on psychophysics in 1863. Studies in the physiology
of the senses had on the whole a great influence on his line of thought. They convinced him, for example,
that the intuition of space is bound to the senses and therefore holds only for that which can be sensed,
so that one has no right to attribute spatial properties to the atoms, which cannot be sensed.

The task of science is to give an \textit{economical} presentation of the facts, i.e. a presentation in
which one uses only the strictly necessary and most appropriate ideas. In this conception of the task of
science Mach recognizes that he has a predecessor in Maxwell, whereas he believes that he has priority
over Kirchhoff, who espoused a related conception\textsuperscript{29}).

This way of regarding the nature of the sciences agrees well with the hypothesis of evolution. Mach points
out that Spencer, even before Darwin, regarded the mental processes
as \textit{adaptations} to the given. He himself has, supported by
% --- p. 89 ---
\opage{89}the history of the sciences, developed this biological point of view for knowledge in his
address \textit{Ueber Umbildung und Anpassung im naturwissenschaftlichen Denken} (1884). Adaptation
consists partly in combining, partly in reduction: partly in a connection of the manifold under one point
of view, partly in an elimination of that which is not necessary for understanding. Which facts are laid
at the foundation depends on convenience, tradition or habit. But the very facts that are laid at the
foundation for orientation cannot be ``understood.'' When, e.g., mechanical relations seem to us more
self-evident than others, this comes merely from our being more accustomed to them. Our understanding
always consists in unusual incomprehensibilities being traced back to usual incomprehensibilities. Of
course, in science every habitual judgment to which we attribute validity in a domain other than that in
which it was originally formed must be tested with regard to whether it fits the new case; otherwise it
becomes a prejudice. Our knowledge works its way forward continually through a struggle between judgments
and prejudices. Just as living beings do not immediately form a new organ for a new function that has
become necessary (as when a vertebrate is to fly or swim), but try to use the organs they already have, ---
so science uses a transformation of the older ideas for the presentation of new experiences. In Newton's
``attraction'' there still lurks something of the old idea of seeking the natural place. What matters is to
think the new in the simplest way and according to the principle of parsimony.

Arithmetic saves direct counting, in that the number itself expresses that two ways of ordering can be
identical even though that which is ordered is different. Maxwell had already, as Mach himself points out,
taught that the validity of the concept of number rests on analogy. --- Geometry investigates the mutual
relations of different scales and thereby saves us direct measurement. Visual space and tactile space are
different; but a displacement in the one corresponds to a displacement in the other. --- Quantitative
orderings are simpler and more comprehensive than qualitative orderings, and they facilitate the survey
and handling
% --- p. 90 ---
\opage{90}of large groups of experiences. Such orderings are made possible in natural science by means of
concepts like force, mass, atom. These concepts are only means of thought, whose significance consists
solely in recalling economically ordered experiences, although most natural scientists attribute reality
to them outside thinking. --- The causal relation means only the closest connection that any description
can express. It has only logical, not physical significance. The same holds of the principle of
continuity: we posit continuity where it turns out that the difference grows with the distance between
the members of a series, while it diminishes the more the members approach one another, and finally
vanishes for us. The advantage of applying the principle of continuity everywhere it is possible is that
even in the smallest distinguishable parts of the system concerned we can make considerations similar to
those for the larger parts. Only experience can decide how far this can be continued.

The economy of thinking under adaptation to experience entails, as is easily seen from what precedes, a
constant application of \textit{analogies.} In his treatise \textit{Die Aehnlichkeit und die Analogie als
Leitmotiv der Forschung} (in the first volume of Ostwald's ``Annalen der Naturphilosophie'') Mach treats
this concept. Analogy denotes such a relation between two systems of concepts that both the difference
between two corresponding concepts in the two systems and the agreement in their mutual relations come
clearly to light. It makes a unitary conception of dissimilar facts possible and is therefore of great
biological and epistemological significance, a significance to which the whole history of natural science
bears witness. But analogy is not identity, and the means of inquiry must be distinguished from the
results of inquiry.

\textit{The mechanical conception of nature} rests on a comprehensive analogy between the motion of masses
in space and the qualitative changes of things (in temperature, electrical relations, etc.). That this
analogy can be carried through to so great an extent gives no right to assume that all physical processes
are ``really'' mechanical. The mechanical laws can serve as
% --- p. 91 ---
\opage{91}formal patterns and pointers. The motions of bodies in space are the simplest and most visualizable
processes, --- those which are easiest to pursue in our imagination. Moreover every physical process,
whether it be heat, electricity, sound, etc., has a mechanical side, which appears as expansion,
trembling, attraction, etc. Therefore one can present and clarify physical processes clearly by means of
mechanical analogies, although one has no right to say (with Wundt, e.g.) that all physical causes are
causes of motion. Purely mechanical processes there are none at all, since magnetic and electrical
effects and heat processes always make themselves felt in addition to the mere motion. Every process
really belongs to all the domains of physics. ---

This conception of our knowledge is no dissolution of science. It deprives us of no really valuable
points of view, only the superfluous ones. And it removes the dogmatic that has crept in with many
inquirers.

The world-view which becomes possible on the basis of such a theory of knowledge cannot be materialism.
The qualitative character of sensations cannot be derived from purely quantitative relations and
connections, the psychical not from the physical. Physics and psychology really concern the same thing,
namely complexes of sensations; --- they merely treat it from different points of view.

But the world is not merely a sum of sensations. Definite functions hold between these, and the knowledge
of these functional relations is a knowledge of reality. The difference between idealism and realism is
theoretically a matter of indifference. We call the elements of existence sensations when we regard them
as elements in a psychical world. The task of science is only to find the lawful connection in everything
that happens; it therefore necessarily strives toward a \textit{monism.} ---

Perhaps Mach here gets over great difficulties too easily. The ``elements'' common to physics and
psychology stand somewhat indeterminate and mystical, like a mass of fog that is not shaped and
articulated. It may well be that the difference between physics and
% --- p. 92 ---
\opage{92}psychology is due to a division of labor which has proved appropriate. But such a division of
labor is not easily abolished again; it takes on the character of a ``prejudice'' which offers resistance
as long as quite definite experiences do not necessitate a new intellectual adaptation. And the very fact
that the division of labor has proved necessary and appropriate points to cosmological conditions.
Existence can hardly be as simple as Mach presupposes, when heterogeneous points of view are necessary in
order to understand it. Monism must be won through a struggle with greater oppositions than those Mach
acknowledges.

To this it must be added that the relation between the elements and their functional relations remains
unclear. It can hardly be a matter of indifference to the elements that they stand precisely in these
definite relations, and conversely the relations must surely be conditioned by the elements for which
they hold. But in what then does reality consist? In the mere functional relations alone it cannot
consist. Here, then, arises a whole series of questions about the relation of the particular to the
coherence, questions whose answer becomes of importance for the whole world-view.

That Mach has not pursued the problems so far does not diminish the value of his work, which has had an
awakening and fruitful effect to a rare degree and seems to attract increasing attention.

3. In his epistemological ideas \textsc{Heinrich} \textsc{Hertz} (b. 1857 in Hamburg, d. 1894 as Professor of
Physics in Bonn) joins Mach. The brilliant inquirer, who passed away so early, is especially famous for his
experimental demonstration of the identity of light and electricity, a discovery that confirmed Ørsted's
presentiment, Faraday's ideas and Maxwell's calculations. In an address at the German congress of natural
scientists in 1889 he gave in clear and accessible form an account of his discoveries \textit{(Ueber die
Beziehungen zwischen Licht und Elektrizität).} His teacher Helmholtz has, in the introduction to the third
volume of Hertz's writings, given a characterization of the distinguished inquirer and of his scientific
course of development. Here we have to do only with the manner in which Hertz conceived the
% --- p. 93 ---
\opage{93}first principles of his science, and the characterization he gives of the work of knowledge. He has expressed himself on this in a clear and weighty treatise, the introduction to his ``Principien der Mechanik'' (Werke III).

The task of science is to derive the future from the past. In order to be able to do this we form images or
symbols (Scheinbilder oder Symbole) of such a nature that their consequences necessary in thought are
always again images of the consequences \textit{necessary in nature} of the objects represented. Experience
shows that this is possible, which points to an agreement between nature and our spirit. But it must
surely be held fast that our images need agree with the real only in this one respect, namely the
agreement between necessity in thought and necessity in nature. Whether they agree with it otherwise, we
can never come to know.

If it is now possible to form several images that are equally logical and consistently applicable, we
prefer the simplest. To be sure, it will not be possible --- since the nature of our own spirit always
shows itself in the formation of images --- to exclude unnecessary and inapplicable elements. Simplicity
can always become greater, and applicability must always be tested anew. Whether nature itself follows the
simplest paths cannot be proved.

For his presentation of mechanics Hertz believes he needs no other symbols than time, space and mass,
whereas such concepts as force (energy) and atom are superfluous. Neither forces nor atoms can we observe.
What we see is always masses in motion. Now what we see does not present full lawfulness and coherence. We
must therefore assume that the manifoldness of the real world is greater than the manifoldness of the world
that reveals itself to our senses. If we want a rounded, lawful picture of the world, we must suppose
behind the things we see other, invisible things. But the simplest will then be to represent this hidden
to ourselves as masses in motion, differing from the visible masses only in being inaccessible to our
senses and our usual means of observation. We speak of force (energy) only when it is a matter of
% --- p. 94 ---
\opage{94}determining the relation between two motions; by force is understood partly the cause of a
motion, partly the effect of a motion: it is thus a derived concept. The concept of atom too is a pure
abstraction; what is given is always only a material system, never the single mass point. ---

It must of course be a matter for mechanics whether it can content itself with the presuppositions at
which Hertz stops. Here it interests us to see how an outstanding natural scientist has sought to give
himself an account of his intellectual work. That this account is determined by his own special field of
work cannot, after all, be otherwise. Hertz's great discovery was that electrical oscillations present
quite the same relations as light and the rays of heat, so that in space far more motions must take place
than our senses directly show us. The invisible masses and motions thereby came quite naturally to play a
leading role in Hertz's scientific imagination; in this respect he had a great predecessor in Faraday, who
already strove to exclude everything that could not be visualized, and therefore doubted the reality of
empty space and of force at a distance. --- Whether it is possible entirely to exclude the concept of
force (the concept of energy) is doubtful; and really Hertz too wants only to reduce it to a subordinate,
derived concept. Perhaps he would have come to see the matter differently had he gone more closely into
the problem of causality. The proposition of causality he sets up merely as a necessary presupposition,
--- in the form of the basic assumption that there exist fixed, lawful relations between the different
positions of material systems. But the problem of causality is posed sharply and pointedly only when we
stop at the single moment in which the preceding state is about to vanish and the following has not yet
set in. Here the concept of force denotes an expectation of what will happen, a trust that what is to come
will stand in a lawful relation to the past. Hertz himself also says that we speak of force \textit{while}
the process by which we derive future experiences from the preceding is going on. (Einleitung p. 14). But
in a certain sense we always live in such moments --- and the ``auxiliary concept'' is
% --- p. 95 ---
\opage{95}thus constantly needed, the more intensively the intellectual work is\textsuperscript{30}).

From Mach Hertz differs in emphasizing the symbolic more than the economical in our basic concepts. It is
the need for visualizable images that makes itself felt especially in him, and when he would gladly do
without concepts like force and atom, it was perhaps because they did not make clear visualization
possible. Either Hertz has been markedly visual, or else it was decisive for him that visual symbols have
the advantage over motor (muscular) symbols that they permit direct measurement. Force is a symbol that
points back to the experience of muscular exertion; but this can be measured only through the very motion
it produces\textsuperscript{31}).

Of greatest interest it seems to me that Hertz's definition of the truth of scientific symbols agrees
closely with the only definition that can be given of the truth of our sense qualities, and which his
teacher \textsc{Helmholtz} had clearly set forth in his ``Physiologische Optik'' (§ 26) in the following
words: ``Insofar as the quality of our sensation informs us of the peculiarity of the external action by
which it was evoked, it may count as a \textit{sign} of this, but not as a \textit{copy} of it ... A sign
need not have any kind of similarity with that of which it is the sign. The relation between them is only
that the same object under the same conditions evokes the same sign.... We call our ideas of the external
world \textit{true} when they give us sufficient direction about the consequences of our actions in the
external world, and let us draw correct inferences about the changes that can be expected in it.''\textsuperscript{32})

4. In a peculiar contrast to Hertz stands \textsc{Wilhelm} \textsc{Ostwald}, the famous chemist, in that he
precisely wants to make energy, the concept Hertz wants to push aside, the basic concept. Perhaps this
contrast is connected with Ostwald's being motor, whereas Hertz seems to be visual. But it may also be
connected with their fields of study, in that chemists have more occasion than physicists to think of
relations of energy, whereas physicists find it easier to carry through
% --- p. 96 ---
\opage{96}visual schemes than chemists, for whom they take on a more pronounced hypothetical character.

Ostwald (b. 1853, since 1887 Professor of Chemistry in Leipzig) wages a war against all materialism, in
that he wants to reduce all matter to energy. This thought first appears in his address at the German
congress of natural scientists in 1895 \textit{(Die Ueberwindung des wissenschaftlichen Materialismus),}
and is elaborately developed in his \textit{Naturphilosophie} (1902), dedicated to Mach. Ostwald has
furthermore sought to work for the discussion of the questions that lie on the border between natural
science and philosophy, by founding a new journal \textit{Annalen der Naturphilosophie.}

Everything is energy, and nothing exists but energy! That is Ostwald's main proposition. All properties of
matter can be resolved into energy: mass is capacity for energy of motion; space-filling is volume energy;
weight is a certain kind of energy of position; the chemical properties are different kinds of the energy
that appears in the transformations of substances. Energy is work, or everything that arises from work or
is transformed into work. It is the most comprehensive concept that science has so far formed. It
comprises both persistence (substantiality) and action (causality). If by substance one means that which
persists, energy is the true substance. It streams out from the sun over the earth, is used immediately or
is accumulated as chemical energy to be consumed later in vital activity.

Concepts like matter and body we have occasion to form only when different combinations of energy bring
about a composite state of equilibrium. If something is to \textit{happen}, there must be differences of
intensity present between the different energies. In reality there are only approximations to complete
equilibrium. --- The proposition of the persistence of matter or of the material elements means only that
the elements can always be produced anew after any composition into which they have entered.

Since Ostwald has wanted to give a philosophy of nature, not merely a theory of knowledge, he seeks to
bring the life of consciousness under the same point of view as that from which he regards the whole
% --- p. 97 ---
\opage{97}organic and inorganic nature. Consciousness too is, according to him, energy. That we represent
everything external as energy finds its very explanation in the fact that all processes of consciousness
are of an energetic kind. This Ostwald seeks to establish by an investigation of attention, memory,
comparison and willing. And he refers to the synthetic character of consciousness, which Kant so strongly
emphasized, and which expresses itself in all psychical lived experiences standing in inner connection as
having arisen in one and the same brain or spirit. The energy of consciousness expresses itself in this
combining. Since our ideas of the external world are determined by the nature of our consciousness, it is
no wonder that they bear a similar stamp to consciousness itself. The significance of the phenomena of
consciousness consists in their making accumulation of experiences possible, so that earlier lived
experiences can be used for comparison with new ones.

With this Ostwald believes he has solved the problem of the relation between the physical and the
psychical. But when he defines spiritual energy as ``unconscious \textit{or} conscious nerve energy,'' it
is clear that the problem constantly poses itself anew, since one must ask how the transition takes place
from unconscious to conscious nerve energy. It is of no use to answer that it takes place through
attention, for if this is presupposed to be bound up with consciousness, we move in a circle, and if it is
not presupposed to be conscious, it is not easy to understand how consciousness can arise through the
cooperation of two unconscious processes. The concept of nerve energy remains in Ostwald a purely
mystical concept.

This is connected with the fact that Ostwald's doctrine of energy is built first and foremost on
experiences from the outer world. Energy he defines as that which follows from work and is converted into
work. But the work spoken of here is the overcoming of resistance to a motion. Only apparently can Ostwald
therefore exclude the geometrical properties from his concept of energy; they are just as original as the
purely dynamic ones. Mass and energy are correlative concepts, which must be determined reciprocally by
one another. Mass is known by the constant relation between energy and increase of velocity, and
% --- p. 98 ---
\opage{98}energy is known when one presupposes mass and increase of velocity as known. Increase of velocity
is known by measurement of the relations of space and time. We thus come by the physical
concept of energy no further than the spatial. Just as little as Hertz and Mach succeeded in excluding the
concept of force, did Ostwald succeed in excluding the concept of matter. --- Up to now it has succeeded
neither for him nor for anyone else to set up a concept of energy that could be laid at the foundation
both of the concept of psychical energy, which expresses itself in the combining and distinguishing of
elements of consciousness, and of the concept of physical energy, which expresses itself in the overcoming
of a resistance in space.

\chapterhead{II.}{THE NATURAL HISTORY OF PROBLEMS.}

By a philosophical path \textsc{Richard Avenarius} came to a conception of knowledge which is related to
that which Maxwell and Hertz arrived at by reflecting on the principles of their science, and which Mach
arrived at through the study of the history of natural science.

Avenarius (b. 1843) studied in Leipzig, where he was introduced by the physiologist Ludwig to the strict
mechanical conception of nature in its application to organic phenomena, and by the philosopher Drobisch
to the Herbartian philosophy with its doctrine of ideas as expressions of the soul's instinct of
self-preservation. The study of Spinoza became of great importance to him by confronting him with a
magnificent attempt at strict systematic reduction of all thoughts to a single one. His first writing
dealt with Spinoza's philosophy (\textit{Ueber die beiden ersten Phasen des Spinozischen Pantheismus.}
1868). He sought here in an acute manner to show how Spinoza's system has come about through the fusion of
three different series of thought, a religious one, expressed in the concept of God, a natural-scientific
one, expressed in the concept of nature, and an abstract metaphysical one, expressed in the concept of
substance. The three concepts express in the end for Spinoza one and the same thing (deus = natura =
substantia). --- Through the reduction of the organic to mechanics, through the conception of ideas
% --- p. 99 ---
\opage{99}as expressions of self-preservation, and through the study of the most thoroughgoing attempt at a
doctrine of identity, there was prepared in Avenarius the conception of knowledge which he first
expressed in his treatise \textit{Philosophie als Denken der Welt gemäss dem Princip des kleinsten
Kraftmasses.} Prolegomena zu einer Kritik der reinen Erfahrung. (1876). At the bottom here lies the fact
that consciousness does not command an infinite power of ideation; therefore it must apply economy in its
thinking. It strives then to trace the unknown back to the known; this takes place in all recognition and
in all grasping. And it strives in general to reduce differences and manifoldness to the least possible.
Philosophy is precisely the scientific striving to think the sum total of what is given in experience with
the least possible expenditure of force. It strives to attain pure experience by excluding from knowledge
all those ideas which are not contained in the given itself; only when this is attained is no more force
expended on thinking the given than this itself demands.

In 1877 Avenarius became Professor of Philosophy in Zurich, and here he worked out his chief work
\textit{Kritik der reinen Erfahrung} (1888---1890). It is a doctrine of problems which, by a purely
biological and psychological path, seeks to demonstrate the conditions for the arising and vanishing of
problems. With regard to the problems Avenarius wants to place himself at the standpoint of the
psychophysical experimenter or the psychiatrist and thus to study the problems purely
naturalistically. He believes he can establish a constant approximation to ``pure experience'' through an
increasing purification of all elements of thought that are not contained in the given itself. It is an
approximation to a purely empirical concept of the world. What such a concept of the world would be like
he has set forth in his writing \textit{Der menschliche Weltbegriff} (1891), which however mainly
supplements the chief work by seeking the ground of the ``impurity'' of experience in animism and in
general in the need to conceive things in analogy with ourselves.

Avenarius' work is the fruit of earnest thinking and built upon a wealth of observations and examples. But
he has damaged its effect by the introduction of a very artificial and
% --- p. 100 ---
\opage{100}unnecessary terminology, which has a repellent effect on many readers. There is something tragic
in the fact that his considerable work of thought is thereby hindered from winning the recognition it
deserves. His health was broken down under his zealous studies. In vain he sought relief at health resorts,
for instance at Skodsborg, where I --- after having earlier visited him in Zurich --- on walks with him in
the Deer Park was for the first time initiated into ``pure experience.'' A couple of years later he died
(1896). A rare energy of thought united with artistic interest and an open, friendly disposition are the
main features of the image which those who came into personal contact with him preserve in their memory.

In the characterization of Avenarius' philosophy I shall not follow his own presentation, nor apply his
terminology. By repeated study of his chief work it has become clear to me that his fundamental thoughts
can be presented in a far simpler manner than he has done, --- and that only thereby do they come into their
full right.

It is the natural history of problems that Avenarius wants to give, in that by a physiological,
psychological and historical path he wants to show under what conditions they arise, are sharpened, are
solved or vanish. The material for this investigation can be fetched from everywhere, for the child, the
savage, the practical man, the religious believer have their problems just as well as the man of science
and the philosopher. All stand alike here, since it is a matter of something that follows from the natural,
universally human conditions. The difference is only that scientific posing and solution of problems comes
closest to pure experience, i.e. to re-presentation of the given without subjective additions.

A problem is always --- wherever it arises --- a sign of a relation of tension (a ``vital difference'')
between the individual and its environment, a relation of tension that arises through the influences that
come from the environment demanding either greater or less energy than the individual commands. --- When
influence (P) and energy (E) quite correspond to one another (which we can express by P=E), the highest
degree of the maintenance of life (``vitales Erhaltungsmaximum'') is present. There occurs recognition
% --- p. 101 ---
\opage{101}of the environment; the individual feels at home and has confidence in its observations and its
ideas. --- When now, by means of changes in the environment, new and greater work is demanded of the
individual than before, without the energy being increased or its manner of application altered, a problem
arises. (This relation can be expressed by P > E). The individual finds deviations, exceptions and
contradictions in the given; this stands as ``something other'' than it was before, perhaps as something
``wonderful''; security is abolished, and the individual feels a stranger in the world; recognition is for
the time being impossible. Every genuine problem is a kind of homesickness and releases a longing to
abolish the strangeness. With increasing culture, education and knowledge the possibility of such
strangeness increases. --- If conversely the energy increases more than the influences, or these decrease
(which can be expressed by E > P), a problem arises from the opposite ground: there is a lack of
application for the energy; this is then released freely, in unusual directions, which are not wholly
determined by the given; it yearns toward the distant; danger and suffering are sought out; there come
times of emancipation or of storm and stress, in general times of practical idealism. --- This last form
Avenarius has dwelt very little upon. As representatives Rousseau and other ``subjective'' thinkers might
be cited. Hume has in his psychological theory of knowledge laid great weight on this point, and --- as we
shall see --- Guyau's and Nietzsche's whole philosophy is built upon this relation.

The relation to which Avenarius especially holds is that in which the influences demand a different or
greater performance of work than, at least for the time being, corresponds to the individual's energy
(P > E). In the closer investigation of what this relation entails, Avenarius considers it most
scientific, i.e. most closely approaching pure experience, to adopt a physiological way of regarding, thus
to regard the processes that the posing of problems releases as processes in the central nervous system.
What an individual states about its conditions during the posing, treatment and solution of problems is
then regarded as symptoms of what goes on in the central nervous system. We thus obtain
% --- p. 102 ---
\opage{102}two series (``vital series''), a subjective one: the states of the individual, expressed in what is
stated, and an objective one: the changes in the central nervous system. The former is a mathematical
function of the latter, and the subjective series can therefore be designated as the dependent vital
series, the objective series as the independent vital series. It is merely for reasons of expediency that
the objective series stands as the ``independent'' one. The relation between the two vital series could in
and of itself be reversed, and Avenarius' definitive standpoint is designated by an oral utterance of his:
``Ich kenne weder Physisches noch Psychisches, sondern nur ein Drittes''. --- On closer examination of his
presentation it nevertheless turns out to be an illusion when he believes he is laying the objective vital
series at the foundation. The first volume of his work is to present it, and the second volume the
subjective series. The presentation in the first volume now has the peculiarity that no special examples
are found; these come first in the second volume and are all of a ``subjective'' kind. The greatest wealth
of facts is thus found in the subjective series, and it is only a very meager and purely schematic
presentation that can be given of the objective series; even this meager scheme becomes intelligible only
when the psychical correlates are sought out. This relation is characteristic of the relation between
psychology and physiology and contains a warning against assuming that one attains greater
``scientificity'' by proclaiming the ``biological'' method as the one right method in the treatment of
questions such as those spoken of here\textsuperscript{33}). ---

A vital series is a process by which a new equilibrium is attained after an earlier state of equilibrium
has been abolished by a posing of a problem. It has three stages, which (if we keep to the subjective
series) can be designated as need, work and liberation. The first stage can in turn express itself in
different ways: as privation and longing, as desire, as doubt, as anxiety, pain, repentance or
disappointment. It must be remembered that we are speaking here of the problems and relations of tension
of practical life, as well as of pure theory. The second stage can appear as acting or striving, as struggle
% --- p. 103 ---
\opage{103}or steady activity, as a venturing in which the personality is put into play, or as a patient
working out of the conditions for solving a task. The third stage is the time of the feeling of victory and
of satisfaction.

The third stage (the solution, deproblematization) can be attained in different ways. If the tension is
less considerable, it can be abolished by restitution, in that either a change again takes place in the
environment by which the old state re-enters, or the individual makes a momentary adaptation. If the
tension is greater, a substitution must take place by the aid of habit or practice or by the development
of new states or capacities. The complete solution requires that functions be developed which correspond
to the constant and universal in the conditions that the environment can present, so that the individual
for the future becomes independent of the most essential changes. In restitution homesickness leads back
to the old home; in substitution it leads to a new home being built. It may be that this solution is
purely individual (as in a view of life that the single personality is alone in holding); but it can also
be social (as in a popular religion), or universal (as in science). The content of formal logic is the
forms in which it has proved that vital series can come to universal solution. The differences are here
reduced to the least possible, and the so-called principle of identity really indicates merely a demand
for such a reduction, so that a ``heterotic minimum'' is attained, i.e. the least possible heterogeneity
in the content of experience. What matters is a separation of those elements that have only individual
and transitory significance, and which derive from heredity or tradition, with a holding fast of those
elements that constantly find application. If those elements are called \textit{α}, these a, then what
matters in y = f (\textit{α}, a) is to make \textit{α} become o. Thereby the individual or social solution
passes over into becoming universal, and we approach pure experience. This takes place the more
observations are available and the more thoroughly they are thought through. Logically tenable in the end
will be only those elements that have universal significance (thus a, not \textit{α}).

% --- p. 104 ---
\opage{104}But there is a difference between logical and biological tenability. It may be that individuals,
peoples and ages cannot do without certain ideas, though these do not have universal validity (thus
decidedly belong to the class \textit{α}, not to the class a). There will then exist ``forms of
protection'' which are withdrawn from the control of experience. To this belong many elements in
individual views of life, in popular religions, in philosophical systems and in natural-scientific
hypotheses and principles. Such forms of protection will of course not be unchangeable; the biologically
indispensable will vary with the conditions. Perhaps their justification is supported directly on their
indispensability; their validity is then ``postulated.''

What is a problem for the one is not always so for the other, either because he believes he has a
solution, or because the conditions for the arising of the problem are not present in him. And what is a
problem for the one is perhaps precisely a solution for the other (e.g. creation out of nothing). The
purely subjective symptom (``certainty,'' ``self-evidence'') is no guarantee that a real solution (an
absolute deproblematization) has been attained, and that a new posing of problems (problematization) will
not set in. New vital series can constantly be released, until the constantly repeated influences lead to
a complete elimination of the unnecessary and untenable accessory ideas (\textit{α}), so that no other
and more energy is released than is precisely needed (a). There constantly takes place such a process of
elimination, by which we approach a heterotic minimum. The standpoint of pure experience will be
characterized by knowledge consisting essentially in a description, as far as possible by means of
quantitative, not qualitative determinations, with constant demonstration of equivalents and with
derivation of the following from the preceding. In the practical respect an analogous development takes
place, in that the most stable social life proves to be that within which the social differences are
reduced to the least possible. Thereby it becomes possible for the single part of society to assert
itself --- not at the expense of the others, but precisely by working for their maintenance. ---

Avenarius himself admits that pure experience can be attained only
% --- p. 105 ---
\opage{105}approximately. He ought certainly to have laid far greater weight on this approximative element;
thereby his conception of the problems would have become more fruitful. For with Avenarius it easily
comes to look as if it is a fault in us that we do not stand at the standpoint of pure experience; and
that it cannot be, when it is a standpoint that can never be fully attained. Pure experience stands with
Avenarius as the pure idea in Plato. Our real knowledge is constantly at work to attain the pure idea or
pure experience. But we must go a step further. Can even pure experience do without auxiliary concepts
(\textit{α}), if it is still to remain experience? We cannot after all make ourselves one with the
``environment''; this continues, even if our experience becomes ever so ``pure'' of subjective elements,
nevertheless to ``surround'' us; the difference between it and us does not fall away. And already the fact
that something is \textit{given} implies that subjective elements are present. Avenarius originally
designated his task as an attempt to give a philosophy as thinking of the world with the least possible
expenditure of force; it is such an attempt that in his opinion leads to ``pure experience.'' But the
economy of thinking necessarily leads to symbolism, since one cannot demonstrate that the economic
consideration (the principle of simplicity) is a principle of purely universal, not partly subjective
validity. In any case the principle of the economy of thinking does not lead without further ado to using
no auxiliary concepts at all, any more than an economical household is one that uses no pots and no fuel
at all. Here is a point where Avenarius deviates, not to his advantage, from the philosophizing natural
scientists mentioned above; these saw clearly the connection between economy and symbolism, though they
laid different weight on these two sides of knowledge. --- With this it is also connected that Avenarius
takes up a somewhat dogmatic standpoint in the relation between physiology and psychology. He does not
see that epistemologically we always know the ``independent vital series'' through the ``dependent vital
series'': we should have no idea of the central nervous system and its states if we did not have at all the
capacity to form
% --- p. 106 ---
\opage{106}ideas; and even our present ideas of the central nervous system are due to posings and solutions
of problems that do not yet satisfy the demands of ``pure experience.''

In another respect too Avenarius does not lay sufficient weight on the fact that we are placed in a
development, and that the idea of an ideal conclusion can have significance only as a measure for the
course and direction of development. He holds predominantly to the case P > E, thus to such cases where
the demand for work is posed from outside. But a condition for the progress of life --- of the life of
knowledge as well as of the life of will --- is that the demand for work can also be posed from within, in
that there is a surplus of force that seeks application (E > P). In all of nature life begins with such a
surplus; it is that which makes growth possible and leads to new shoots. Precisely when life stirs
vigorously, it yearns beyond the given, seeks out or evokes new experiences, sets new beginnings\textsuperscript{34}). It happens when one ventures, and it happens especially when one values. If
Avenarius, to whose plans it belonged to write an ``Ethics of the Open Air,'' had come to treat the
problem of valuation, he would certainly have laid weight on the fact that in all valuation one goes
beyond the merely given. Here the subjective elements can least of all be eliminated, even if they must
constantly be rectified through objective criteria. In Guyau and Nietzsche we shall come to occupy
ourselves with this question. ---

These critical remarks were meant only to contribute to eliminating some of the subjective accessory
ideas (\textit{α}) which have hindered the energetic inquirer from seeing matters clearly. His excellent
idea of a natural history of problems\textsuperscript{35}) would without them --- and without the repellent
form in which it is presented in his books --- quite certainly have made a greater impression on the age.
It is an idea that will be taken up again, and the work that is put into it will not be wasted.

% --- p. 107 ---
\opage{107}\grouphead{THIRD GROUP.}{THE PHILOSOPHY OF VALUATION.}

The three relations between the inner and outer conditions of knowledge
that Avenarius distinguished can be used to
characterize the difference between the three groups of philosophical
endeavors that appear in the most recent period. The systematic
group is essentially characterized by presupposing an equilibrium
between need and satisfaction or between capacity and
work (thus E = P). The epistemological-biological
group lays the main stress on the fact that existence can
make other and greater demands on our capacity than this can fully
satisfy (thus E < P). The third group, finally, emphasizes that there can be a surplus of need and capacity
(thus E > P), which leads to a valuation of existence
and to a striving beyond the given. As representatives of
this last outlook I shall here consider Guyau and Nietzsche.
It would be justified to mention other modern philosophers of valuation
here as well, notably William James and Rudolph Eucken;
but of these I have recently had occasion to speak in another work \textit{(Om nogle religionsfilosofiske
Arbejder fra den nyeste Tid.} University program
1903).

Guyau and Nietzsche both stand on the ground of the doctrine of evolution
and hope for the emergence of higher forms of life. They base this
hope on the power and health of life, --- on the conviction
that there is an overflowing fullness of energy which finds neither
% --- p. 108 ---
\opage{108}expression nor outlet in our present experience and our present
conditions of life. They pronounce a great Yes against the No of pessimism,
--- the one in a mood of deep feeling and gentle resignation
before the shifting and transience of values, the other in
defiance and contempt of the past and in an unbridled, in the end convulsive,
hope for the future.

A main problem in both is the relation between instinct
and reflection, --- between the undivided energy of life at its earlier
stages and its divided workings under the advance
of culture and reflection. Since this problem was first taken up toward the end
of the eighteenth century, by Rousseau, Lessing and Kant,
it has not been posed as energetically as by the two authors
with whom we are here concerned. Both polemicize against
one-sided intellectualism and rely on the life of feeling and will
as that which can never be wholly shaped into clear, rational
ideas. Whereas in the other group it is the fullness
given from without that may threaten to overwhelm thought,
here it is the fullness rising from within, from the inner world of the mind,
that drives thought to its limit.

Connected with this is the fact that their presentations stand
on the borderline between philosophy and poetry. Mood and passion operate
at every point, not always to the benefit of the clarity and consistency
of the inquiry, but to the benefit of the literary,
or even agitatorial, effect their writings exert. Their
thoughts acquire, for the most part, more significance as symptoms
of what stirs in the minds of those who are most strongly
gripped by life in our day than as active contributions to the treatment
of problems. This holds in quite special measure for
Nietzsche.

There is one further point of likeness: both are ill, and
their thoughts and works have for the most part come into being in
constant struggle with illness. This decides nothing with regard to
the value of the thoughts. There could be thoughts that cast light
on life and that could come into being only under such
conditions. The border territory between health and illness might
well be quite particularly fruitful. We must therefore first of all
discuss the significance and validity of their thoughts, and only
% --- p. 109 ---
\opage{109}when this test has been carried out can we use what we
know of their individual circumstances to explain to ourselves their
bent of mind.

\chapterhead{I.}{JEAN MARIE GUYAU.}

Guyau is an example of early maturity. At only twenty
years of age (he was born in 1854) he won the prize of the French Academy of Sciences
for a work on the history of the morality of utility
from Epicurus to our own day. His first studies concerned chiefly
Plato and Kant, and he was also strongly influenced by his stepfather
Fouillée. A brilliant gift of presentation and a penetrating
critical energy attracted attention to his work,
and won it recognition even from those who were sharply
criticized in it, such as Spencer. In the course of his further studies
he drew somewhat closer to the outlooks he had
attacked, and in his positive presentations he often
has difficulty slipping past objections he himself had
raised against the English school. It is aesthetic, ethical
and religio-philosophical problems that have chiefly occupied him.
Gnawing doubt about the significance of our values in existence
pursued him from his youth and found expression not only in
philosophical writings but also in poems \textit{(Vers d'un philosophe).}
His original, Platonizing faith in the rationality of the world
paled under the combined influence of reflection and the lung disease.
It does seem that later he again held fast to the possibility
of an ideal significance of existence, but the clear, philosophical
view of the difficulty of the problem he did not lose. One
has found a French trait in the way in which, with him,
inwardness of mood and clarity of thought so often go together.
He knows how many of his visions and hopes
are illusions; but he holds fast to them when they are merely
fruitful illusions, capable of setting thought and will to
energetic work. (The poem Illusion féconde). He sees the end
clearly before his eyes, when his illness has advanced;
but he wishes to put all his love of human beings into his
% --- p. 110 ---
\opage{110}last song (La cicade). Hardly any other author
has shown as he has what inwardness and loftiness life can
have even if all dogmas are regarded as illusions. He
lived, moreover --- apart from the illness
--- under fortunate circumstances.
Researching and writing poetry, loving and loved, he approached
death in full consciousness of the outcome of his illness. He
died in Menton in 1888. On his grave were set the following
lines from one of his writings: ``Our highest endeavors seem
precisely to be the most useless; but they are to be likened to waves
that have been able to reach us and therefore will also be able to
reach farther. I am certain that the best that I possess will
outlive me. No, perhaps not one of my dreams will be
wasted; others will take them up, will dream them again
after me, until some day they are fulfilled. It is by the help of the
ever-dying waves that the sea is able to shape its shore
and the vast bed in which it moves.''

\sectionhead{A. CRITIQUE OF ENGLISH ETHICS. \textit{(La morale\\ anglaise contemporaine.} 1879).}

Already in Guyau's critical work we find hints
of his definitive conception. What he objects to in the ethics of
the English school is chiefly three things.

The English teach that morality is essentially due to prudent calculation
(Bentham), or association of ideas (Hartley, the two
Mills), or natural selection (Darwin), or adaptation
to the conditions of life (Spencer). They base the validity and
persistence of morality on definite inner and outer conditions of development. The
moral feeling is a natural phenomenon, which man regards only
as an objective power because he has not noticed how
it came to be. The involuntary and unconscious growth of
the moral feeling is the main subject of English ethics. But
here it now becomes entangled, Guyau holds, in a self-contradiction
and works against itself. For these very English theories must
open man's eye to how matters stand
with the nature of the mysterious and moral feeling.
It
% --- p. 111 ---
\opage{111}comes to stand as a deception that one seeks to free oneself from.
Reflection will dissolve what has formed itself quite involuntarily.

But --- has a real development taken place here?
The English presuppose as the ultimate foundation the egoistic
drive to self-preservation, and does it not remain the same,
however many calculations, associations and adaptations
may be added? The capacity for disinterested devotion cannot
be explained through such factors. It does not follow from this
that disinterested sympathy is an original tendency. But
what the English conceive as original egoism is, according to
Guyau, something else and more, lying deeper than both
egoism and sympathy. It is a need for activity, for unfolding
and expansion, which is present everywhere that life is
healthy and vigorous. This need does not have to go out against
others, but can on the contrary lead to association with others and to
work for them. It is thus not first through many intermediate links,
but already in its first germ, that the human need
contains a possibility of devotion to a more comprehensive
life than the purely individual. --- Guyau here develops a thought
that had already emerged in Vauvenargues and Rousseau,
without however seeming to have it from them. ---

When weight is thus laid on the inner need for activity and development,
a need partly independent of external
occasions, our striving does not --- as is the consequence in the morality of utility
--- take on the character of a mere means.
It can then hold that we live in order to will and act, not merely
the reverse. Il faut vivre pour vouloir et agir! The present is to be sacrificed
neither for the future nor for the past. In the very act of will
there lies an end: il y a dans la volonté quelque
chose de définitif! ---

\sectionhead{B. ETHICS. \textit{(Esquiss d'une morale sans obligation ni\\ sanction.} 1885).}

As the foundation of his ethics Guyau lays the concept of life in the
widest sense of the word, whereby the opposition between the conscious
and the unconscious and between the egoistic and the
% --- p. 112 ---
\opage{112}sympathetic falls away. Life manifests itself as a need for growth,
persistence, reproduction, expansion. It is biology by whose
help he wishes to remedy the difficulties in which the English
entangled themselves through their psychological grounding:
yet he concedes that Darwin and Spencer have in fact already
demonstrated a biological foundation. It is a main point for
him that ends are in the last resort not produced by consciousness
when it sets itself ends: what happens is that one
becomes conscious of, or makes an object for oneself of, what has already
asserted itself unconsciously. Ends thus do
not come wholly from without.

Ethics then becomes the doctrine of the means by which the end
set by nature itself --- the growth and unfolding of life --- can
be attained. And the supreme ethical commandment becomes: Unfold your life in all
directions! Be as rich a human being as possible, both as regards
the inwardness and as regards the scope of your endeavors! And
therefore --- live in society with others and train your capacity
to live in society as much as possible! --- Duty is only
an overflowing of life, which demands to be used and to give
itself up. It is not due to compulsion and external necessity, but
is the expression of a force. In sympathy or altruism
this overflowing force emerges especially. Egoism, on the other hand, is
a sign of a narrowing, an isolation, and it rests in the end
on an illusion. The highest virtue, and moreover the one most
firmly grounded in nature, is magnanimity (générosité). ---

Guyau has sought to avoid falling prey to the critique he himself
has directed against the English by going back to the unconscious
and involuntary. But he overlooks the different possibilities
--- and with them the different problems --- that lie in the
transition from the unconscious to the conscious and from the involuntary
to the voluntary. It does not follow that the conscious
exhausts or wholly corresponds to the unconscious, and a similar
relation obtains between the voluntary and the involuntary.
Means presuppose an end, and ethics cannot seek means
when it does not know what end they are to serve. What we
make an end depends in turn on what we ascribe value to.
Mediate values (the means) presuppose immediate
% --- p. 113 ---
\opage{113}values. With these one must therefore begin, and they cannot
be found without the help of psychology. An empirical doctrine of
values must form the foundation, and if there are several fundamental
values that cannot be reduced to a single one, then
a difficulty lies here for all ethics right at the very beginning.
It is only so long as one keeps to vague generality
that an appeal to biology is sufficient. As soon as
definite questions are posed, the concept of value emerges. Guyau
himself has furnished an example of this. Against Bourget's assertion
that it is arbitrary and dogmatic to distinguish between natural
and unnatural feelings and impulses, Guyau maintains (in his work
\textit{L'art au point de vue sociologique.} p. 375) that there may very well
be a difference in value even if the different states
arise with equal necessity. There is, he says, a
criterion of natural value (valeur naturelle), which must be sought
in the power and scope of life and in the consciousness and the feeling of pleasure
that is the inner revelation of it. --- It is clear
that however powerful and comprehensive ``life'' might be, --- if
no ``consciousness and feeling of pleasure'' were attached to it,
no concept of value could be established, and no ethical discussion
could take place\textsuperscript{36}).

But the need that is satisfied by life's increase in power
and scope is, after all, involuntary and instinctive; will it not
be weakened by being drawn forth as an object of attention
and reflection? --- That was the main difficulty Guyau had
found in the English, and from which he must therefore try to free
his own theory. He declares that reason paralyzes an
instinct only when it goes in a different direction from the latter, or
when it is an advantage that it take its place. But neither
holds here. The involuntary devotion and unfolding
has its great practical significance. The desire to venture
springs from life's surplus of force, which begets hope and
faith, and thereby makes possible something that would not otherwise be possible.
In bold movement of thought, in enthusiastic self-sacrifice, this
power and health of life express themselves. The only sanction that the ethics of the
future will recognize will be the satisfaction of venturing (le
plaisir de risque). --- I cannot see that the English could not
% --- p. 114 ---
\opage{114}just as well appeal to these grounds; in fact they do
so, and to a higher degree than Guyau, since their ethics is more empirical
than his. If prudence, association and adaptation
could not lead to anything of value, they would not be the steady
paths along which development advances. And besides: if
reflection can dissolve instinct, can it not also inhibit
expansion?

\sectionhead{C. AESTHETICS. \textit{(Les problèmes de l'estétique\\ contemporaine.} 1884).}

Guyau's aesthetics offers an interesting example of
how a problem can present itself differently according to the different
conditions of different times. For earlier inquirers, such as Kant and Schiller,
and still much later for Spencer, the task was
to make room for art in the midst of the struggle of life and in spite of it.
They took refuge in the same thought that plays so
great a role in Guyau, in that they pointed out that when there comes
a surplus of energy, it is naturally used for functions similar
to those employed in the struggle of life, only that it
is now image and play that matter, not the reality of life itself.
Guyau objects to this whole tendency in aesthetics:
in art we seek precisely more and fuller life than the ordinary
life can bring; art is an expansion of life, and
it is only a pity that we must resort to image and play in order to attain
it. The unreal is by no means a condition of the beauty of art,
but precisely a limitation it must put up with.
In artistic contemplation arises the wish to become what
we behold, and it is the poet's sorrow that he cannot make himself
one with the whole fullness of life (the poem Le mal du poète). Therefore,
in Guyau's conception, beauty of nature stands above beauty of art.


From this it follows that Guyau asserts that all
the senses cooperate in the beautiful. A drink of milk at a
mountain hut can have the effect of a whole pastoral symphony. All sides of
our being cooperate --- through sensation, through memory,
through imagination. Beautiful Guyau calls an influence that at
once sets feeling, thought and will in motion and awakens
% --- p. 115 ---
\opage{115}pleasure in this combined working. More and more will all enjoyment
and feeling of pleasure take on the stamp of beauty, so that the distinction that
at an intermediate stage of development must be made between the agreeable
and the beautiful will fall away; this distinction is conditioned
only by the animal still being so strong in man.
Then the aesthetic and the ethical will also coincide in the assertion
of life's natural values. In the end it will turn out that
no artist can really present anything but his own life
--- and the value of the work will rest on the value of this life!
--- Readers of Guyau's writings will, on this last remark, think
that it fits Guyau's own authorship, inspired
as it is by the great Eros that was the soul of his life.

\sectionhead{D. PHILOSOPHY OF RELIGION. \textit{(L'irreligion de l'avenir.} 1887).}

Guyau's most remarkable and most significant book is the one that deals
with ``The Irreligion of the Future.'' This book is at once
written with the most searching critique and with the most heartfelt
feeling, with the fullest certainty that the time of religion is
over, and with the most cheerful hope that a high and ideal
view of life will never lack conditions for developing. In its title it
is directed against the attempts to construct a religion of the future,
but at the same time against ``antireligion,'' in that it will assert
the inner connection between the religious standpoint and that
which is to replace it.

The religious problem is for Guyau closely connected with
the ethical and the aesthetic. All three domains meet for him
in the concept of life. In concentrated personal and social life
all three elements appear, and only thereby does each
come into its own. Religion begins, it is true, as a kind of
mythological physics; but the intellectual elements are not
the predominant ones. What lies at the base is man's experiences
in the community of life in which he lives. This community of life
he involuntarily extends to the whole of existence, so that
gods and men form one great society. Religion is
not merely anthropomorphism, but in particular sociomorphism. The
difference between religions rests chiefly on what kind of society one imagines
the universe to be, a family or a realm? --- and in
% --- p. 116 ---
\opage{116}the latter case: ruled by a father or a tyrant? --- The social
character of religion receives its most pronounced expression in cult,
through which man enters into society with the gods.

Gradually both the religious cult and the religious representations are
sublimated and idealized. Inner cult tends to
replace outer cult, and mysticism replaces mythology. The culmination
is reached when God stands as the personified moral
ideal. --- As essential features in all religion Guyau emphasizes:
mythological explanation of nature, which in the higher religions is continued
in the acceptance of miracles, --- a circle of dogmas as absolute
truths, --- and cult as a circle of actions with
supernatural effects. Where these three things are not found,
Guyau will not speak of religion.

The dissolution of a religion does not occur directly and from without, but
occurs by its inner conditions of life falling away. And this occurs
again slowly, according as industry, science, individualism
and independent ethics develop. Of these factors Guyau lays
most weight on independent individual conviction: it
is this factor that has especially led to the victory of one religion over
another; it will in the end lead to the gradual dissolution of
every dogmatic system, and thereby to the final cessation of
religion. This process of dissolution may last long; but the ideas
work uninterruptedly and imperceptibly; --- they advance even if they
--- like tired soldiers --- sometimes sleep on the march.

But the final absence of religion (l'absence finale de
religion), irreligion, is not the same as antireligion, the complete
opposite of religion. It is on the contrary to become
a higher degree of religion and of civilization, in that the best in
the religious life, after the falling away of dogma, will grow in strength and
in scope. The eternal in religion is the very striving that
has produced it, --- the need to go beyond the bare
fact and find a greater connection. The human
spirit is like the swallow: its long wings are not suited to a
flight along the ground, but to a high and bold flight in the free air;
the only thing that matters is that it can lift itself up from the earth, and that
often proves hard, but the imperishable longing toward the ideal
will again and again provide it with air under its wings. This longing
% --- p. 117 ---
\opage{117}will stir more strongly than before when it is separated from religion.
Only one will feel less need to find a definite
language in which the eternal riddles could be expressed; they will speak
for themselves. The ultimate hypotheses will always retain a pronounced
individual character. No religion can express everything
that the single personality in its whole peculiarity contains
or needs. Everyone must produce his own God
or his own Bible. And different convictions may very
well exist side by side among human beings, just as different
plants can thrive in the same soil. --- In particular will the
ethical elements in religion be preserved and set free. Already
now, religion in the noblest human beings tends to
become one with love. --- The mood toward life and existence
will not be simple and constant; it will shift and waver.
Development has lasted an eternity --- and has then led
only to a world like this, with all its disharmonies and with
all the uncertainty in its course! The more man, by a kind of ideal
sociology, lives himself into the whole of existence, as he is able to
know it, the more will melancholy grow; thought is pain,
not merely light. God himself would, if he existed, have to feel the
greatest pain, since as an infinite being he would have to feel most bitterly
his impotence to help. But the mood of the whole will have
the stamp of sublimity, so that sorrow in the striver becomes only
an element in the feeling of life. The surplus of force will beget hope,
faith and desire to venture. Even in the face of death, ``the greatest novelty
in the intellectual life next to birth,'' the striver stands,
despite all pain, with unshaken intellectual firmness: Notre
dernière douleur reste aussi notre dernière curiosité! ---

This last remark is characteristic of Guyau's tireless
work of thought, which went hand in hand with his deep seriousness,
melancholy mood and great resignation. It has been said of him
that he sails out upon the infinite sea without ever
casting anchor. But can one ever cast anchor once
and for all? --- As regards the religious problem, he has dug
as deep as few, and he has clearly seen the necessity of
asking whether the concentration of life that the religions conditioned in
their classical periods can vanish without an
% --- p. 118 ---
\opage{118}equivalent coming in its place. He has here sought to assert continuity in the
domain of spiritual values, although he holds that the concept of religion
applies to these values only when they appear in a
certain definite form.\textsuperscript{37}) The continuity between religion and irreligion
he emphasizes in any case so strongly that the best in
religion (and also what has been the real tendency
in its coming into being) is to be preserved in irreligion.

\chapterhead{II.}{FRIEDRICH NIETZSCHE.}

\sectionhead{A. \textit{Characterization and Biography.}}

Whether it is essentially a poet or a thinker that we have
before us here may be a subject of dispute, and has been a subject
of dispute. It can in any case be claimed that it would
have been fortunate for the lasting worth of Nietzsche's works,
and that there would have been greater harmony between the poet
and the thinker in him, had he used the form of dialogue
or of drama, or perhaps a similar form to Kierkegaard's
pseudonyms, to present his ideas. For it is different,
conflicting thoughts that stir in him, and stir with
the same passion; it is as if it were different personalities
that appear, not only in his different writings, but
often in one and the same writing. But his very main thought --- in any
case what I shall try to show to be his main thought
--- clothed itself for him in poetic form, in that only in great, pictorial
strokes could he give vent to what lay deepest in his
heart. The most significant man among his youthful friends,
Erwin Rohde, regarded him essentially as a poet, and
therefore rejoiced when ``Zarathustra'' began to appear.
``I have,'' he says, ``long had the feeling that Nietzsche suffers
from a fullness of poetry that will not unfold into real poetic creation
and now causes fever and distress in his inner being.'' In a very
clear manner Nietzsche himself has spoken of the relation between
the poetic and the philosophical element in him, and of the significance
of both for him, in a preface (of 1886) to a new edition of
% --- p. 119 ---
\opage{119}his early work ``Die Geburt der Tragödie,'' where he designates
it as his constant task to see science from the standpoint of art and art
from the standpoint of life. This
statement shows at the same time that I am right in designating his
tendency, together with Guyau's, as philosophy of valuation. It is,
as with Guyau, the power and fullness of life that is his innermost
basic thought; and Nietzsche did in fact read with great interest
several of Guyau's writings. What he strives for is a new, positive
valuation of life on the basis of a consideration from the standpoint of the history of culture.
In relation to this, that which in wider circles has most caught
the eye, the doctrine of master and slave morality, the social dualism
and the contempt for the multitude (the ``herd''), is derived and subordinate
--- and, as I shall try to show, wrongly derived. It is
on this point, in the relation between his real basic thought
and his special outlooks and antipathies, that the greatest contradictions
are found in him. But those of his writings in which these utterances,
derived (and wrongly derived) in relation to the basic thought,
are found have been the most read. His literary influence rests
on this side of his authorship. Without wishing here to
renew the dispute that I had some years ago on the occasion
of Nietzsche with Georg Brandes, I cannot refrain from
remarking that it was unfortunate that Nietzsche was introduced into our
literature from this side. The weight was laid on what is most striking
in him, but not on what is fundamental. Now that one
--- especially after the publication of so many of Nietzsche's posthumous
writings and drafts --- can survey his production tolerably well,
there appears to me the humorous fact that nearly the only thing
on which my opponent of that time and I were agreed was something
in which we were both wrong. We were agreed that Nietzsche was
an opponent of the morality of welfare, but it now turns out that Nietzsche
was in reality a moralist of welfare on a grand scale. This must be emphasized
here, where we are not concerned with what is perhaps most epoch-making
in literary terms, but with what is philosophically
decisive: the basic thoughts, their grounding and their
struggle.

The significance that the literary and poetic element has in
Nietzsche entails greater difficulties for presentation and
% --- p. 120 ---
\opage{120}characterization than in other philosophers. In order to give a tolerably
correct picture of him, my presentation must be fuller
than it would need to be if it were merely a matter of a purely intellectual reckoning.
Yet it will be of great psychological interest to study
the relation between thinking, art and life in a genius whose
striving was to let all thinking and poetry serve the deepening
and elevation of life. A psychological yield here will also be of
philosophical significance. ---

Friedrich Nietzsche was born in Prussian Saxony on the 15th of
October 1844, of an old family of clergymen. He himself would gladly
have been of nobler origin, and believed that his great-grandfather
had been a Polish nobleman. Some of his critics have
on the contrary precisely thought to find the clergyman's son in his line of thought.
Very early he was drawn under other influences than the
ecclesiastical, in that at school and at the university he zealously cultivated
classical studies. First-hand acquaintance with Greek spiritual life
and the enthusiasm over it became decisive for the whole striving and poetry
of his life. But alongside this the study of Schopenhauer,
which he pursued simultaneously with his philological studies,
was of deeply penetrating significance. Of other philosophers
Dühring and Albert Lange have had influence on him; but Schopenhauer
became his real ``educator.'' On the occasion of his
study of Schopenhauer he wrote (1869) to his friend Paul Deussen:
``A philosophy that we adopt only from a pure urge for knowledge
never becomes wholly our own, because it never \textit{was} our own.
Every individual's true philosophy is \textit{ἀνάμνησις} [recollection].''
The urge for knowledge expressed itself in Nietzsche always in inseparable
fusion with the urge for art and the urge for life. That
is his strength, and it is his weakness --- as a thinker, especially
since his thoughts, on account of his impulsive and disharmonious
nature, were so often dictated by the mood of the moment or by a
reaction against tendencies in the surrounding world.

Nietzsche's second great educator was Richard Wagner.
Already during his years of study in Leipzig he had made personal
acquaintance with the great Master, and when in 1869 he became
professor of philology in Basel, he associated intimately with Wagner
and his wife, who in those years lived near Lucerne.

% --- p. 121 ---
\opage{121}A third educator he also had: illness, which for
many years consumed his strength and in the end became incurable
insanity. He is said to have contracted it while serving as
a nurse in the Franco-German War. It was contracted
by infection, and \textsc{Møbius} \textit{(Ueber das Pathologische bei Nietzsche)}
therefore maintains that the brain disease from which he came to suffer
cannot be brought into connection with his original personality,
although it obtained a determining influence on the latter's later
development, and although he must be regarded as having been hereditarily
predisposed. According to Møbius Nietzsche's condition was aggravated by
wrong medical treatment and by his own abuse of chloral. When
Møbius also thinks he can point to a definite place in one
of Nietzsche's writings where the illness is supposed to have begun to
exert its influence, I believe that this is wholly unfounded; but
of this later. --- After a long struggle the illness made it
impossible for him to continue his university activity; philology besides had
become too narrow for him as the great problems
gathered momentum in him and set imagination and thought in motion.
In 1879 he took his leave, and lived in the following
years mostly in the Engadine in summer and on the Riviera in winter.
In the wonderful mountain nature of the Engadine, with its light
air and its remarkable illuminations, high above the ordinary
human life in the valleys, --- up here ``on the heights'' he
conceived, by his own account, some of his most significant
ideas.

It is peculiar to Nietzsche that all his three educators
worked most decisively upon him through the reaction
they released, or the effect of contrast they conditioned. In illness
his will was hardened, and bravery in suffering became his
chief virtue. He learned the pain of life at first hand, but also
life's overflowing power to assert itself in spite of everything.
The education of illness he never got free from; the two other
teachers he could part from, in any case apparently.
With the living educator, Richard Wagner, the
break was painful. With Schopenhauer it was pessimism
that at a certain point in his development repelled him; with
Wagner he was repelled both by the pessimism and by the
% --- p. 122 ---
\opage{122}turn to medieval Christianity that he thought he found in the
great tone-artist,\textsuperscript{38}) whose human weaknesses he gradually
acquired a sharp eye for. Siegfried he had admired
and loved; already at
Wotan --- the god as pessimist!
--- he had become doubtful; Parsifal he hated. From the first he had
admired Wagner's art as ``Dionysian''; later
it stood for him as an expression of decadence.

In studying Nietzsche's biography\textsuperscript{39}) one naturally looks
for experiences that could have brought him into contact with
the people, the great mass, the ``herd,'' as he later called it.
It was natural to suppose that he had particularly often
had occasion to see human beings appear en masse. Only two
such occasions are found, and both times his contempt was aroused.
The one time it was beer-drinking German students;
the other time it was the theater public in Bayreuth, which
did not answer to his expectations of the ideal audience
at the rebirth of great art. --- On the other hand Nietzsche
never had occasion to live among the people and to learn to know their striving,
their capacity for development and organization; he stood wholly
alien to the greatest social events of the
nineteenth century. ---

In the course of the illness Møbius thinks he can distinguish several
periods. A rise of the morbidity takes place from 1881
to 1884; then a decline in the following years, and then a new
rise, until at the beginning of 1889 a catastrophe came.
He sent letters round to several acquaintances, now signed
``Dionysos,'' now ``the Crucified.'' One day he fell in a faint
to the ground in the street in Turin. A friend brought him to Basel,
whence he was later taken to his home region. The last years he spent
in a state of complete torpor in Weimar, carefully
tended by his nearest. On the 25th of August 1900 he died.
--- As no autopsy took place, the nature and seat of the illness
have not been able to be ascertained with full certainty. Møbius
assumes that it was progressive paralysis.

Nietzsche was a man of opposites. The opposition
that appears especially in him, stamping his whole conduct and
fateful for his life and his significance, is the
% --- p. 123 ---
\opage{123}opposition between high enthusiasm and deep repugnance. Under
healthy conditions these two opposites are two different sides
of the same basic mood; in Nietzsche they were hostile powers
fighting for dominion in his mind and in his will.
In this struggle he perished. Admiration for the great
and lofty, longing for it, the wish to work for it was no doubt
the innermost thing in his soul; but antipathy and contempt for the
low and weak overwhelmed, hypnotized him, so that that for which he
wished to fight came to be overshadowed by that against which he wished to
fight. The ``reactive'' moods --- which he yet himself
despised in others --- gained the upper hand in him again and again.
He himself felt what danger here threatened his spiritual
life. In the year 1874 he writes to Miss v. Meysenbug: ``Wie
wird mir zu Muthe sein, wenn ich erst alles Negative und Empörte,
was in mir steckt, aus mir herausgestellt habe, und doch
darf ich hoffen, in fünf Jahren ungefähr diesem herrlichen Ziele
nahe zu sein.'' This hope was not fulfilled, even after
he had begun to work at the positive presentation of
his ideas. Even in the last year before the catastrophe he broke off
his connected work in order to give vent to his antipathies,
although he had opened his valves a sufficient number of
times before.

The predominant place that the antipathies came to play in
his production stands in opposition not only to his enthusiasm,
but to his loving and heartfelt disposition as it
expressed itself in his closest relationships. This too he felt himself.
He writes to Miss v. Meysenbug (1878): ``This autumn
I resolved to begin each day by asking myself:
Is there no one for whom you could do something good today?
. . . . With my writings I cause so many people harm
that I naturally must wish to make it good again in one
way or another.'' And some years later (1883) he wrote to
Erwin Rohde that he well knew that the picture his books now
[in the most polemical part of his authorship] gave
of him was not the picture his friend bore of him in his heart.
He had, he continues, really an ``other nature'' [than that
% --- p. 124 ---
\opage{124}which at this time expressed itself in his books]; by the ``first''
alone he would long since have perished!

Not only against enthusiasm and love did the
strong predominance of the antipathies strive, but also against the optimism
for which Nietzsche worked so passionately in his later
years. He demanded in increasing degree and in excessive terms an
affirmation of life, a recognition of its value. But
how could he demand this when life was so filled
with loathsomeness and wretchedness, as he again and again declared,
and could not finish declaring?

The conflicting moods could not be reconciled in him,
either in the form of great humor or of deep melancholy. Of
Shakespeare it has been said that his humor is an expression
of his fidelity to life. In this way
Nietzsche was not able to show his fidelity; perhaps it was connected
with his being neither poet nor thinker in a decided way,
but with two natures here too being at strife in him. His contemporary
and in many respects kindred thinker Guyau found in melancholy
and in the eternal striving of thought the harmonious balance.
This form too was denied to Nietzsche. A timbre that
gathered in itself the conflicting tones he did not attain. Oscillations
between Dionysian intoxication and consuming disharmony led him ever
nearer to dissolution. Here it will be justified to find
an effect of the illness. He fought honestly against
it. For he was a brave man, a man of will. He felt
deeply ``die Vehemenz der inneren Schwingungen''; but when his
critics called him ``eccentric,'' he said with pride:
``These gentlemen, who have no concept of my center, of the
great passion in whose service I live, will hardly be able
to have an eye for where I was really outside my center.''
(Letter to Deussen 1888).

His bravery he showed in face of work, suffering, misunderstanding
and loneliness. ``No pain,'' he says in a letter, ``has
been able and shall be able to seduce me to a false testimony
about life, as I know it.'' ``My existence,'' he writes
at about the same time, ``is a terrible burden, which I would long ago
have thrown off had I not made, precisely in this
% --- p. 125 ---
\opage{125}condition of suffering and almost absolute renunciation, the most instructive
trials and experiments in the spiritual and moral
domain; this joy that longs for knowledge leads me up
to heights where I triumph over all torment and all hopelessness.
On the whole I am nevertheless happier than ever before.'' ``Take
note,'' he wrote later, ``that the years in which I ceased
to be a pessimist were precisely the years in which my vital force was
lowest.''

Only one thing did he fear: pity. This
he had, in his pessimistic time and under Schopenhauer's ``education,''
sinned against. This fear, which stands in a strange
contrast to his bravery, finds partly its explanation in
his own hardness, but was also connected with his
own soft disposition, for which pity could be a temptation.
By pity he thinks above all of the passive,
sentimental form of sympathy, which weakens both the one
who has it and the object of it. It ended with pity
standing for him as the last great sin into which one could still
fall when other sins had become impossible. The
peculiar hypnotic effect that Nietzsche's antipathies had
on him expresses itself here too. In the end he sensed the odor of sickness
all around Europe, wherever he went and stood.

The form of Nietzsche's authorship was determined by
his temperament and by his morbidity. He had a need
for concentration, and there are in his production certain positive
main thoughts whose complete working-through he longed
for. But he was not able to bend the excessive mood of
the single moments under what the basic thought and the basic mood
demanded. The limitation that makes the master was
not in his power. Hence the discordance in his utterances; the
intoxication of one moment did not always accord with that of another. The aphorism
became his favorite form. Each single aphorism lets a
standpoint come to full, often striking expression, only in order to
give place to other aphorisms, which without more ado say the opposite.
The real work of thought does not get done. With justice
the aphoristic form has been called dangerous for Nietzsche.
He came to it partly by imitation of older French
% --- p. 126 ---
\opage{126}authors, especially Larochefoucauld, whom he highly admired,
partly on account of the morbidity, which made connected
work impossible for him. He produced on walks and
set down his thoughts in loose form; others then had to edit
them. ``Ich kritzele auf meinen Wegen hie und da etwas
auf ein Blatt; ich schreibe nichts am Schreibtisch; Freunde entziffern
meine Kritzeleien.'' (Letter to Eiser 1880). A great part
of Nietzsche's popularity is no doubt due to this form, which makes
it so easy to get at least a ``nodding acquaintance'' with
his books. Of one of his books written in this form he himself says:
``Ein Buch, wie dieses, ist nicht zum Durchlesen
und Vorlesen, sondern zum Aufschlagen, namentlich im
Spazirengehn und auf Reisen .'' (Morgenröthe. Aph. 454). It
is, as he himself has said, not easy to find his real center;
but his own form of presentation has made it more difficult
than it is in itself.

\sectionhead{B. \textit{Nietzsche's Writings.}}

It is a consequence of Nietzsche's disharmonious and wavering
temperament that his writings cannot be ordered purely chronologically,
but must be divided into groups, which must be characterized by
the different relation that appears between the basic thoughts
that on the whole are the ruling ones in him.

1. \textit{First Group.} The main problem is posed, and the definitive
solution is hinted at. Characteristic of this first group is
the cultural-historical foundation of the posing of the problem, which Nietzsche
draws from his philological studies. He relies on sociology,
whereas Guyau relies essentially on biology. \textit{Die Geburt der
Tragödie} (1872) is not only Nietzsche's first significant work,
but ---
as regards real treatment of problems --- perhaps the
most significant of all his writings. His main thought of a new
valuation of life emerges already here. In one of his last
writings (Götzendämmerung) he himself says: ``Die `Geburt der
Tragödie' war meine erste Umwertung aller Werte.'' The tragic-pathetic
conception of life --- symbolized in Dionysus and Apollo
--- is set in opposition to the intellectual optimism that
% --- p. 127 ---
\opage{127}is represented by Socrates.
The relation between science, art
and life appears here between these three figures --- and Dionysus
stands highest, although the Apollonian, shaping and limiting
element is recognized more than in the later books.

The standpoint thus won is applied in \textit{Unzeitgemässe
Betrachtungen} (1873---1876) to contemporary German spiritual culture. Of
the four considerations, the first contains a sharp critique of
Strauss, the second a polemic against the overvaluation of the
historical mode of consideration, the third and the fourth glorify
Schopenhauer and Wagner as the great educators. It
is the striving, tragically attuned personalities that are set in opposition
to the inquirer self-satisfied by his critical results
and to the purely objective occupation with the past. --- Later
Nietzsche held that he had read his own hopeful mood
into Schopenhauer and Wagner, who were in reality
decadents, and in his later production he never finishes
reacting against his own earlier admiration. He burns
what he has worshipped, and constantly lights the bonfire anew.

In this first group there emerge --- besides the basic relation
between science, art and life --- still other of Nietzsche's
main thoughts already: the radical aristocratism, which in
the great men sees the goal of history, and the social dualism
between masters and slaves connected with it.

2. \textit{Second Group.} Here the oppositions within
Nietzsche's line of thought emerge more strongly, and also his opposition
to other standpoints. The break with Schopenhauer and with
Wagner became here of particular significance. He speaks later of an
anti-romantic cure to which he subjected himself, because he suffered from
the most dangerous form of romanticism. The point was to defend life
against the conclusions one is accustomed to draw from pain, disappointment
and isolation. He surely means here the same as when
in a letter of 1883 he speaks of a long and heavy spiritual asceticism
to which he has subjected himself in the last six years. The cure or
asceticism he mentions consisted in realistic studies. He made
acquaintance with natural science, with older French philosophy and
with modern English philosophy. Thereby a new element came
into his line of thought, which was difficult to reconcile with
% --- p. 128 ---
\opage{128}the earlier ones; and there arose occasion for new polemic and new reaction.
It was at this time that he came upon the aphoristic form,
from which he was not again able to free himself. Its
dominion was bound naturally to become the stronger the more various
tendencies, sympathies and antipathies came forth in
him and each demanded expression.

It was not merely a passing cure, however, to which he here subjected
himself. The ethical problem worked its way forward for him more determinately,
as a special form of the general problem of culture
that the first writings had posed. A thorough doubt about the
moral presuppositions (``valuations'') hitherto holding
in life and in thinking arose in him. ``Ich begann, unser
Vertrauen zur Moral zu untergraben.'' (Preface of 1886 to
``Morgenröthe'').

To this belong the writings \textit{Menschliches, Allzumenschliches}
(1878---80), \textit{Morgenröthe} (1881). \textit{Fröhliche Wissenschaft} (1882).

At the end of ``Fröhliche Wissenschaft'' there is found for the first
time a thought that in Nietzsche's later years played a great
role with him, the thought that existence was rhythmic in its development,
so that what had once happened would come again
in the same manner in all particulars in a following period. In
connection with this the figure of Zarathustra emerges as the great prophet
of the affirmation of life, who will descend to men in order
to proclaim to them that life is so beautiful that its recurrence can
and must be willed. (Fröhl. Wiss. Aphr. 341---342). Shortly afterward
Nietzsche worked out the first part of his Zarathustra book. He writes
to a friend (June 1883): ``What I still have left of
life (only little, I believe) shall be used to express fully and completely
that for the sake of which I have at all endured life.
The time for keeping silent is now past; my Zarathustra, which in the course of
a few weeks will be sent to you, will betray to you how high
a flight my will takes. Do not let yourself be deceived by this book's
legendary form; behind all the plain and strange words lies my
deep seriousness and my whole philosophy.'' --- I do not understand why
Möbius believes that at that place in ``Fröhliche Wissenschaft''
Nietzsche's insanity had begun to attack his life of ideas,
after having previously attacked only his life of feeling.
% --- p. 129 ---
\opage{129}The figure of Zarathustra had occupied Nietzsche from his childhood
on, and it is moreover prepared by the figure of Dionysus in
``Geburt der Tragödie.'' The thought of the recurrence of all things had
Nietzsche from his Greek studies, since it occurs among
the Pythagoreans and Stoics; it comes forth now quite naturally
in him as a touchstone of how great the need for affirmation of life
is. The idea may seem glaring; but it is psychologically
intelligible in Nietsche precisely at this time. He
says himself in an autobiographical note: ``The basic composition
of `Zarathustra,' the thought of eternal recurrence, this highest
formula of affirmation that can at all be attained, belongs
to August 1881; it is thrown down on a sheet with the inscription:
6000 feet above men and above time! In those days I walked
through the forests by Lake Silvaplana and halted by a mighty
pyramidally towering block of rock near Surlei.
There I got this thought.''

The Zarathustra book itself I do not count in this group.
Work on it was again and again interrupted by the urge to
give vent in new aphorisms to polemical and antipathetic moods.
He did not get to have his ``great seriousness and whole
philosophy'' expressed. \textit{Jenseits von Gut und Böse} (1886) was to be
a commentary on the Zarathustra book (which was not yet finished
--- and never became so). But it is a commentary that presupposes
that one understands the book itself, the unfinished book that
is to be commented on. Nietzsche writes in a letter (Oct. 1886):
``Hast du dich in meinem `Jenseits' umgethan? Es ist eine
Art von Commentar zu meinem `Zarathustra'. Aber wie gut
müsste man mich verstehen, \textit{in wie fern} es zu ihm ein Commentar
ist!'' --- One cannot help asking how many
of those who as good Nietzscheans confidently began to live
``beyond good and evil'' have really satisfied this
demand that their Master makes. But neither is it an easy task
that is set them. And it is a new form of the tragedy in
Nietzsche's life and thinking that he came to write commentaries
on works that were not carried out, commentaries that nevertheless
themselves needed the works for --- a commentary! --- \textit{Die Genealogie
der Moral} (1887) was again to be a supplement to ``Jenseits.''
% --- p. 130 ---
\opage{130}It is an attempt at a historical grounding of the social dualism,
which in Nietzsche's aphoristic writings stands so strongly in the foreground
with him.

In these two writings the theory of slave revolts in morality
and of the necessity of an absolute new creation of values
comes clearly forth, although Nietzsche himself traces them back to ``Menschliches---Allzumenschliches''
and to his stay in southern Italy in the year
1876. These are the ideas in Nietzsche that have caught the eye
most in wider circles. Most strongly there appears here Nietzsche's
great need to feel himself at a distance from the multitude, --- what he
calls ``das Pathos der Distanz,'' --- and his caste morality is expressed
in the crassest form. But it must be held fast that the
two writings are only excerpts from a work planned already from the year 1881
on, and that they therefore cannot be regarded as Nietzsche's
last word. The mood that finds expression in them stands
precisely in the end in a decisive conflict with his deepest striving,
a conflict he has sometimes himself felt. --- Mrs. Elisabeth Førster-Nietzsche's
information in the preface to the fifteenth volume of Nietzsche's
works shows us the literary connection in his production
from 1881 to 1888, and in particular the relation of the various works
to the main work and its positive conclusion. From this it is also seen
that a series of violently polemical works from his last healthy
year (1888) are likewise excerpts from the main work. Thus: \textit{Der
Fall Wagner, Götzendämmerung, Nietzsche contra Wagner} and
\textit{Antichrist.}

3. \textit{Third Group.} His antipathies and his polemical moods
Nietzsche had difficulty giving full expression to. He
breaks off, right at the end, his main work in order with feverish haste to
send out his protests. And yet it ought to have lain near for him
who will undertake a revaluation (even of \textit{all} values!) to
make clear to himself and others on what foundation and by
what standard this revaluation is to take place. Such a
foundation Nietzsche did also seek, although he did not get it fully
presented. But he did not see clearly that the radical aristocratism
and the social dualism that he had proclaimed in his writings,
and to which he owed his greatest literary triumphs, stood
in irreconcilable conflict with his main thought, his basic value:
% --- p. 131 ---
\opage{131}the power, health and happiness of life. --- In two unfinished works
he worked at bringing this main thought forth as the ultimate measure of value.
The one is \textit{Also sprach Zarathustra. Ein Buch für
Alle und Keinen,} of which four parts appeared 1883---1886. With
these four parts must be compared the extraordinarily significant
draft of a concluding part, which is printed in the twelfth volume
of the works (p. 321). The other work, which in more philosophical
form was to give what ``Zarathustra'' gave in poetic form, is
\textit{Der Wille zur Macht. Versuch einer Umwerthung aller Werthe;}
it exists only in fragments, which are now published in the fifteenth
part of the works. It was to close like ``Zarathustra,'' in that
the last book was to have as its heading: ``Dionysos. Philosophie
der ewigen Wiederkunft.'' --- But Nietzsche had given
his antipathies vent to such a degree that he now lacked strength
to form his own positive ideal in image and in thought. And in particular
he was so unaccustomed to all connected thinking and
so suspicious of it, and of every definite
form whatever, that in the end Apollo no more than Socrates could express
what Dionysus meant. And in his ecstasy Dionysus
then tore himself to pieces.

\sectionhead{C. \textit{The Starting Point in the History of Culture.}}

Nietzsche wanted a new Renaissance, and like the Italian
Renaissance he begins by bringing the Greeks forward. His
philological and his philosophical interest were united here. When as a quite young man
he had become professor of philology in Basel,
he feared becoming a Philistine and ``Heerdenmensch'' if he
lost himself in specialties. He wished to permeate his science
with new blood by using it to illuminate the great problems of life.
And in this attempt Schopenhauer and Richard
Wagner were to be his teachers.

In \textit{Geburt der Tragödie} there is emphasized as the most important moment
in the history of Greek culture the union of Dionysian
enthusiasm with Apollonian harmony. The Bacchic frenzy,
which was an expression of overflowing fullness of life, was, through Delphi's
influence, brought into clear forms by receiving pictorial and
% --- p. 132 ---
\opage{132}artistic expression. Tragic art arose from the Dionysian
satyr choruses; the content of the choral songs was from the first interpreted by the chorus's
visions, in that it saw before it in great images its god's persecution, dismemberment and
resurrection. Thus tragedy was
born of the spirit of music, as an interpretation of the great destinies of life.
The prevailing conception of Greek antiquity, which is due especially to Winckelmann,
Goethe and Schiller, overlooked the Titanic background
of the antique harmony: Dionysus as the predecessor of Apollo!
Image, word and thought are secondary, derived in relation to
the great urge for life.

Greek tragedy perished through Socratism, which
inculcated that what mattered was knowing, understanding. Under
its influence Euripides became the first sober tragedian;
tragedy had gradually to yield its place to the sober,
bourgeois comedy, and Plato ironizes over the poetic enthusiasm
that does not itself know what it does. A new type enters
here into humanity: the theoretical man, the inquirer,
for whom perhaps in the end seeking stands above finding.
From that time on the struggle between the tragic and the theoretical
view of the world is perpetual, until deliverance comes --- through
Goethe's Faust, Kant's critique of reason, Schopenhauer's view of life
and Wagner's music. But highest among these liberating
powers stands the art of music, as Wagner has practiced it: in it
we lay our ear to the heart-chamber of existence; in it stirs
the Dionysian, the Bacchic never-resting striving, as
innermost force and ground in the whole of existence. It is of importance
that the urge for images is awakened; but there is a striving beyond
every image. Art is not merely imitation of nature,
but supplementation and overcoming of it. ---

Nietzsche uses the history of culture in broad strokes; this holds
not only of his utterances about Greek culture, but also of his
discussion of Christianity, the Reformation and the Revolution in
later writings. In Dionysus, Apollo and Socrates he has personified
three tendencies that constantly struggle and break against each other in human life.
In purely historical terms he is not right in seeing Socratism
as the chief cause of the decline of Greek culture. The chief cause
lay in the Greek life being drawn into the great world

% --- p. 133 ---
\opage{133}the whirl of events, and that in its (if one likes, Apollonian) form
in small societies it lacked the condition for asserting
itself under the new circumstances. And psychologically it is also one-sided
to let energy be formed by imagination and both be trained by thought.
In the individual and in the race there goes on continually
the opposite current as well, from thought through image to feeling
and will. And Socrates was himself a positive force, the model for
the great, wholly cast characters of the closing period of classical antiquity,
the time when it was no longer a matter of singing in chorus,
but when the voice of the individual (as in the fully dramatized
tragedy) was to bear witness to what he had got out of life.\textsuperscript{40})

\sectionhead{D. \textit{The Goal of History and the Social Dualism.}}

Against the exaggerated importance which, in Nietzsche's opinion,
historical education has acquired in recent times, he exclaims:
Remember to live life! (Memento vivere!) Do not let the past
weigh too heavily upon you, so that instinct,
personality, art, and thought suffer for it! Otherwise there will --- as already
Hesiod feared --- come a time when human beings are born
gray-haired!

The significance of history does not lie where one usually
finds it, in the endless process of development, or in the fate of the great
masses. In the individual, truly great human beings the whole value of
history is concentrated; in them it has reached its
goal, and the long historical process has only the significance
of giving the occasion and the force which are needed for the arising
of such human beings. The final end of humanity (Ziel) cannot
lie at the end of history, but can lie only in its highest
specimens. The great talk of historical development is of
evil; it wastes our force, and it excites the masses. The
great masses of human beings are only means or hindrances or
copies --- im Uebrigen hole sie der Teufel und die Statistik!
--- Greatness does not rest on mass effect; the noblest and best
has no effect at all on the masses.

Here Nietzsche expresses his radical aristocratism (aristocratic
radicalism is an incorrect designation). (Unzeitgemässe
% --- p. 134 ---
\opage{134}Betrachtungen. Zweites Stück).\textsuperscript{41}) It is radical aristocratism
in the sense that it is thoroughgoing aristocratism.
The noun radicalism Nietzsche would hardly acknowledge as his own, since
on the contrary he is very conservative in everything that concerns ``the herd.''
Here one must choose; one cannot be an adherent of
Nietzsche and at the same time be ``radical'' in social and political respects
--- nor really in ecclesiastical respects either. Of his
adherents the pseudonymous Peter Gast, who has edited, published,
and commented on several of his books, has clearly seen this.

The first motivation of aristocratism in Nietzsche
is thus connected with his conception of history, whose
significance he wants to reduce to as little as possible, which indeed
undeniably succeeds when the great human being stands there, raised
above past and future, not an effect of the former, not a cause of
or means for the latter. There is nonetheless an inconsistency, since the
great mass, besides being hindrance or copy, is yet supposed to be able to
be means. And what interest statistics ``im Uebrigen''
could have in occupying itself with the mass, when it is
not even seen as means or hindrance or copy, --- that
is not easy to see; nor what pleasure the Devil can
have in it, when it is not seen as hindrance. --- Schopenhauer's
influence, from which Nietzsche never really freed himself,
lies behind his philosophy of history. It is strange that
Nietzsche did not see how pessimistic this conception is,
although he worked throughout the whole last part of his active career to
free himself from pessimism. A sharp contradiction therefore appears,
as we shall see in what follows, in him on this
point.

Later he gives a somewhat different motivation of his
aristocratism from this philosophy-of-history and pessimistic one. In
\textit{Menschliches---Allzumenschliches} he starts from the force of life
and the originality of culture as a standard: If life is at all
to go forward and upward, its progress and ascent
can be represented only by a higher caste, which can perform free
work, and which presupposes a lower caste, which --- like a
kind of Cyclopes --- can perform the material work, which is always
more or less forced labor. Only that higher caste
% --- p. 135 ---
\opage{135}has the possibility of freedom of spirit. ``The higher caste,'' says Nietzsche
in one of his later writings \textit{(Antichrist,} § 57), ``I call
the fewest, and, being the perfect one, it also has the
privileges of the fewest, among which is that of representing happiness,
beauty, and goodness on earth.'' The two castes ought to be kept
well apart from each other. Their antagonism is necessary. The
lower caste is as far as possible to be equipped with machine virtues or
herd virtues. Religion, current morality, and the bourgeois
virtues suit the herd, not the chosen. ``Ich möchte,''
says Nietzsche \textit{(Der Wille zur Macht.} Aphor. 472), ``die liebenswürdigen
Tugenden nicht unterschätzen; aber die Grösse der
Seele verträgt sich nicht mit ihnen.''

The higher caste is only end, not also means. It is
according to Nietzsche a sign of corruption when an aristocracy no
longer regards itself as the real meaning and justification of society
and does not trust its right to make countless
other human beings into incomplete human beings, into slaves and tools.
Society exists only as substructure and scaffolding for
the select beings of whom the aristocracy consists --- like
a climbing plant on Java which clings to an oak in order to reach up into
the light with its crown and there display its happiness!
\textit{(Jenseits von Gut und Böse.} Aphor. 258). By its mere existence
this aristocracy expresses what gives life its value.
On the other hand its significance does not lie (as it does e.g. for the
aristocracy that Plato and Comte depict in their ideal states) in its
being to lead and help forward the other estates. Quite the contrary.
Right at the very end \textit{(Wille zur Macht.} Aphor. 12) Nietzsche declares
it a ``chief point of view'' that the task of the higher race
does not consist in the leadership of the lower, but that the lower is
the basis on which a higher race can live for its own
task. Yet even this higher race has a goal in the future: to
work for the coming of the overman. ``O my brothers,'' says
Zarathustra \textit{(Also sprach Zarathustra.} III), ``I consecrate you
to a new nobility: you shall be begetters and rearers,
sowers of the future! ... Your children's land shall
you love, and this love shall be your new nobility!'' ---
Thus there is after all still a goal in the future, and the goal of history
% --- p. 136 ---
\opage{136}is not reached in the aristocracy itself; a new formation of race
within the higher race is required. The problem proves more difficult
than it seemed at first.\textsuperscript{42})

The lower race, democracy, serves as material foundation,
as performer of the material work. Democratic institutions
have --- despite their tediousness --- the utility of keeping
tyranny away. But how it is possible to keep the
lower race within definite limits, we are not told.
Nietzsche himself says in his first writing that nothing more terrible
exists than a barbarian slave class that has learned to regard
its existence as an injustice and is preparing to take revenge.
But can one then prevent the Cyclopes from becoming knowing and from
comparing? It will surely always be easier to convince
the masters that they need slaves than the slaves that they
need masters. Nietzsche is here all too childishly romantic
--- or perhaps too barbaric. The maintenance of the military state
is declared by him to be the last means of maintaining the great
tradition of the high and strong human type \textit{(Wille zur Macht.}
Aphor. 327) --- a remark that casts a light upon Nietzsche's
inner connection with modern German conditions. His friend
Peter Gast accordingly also praises the military as (Such things must be said
in German!): die herrlichste, straffste, männlichste Einrichtung unserer
plebejisch- und mercantilistisch effeminirten Zeit! --- Popular enlightenment
is then of evil, since it makes the lower race disinclined
to put up with its fate --- as means or copy!

Yet the higher race is required to do strict work and
to exercise great self-control, in order to be able to keep its position.
The nobility must regulate its life in a definite way, if respect
for it is to be preserved. The higher classes must not themselves live as
rabble. They must live moderately and without ostentatious luxury.
Gilded coarseness and foolish pomposity break down respect
for culture. Likewise sudden enrichment without
work and bad distribution of property are unfortunate. (\textit{Menschliches---
Allzumenschliches}). Here a regard is thus required
which makes the higher caste or race unable without further ado
to live life out of itself. It must ``govern'' itself, and
perhaps also the other castes (namely when the military is called in).

% --- p. 137 ---
\opage{137}It is not easy to see why greatness should become less
by exercising a function. As if the dignity of force suffered
by force's being used! Where there is really a surplus of
force, there will be the capacity to be end and means at
once, and both in a grand style. And in the reach of the effect
we have after all the only measure of force.

On this point the fundamental difference between Nietzsche
and Guyau emerges. They agree in laying great weight on
the force of life and the corresponding fullness. It is the object of
their common faith. In his copies of Guyau's books Nietzsche has
furnished his statements in this direction with marginal notes
expressing recognition and assent. But when
Guyau in involuntary expansion finds a basis for devotion
and sympathy both toward human beings and toward ideas,
Nietzsche protests in his notes, and when this thought returns
in Guyau, Nietzsche writes in the margin: ``Fix Ide!''
He himself interprets the involuntary urge to action, conditioned by a surplus of force,
as a ``will to power,'' and more closely as
a need to let one's power go out over others. With Larochefoucauld
he wants to trace all feelings and drives back to egoism;
and egoism he conceives especially as a need to discharge
power (Macht auslassen).\textsuperscript{43}) Undoubtedly Guyau has here
seen rightly in regarding that involuntary urge as lying
deeper than the opposition between egoism and altruism.
Nietzsche is here more dogmatic than the French philosopher. ---

There is a great contradiction in the social dualism when the lower caste
is required to show admiration and reverence for a
culture which it is itself to stand wholly outside. How can the herd
have respect for the great phenomena when there is absolutely no
spiritual connection between them? --- \textsc{Carlyle} has here seen
more rightly than Nietzsche when he \textit{(Past and Present} I, 5), on the occasion
of the slogan of ``the aristocracy of talent,'' remarks: ``The
true eye for talent presupposes the true respect for it, and
--- O Heavens! --- presupposes so many things!'' These many
things Nietzsche calmly lets lie, absorbed as he is in ``the pathos
of distance.''

% --- p. 138 ---
\opage{138}\sectionhead{E. \textit{The Slave Revolt in Morality.}}

According to Nietzsche the social dualism is not only the right
and desirable thing, but also the natural. It finds its explanation
in the historical origin of morality \textit{(Genealogie der Moral).}
Moral concepts and names derive from the ruling,
the powerful, and the happy, and express their feeling of happiness
and self-feeling in contrast to the weak and the powerless. The ``good''
were originally the noble, the mighty, the high. ``Good'' means
``of the first rank''; by contrast it has from the beginning nothing at all
to do with unegoistic actions. The higher culture often begins
with a powerful race overwhelming a weaker one, whereby a
relation of distance, by whose help the advantages can be felt, comes
into being. The political concept of precedence becomes a mental concept of precedence.


An eager study of Theognis, which Nietzsche pursued in his early youth,
has left its traces here. In Theognis, a fierce aristocrat
who was driven out of his native city of Megara, ``the
good'' means the men of the aristocracy and ``the bad'' the men of the people. Moreover
it is the migration of peoples, the Renaissance, and Napoleon that serve
him as historical examples. Possibly impressions of childhood have also
had a part. His native town of Naumburg was a bureaucratic
center in the province, and the self-feeling of the circles of privy councillors
has been the first kind of ``nobility'' that the future
apostle of the aristocracy of spirit came to know. It is characteristic in any case
that he considers and examines the relation of distance
only from the one side and does not think at all of
how the same distance works on the other party,
--- what feelings of fear and awe, trust and admiration the
great and radiant figures may awaken. It is after all most likely
from this side of the relation that moral ideas have
their historical origin.\textsuperscript{44}) And it would also for Nietzsche
himself have been most consistent to assume this, since the ``noble''
natures in his conception are to rest in themselves,
live out of themselves, without reactive feelings. Then it is not
consistent to let them form relational concepts, or even
become aware of the distance. --- Moreover the ruling
% --- p. 139 ---
\opage{139}race must already before the conquest have formed a society, and within
this there will already have formed moral ideas
on the basis of a valuation of the qualities (e.g.
courage) which were of the greatest value for the society in question. ---

The natural valuation, which derives from the powerful
and happy, in whom the will to power can express itself freely, has
now, according to Nietzsche, been overthrown by a moral slave revolt,
from which our culture still suffers. It is due
especially to the Jews, --- this people that called itself the chosen,
but which Tacitus with greater right has characterized as born to
slavery; --- it was their prophets who made being rich
and mighty one with being ungodly and evil. ``The world'' became
a term of abuse, and poverty and suffering were praised. The continuation
and full development of this came in primitive Christianity. In the Revelation of
John the hatred of impotence is especially traced. War was declared
between Judea and Rome, between weakness and strength.
Freedom, pride, and spiritual selfishness were sacrificed. Paul in particular
carried this line of thought through. The apostles did not understand the
death of Jesus, which was precisely a great act of freedom; they made it
an atoning sacrifice, a glorification of suffering. --- In his later
writings Nietzsche piles up such a quantity of abuse in his
discussion of Christianity that his sister can explain
it only by the excessive use of chloral in the last year before
the catastrophe.

According to Nietzsche, however, not merely one moral
slave revolt has taken place, but several. Before Christianity Buddhism
and Socratism are already such revolts, only of a more noble
kind. Later the Reformation is a slave revolt against the Catholic
Church's and the Renaissance's noble worldliness; --- yes,
free thought, the Revolution, democracy, and --- natural science
(which is a democratization of nature by its principle
of universal lawfulness!) are also such revolts. At an increasing
pace Nietzsche extends the concept ``moral slave revolt''; after
he has once got his eye on it, it becomes like a fixed idea,
which is applied to the most diverse, mutually conflicting endeavors.
--- If it is anywhere justified to find
% --- p. 140 ---
\opage{140}evidence of
Nietzsche's morbidity in his ideas, it is on
this point, where he saws off the branch on which he himself
sits. ---

But what Nietzsche has not seen is that ``life'' is stirring
richly and powerfully precisely in those whom he calls slaves. Primitive Christianity
is precisely what Nietzsche calls a Dionysian movement,
a spiritual breakthrough, partly in the form of ecstasy, an expression of
the great longing, an Excelsior! And a role similar to that which Delphi
played toward Bacchantism has been played by the hierarchy and the ecclesiastical
organization, damming the violent stream,
--- not always in Apollonian fashion, despite the nobility by which the
organized hierarchy is marked in Nietzsche's eyes. --- And analogous
remarks could be made with regard to the other ``slave revolts.'' ---


A grounding of his theory of the slave revolts Nietzsche has
found in the alleged observation that the happy are
better human beings than the unhappy. The happy live out of
themselves, unfold their activity involuntarily, as activity
determined from within. They need neither hatred nor distortion, since they are not
dependent. The activity of the unhappy is determined from without,
is of a reactive kind; envy, suspicion, hatred, and deceit
lie close to them, dependent as they are. The happy say Yes,
the unhappy No. --- It cannot of course be denied that when
the active or happy natures ``discharge their power'' in such a way
that it causes suffering around them, they find it easier to
forget than those on whom the discharge of power has fallen. But
it is inconsistent when contempt is to play so great a role
in the psychology of the happy; for it is after all a decidedly reactive
feeling. ---

Yet this grounding is not the definitive one. The decisive
grounding of the social dualism and of the theory
of the slave revolt Nietzsche finds --- strangely enough --- in
the principle of welfare.

% --- p. 141 ---
\opage{141}\sectionhead{F. \textit{The Philosophical Grounding of the Social Dualism.}}

Nietzsche often calls himself an immoralist and declares that he
wants to abolish all morality, and he has, as is well known, called one of his
books ``Hinsides Godt og Ondt.'' It must nevertheless be clear from what has gone before
that it is only slave morality that he wants to eradicate.
And in one place he also says (in analogy with a well-known
epigram of Schiller) that he wants to eradicate morality --- ``aus Moralität.''
He wants to revalue --- but for that a standard is required, a principle that must necessarily express a basic value.
If we ask what this principle of valuation is, we find, to our
surprise, that it is the principle of welfare (das menschliche Gedeihen,
see the Preface to ``Geneal. d. Moral''). His fundamental question
is whether the earlier moral valuation has served the welfare
of the human race or not. His chief point of view is whether
an action or a personality expresses rising or declining
vital force. And he sets himself --- admittedly rather late
--- the task of ordering values scientifically in a graded series
determined by the strength and extent of force (Zahl- und
Maass-Scala der Kraft. \textit{Wille zur Macht.} Aphor. 353).
Unfortunately he has been so occupied with declaiming against
slave morality and against the herd that he has not got around to treating
this task in earnest and thereby developing in more definite form
what he calls his \textit{moralistic naturalism} (W. z.
M. Aphor. 192). What is required is a Dionysian valuation,
--- an ethics of vital force and joy of life. All the way from that height of
joy where man feels himself through and through a deified
form of and self-justification of nature
down to the joy of healthy peasants and healthy half-human animals
--- this whole long, immense scale of light and color of happiness the
Greeks called, not without the grateful awe that accompanies
initiation into a mystery, Dionysus! (W. z. M. Aphor.
482). The point is thus to bring about the happiness that corresponds
to the given stage of life. The ethics of welfare, which Nietzsche so often
scorns, can --- certain ``noble airs'' that conceal unsolved problems
set aside --- accept this conception, especially since
it is expressly declared that it is not a question of any individual's
% --- p. 142 ---
\opage{142}or any caste's happiness, but that the vigorous unfolding of life
and happiness of life which is required is required for the sake of the race
--- nur des Gesammt-Lebens willen. Only for the sake of the race as a whole
can value be attributed to the self-preservation of the individual; properly
the concept of the individual, of the single person, as
it has hitherto been understood, rests on an error. \textit{(Götzendämmerung.}
Aphor. 33).

The social dualism and the whole of master morality thus find in this
their final grounding. The masters and the overmen have then
not merely their end in themselves; the ultimate valuation of them
depends on whether they lie on the ascending line of the life of the race or
not, --- depends, that is, on their contribution to the development of human life.
--- Nietzsche's disciple and interpreter Peter Gast concedes
that the Master is in a certain sense a utilitarian; only he differs
from the usual utilitarianism, which takes account only of
the nearest utility. (Where this usual utilitarianism
is to be found, we are unfortunately not told. Utilitarianism demands
in all its forms that the effects of actions be
pursued as far as possible).

Now it was in itself quite possible that the welfare of the race
might require a sharp opposition between masters and slaves,
although it is difficult to see how such an opposition
could be maintained and carried through at the West European stage of culture.
But no serious attempt to demonstrate this has been
made. Stuart Mill in his day, proceeding from his democratic fundamental view,
made a far more serious attempt to show the necessity
of individual freedom, peculiarity, and greatness than
Nietzsche, proceeding from his aristocratic fundamental view, has made
to show the necessity of a slave class. And yet Nietzsche scorns
Mill and the other English thinkers, whose shoe-latchets he
--- considered as a strict, methodical inquirer --- is not
worthy to loose.

One thing stands fast in any case: it was conceivable that
a social dualism served the human race best; but it is
inconceivable, because it is self-contradictory, that one can at one and the same time
see the final end in single individuals or castes alone and
% --- p. 143 ---
\opage{143}yet in the same moment see the welfare of the whole race as the final
end.

But on the last, decisive points Nietzsche is not
a philosopher, but a poet. What he did not get as far as working through
with his thought has stood before him in a great vision of the future,
and to this we must turn to hear his last word.\textsuperscript{45})

\sectionhead{G. \textit{The Last Yes --- and Zarathustra's Death.}}

What drives Nietzsche on to a thought or a vision
that stands in such sharp contradiction to what he previously or
simultaneously maintains as a polemicist is a great contrast effect,
a mighty urge to make a complete end of the conception of life
which he had entered upon in his early youth under the influence of Schopenhauer
and Wagner. As a polemicist he reacts
against his contemporaries; in his visions he reacts against himself,
and he cannot get this reaction expressed strongly enough. The form of
thought does not suffice here; but in return a lofty poetry,
borne by mighty mood, is at his disposal. It is
the fundamental thought of his youthful work on the birth of tragedy that he
now returns to. Life is to be glorified. The urge for life expresses
itself as will to power, and this will can and shall be so
strong and courageous that it can even choose pain --- and choose
it not merely as means, --- choose to repeat life, the very same
life it has once come to know, with its toil and
pain.

Zarathustra, the prophet, has gathered in his mountain cave the
higher human beings, those who have suffered under the misjudgment and persecution
of the democratic rabble. He explains to them that they
will always have it worse and harder; always more of them will
perish. For it is not the subsistence of man that is
at stake, but the genesis of the overman. What we can
love in human beings is that they are a transition and a going under. There
are such great, distant, and high forces at work that it is no
wonder if the fragile vessel in which they develop is shattered.
Pity for these higher human beings is Zarathustra's
last sin. \textit{(Also sprach Zarathustra.} IV). ---

% --- p. 144 ---
\opage{144}What, then, is the overman? It is a form of existence that
stands in relation to man as man to the ape. Its genesis
is the meaning of the earth's existence (der Sinn der Erde). It
is a new type of life that is to be realized, and by the word overman
the concept or symbol of this type is expressed. \textit{(Wille
zur Macht.} Aphor. 390). The realization of this type humanity will,
when the great noon hour comes in the middle of the way of development,
see as its hope and as its task. The goal
thus comes after all to lie in the future and in a distant future,
which Nietzsche previously protested against. He no longer assumes
that history has reached its goal in the great ones of the past; that
is not enough for him. But what is especially of importance: the hints
that Nietzsche in his more positive developments gives of
the new type of life point in another direction than his depictions
of master morality. It is no longer mere force, but
magnanimous love of humanity that he demands. Even in Christianity
he now finds too little love, in that it cried woe
upon those who laugh. --- As soon as Nietzsche is positively to develop his
ideal and keeps the reactive moods away, the concept of the
overman undergoes a displacement. This is seen first and foremost
in Zarathustra himself. When he has got the higher human beings
to enter into his thought, he senses that the great noon
is near. Now he cannot remain in his lofty mountain solitude.
He needs, like the sun, to let his light shine for the world.
``You high star,'' he says to the sun, ``you great eye of happiness,
what would all your happiness be, if you had no one
to shine for!'' (Also sprach Zarathustra. IV. Das Zeichen). ``Once
more I will go down to human beings. Among them I will perish,
dying I will give them my richest gift. That I
have learned from the sun, when it goes down, the immeasurably rich:
gold it then pours out into the sea from its inexhaustible wealth, so
that the poorest fisherman comes to row with golden oar. That
I once saw myself, and I could not get enough of tears over it!''
(Also spr. Zar. III. Von alten und neuen Tafeln). Now Zarathustra
hates his own hatred. ``I have become one who blesses
and says Yes! Long I fought and was a warrior, so that
I might get my hands free to bless!'' (ibid. Vor
% --- p. 145 ---
\opage{145}Sonnenaufgang). --- Here the will to power expresses itself undeniably in a
different way than where contempt for the herd has the word. And it
is the very urge for life, the urge to drive pessimism out of
all corners, that leads Nietzsche to this summit. It is
the consistent unfolding of Nietzsche's fundamental thought that is here
fulfilled. ---

Now that Zarathustra has gained so clear a view of the goal and
of his proper task, he strikes up his midnight song. This
song (Nietzsche calls it ``the loneliest song that was ever
written'') Nietzsche had written one night on a loggia
on the Piazza Barbarina in Rome; in his poem he sets it
at the midnight hour before Zarathustra's cave. In this strophe
a mood of wholeness is struck and a comprehensive
thought is hinted at, which he otherwise was unable to hold fast and work out.
He has here come as near as possible to a great, concentrated
expression of enthusiastic resignation and hard-won hope. ---
I make no attempt to translate it, but render the thoughts
as well as I can.

``Midnight proclaims: The world is deep, far deeper than
the day thinks! --- Deep is the world's woe. But pleasure is deeper
still than heart's woe. Woe means perishing. But pleasure wants
eternity, wants deep, deep eternity!''

Through pain, sorrow, doubt, and contempt Nietzsche had
worked his way into a deep faith in life. Here he expresses
in poetic form the old biological theory of the feeling of pleasure
as expression of the force and progress of life. And he sees in every
stirring of pleasure a will to the continuance and continuation of life.


But this last thought, the will to the continuance of life, assumes
in him a strange form. It becomes one with a will
to repeat life. For to Nietzsche it had come to
stand as a necessity that what happens must come again
in the same way in the future. He proceeds from the view that the world
consists of a definite number of elements, so that only a definite,
finite number of combinations becomes possible. When it is exhausted,
recurrence must come. The course of the world is a cycle which has already
repeated itself infinitely often and likewise will repeat itself
% --- p. 146 ---
\opage{146}in the future.\textsuperscript{46}) When this thought dawned on Nietzsche, it seized
him with horror. He regarded it as a natural-scientific
necessity, and it put his optimism to a
hard test. Not only was pain now to be chosen and accepted,
but all that had been suffered was to be suffered again and again.
The greater was the jubilation of victory when he overcame the fear of recurrence,
and in this overcoming found the highest degree of
affirmation of life that was at all possible. In an interesting
way Nietzsche here meets Kierkegaard, for whom too
recurrence was a touchstone of the strength and
seriousness of a view of life. Kierkegaard's saying that ``he who wills recurrence, he
is a man,'' is also one of Nietzsche's fundamental propositions.
--- This thought is the fundamental thought of the Zarathustra book, where it
is, however, often obscured by aphorisms that stand in less clear
connection with it forcing themselves forward. Had Nietzsche
been able to work out the end of the book, the idea of the book would
have come into its own better.

This thought of recurrence Zarathustra proclaims first
to the higher human beings whom he has gathered around him at his mountain cave.
After some resistance he gets them to enter into
it, so that even ``the ugliest man'' exclaims: ``War \textit{Das}
--- das Leben? --- Wohlan, noch einmal!'' After that it is that
Zarathustra descends to proclaim the conditions of life to the
great mass of human beings. He gathers human beings to a feast and gives
them new laws. An order of rank is established in which the real
values of life are made the basis. The struggle of the estates is now
happily over. It is expressly recognized as the task of the ruling
caste to win the deep, unconditional trust of the ruled.
Hatred of democratic leveling has had its time,
which is now past. Zarathustra then proceeds to his proper task.
First he proclaims the great hope of the coming of the overman,
which becomes possible through the carrying out of the new valuation.
And then comes the great, terrible moment when
he proclaims to human beings that everything will come again. Now
this thought can be borne, not only by himself, but also by
human beings. For when he asks them: ``Do you want all this once
more?'' --- all answer yes, and Zarathustra dies of joy
% --- p. 147 ---
\opage{147}over it. (See the drafts for the 5th and 6th parts of ``Also sprach Zarathustra.''
Werke XII p. 321 ff.). ---

In poetic form Nietzsche has here taken back the social dualism
and master morality. Already the fact that Zarathustra
descends to proclaim the highest truths to the multitude of human beings
shows this, besides its being implied in the fact that the struggle of the estates
is concluded. There is now a common goal for the race.

One cannot deny that it is a hard test
to which Nietzsche puts our faith in life. It is after all possible to have
a great faith in life even if one does not wish to commit again
the same follies or rascalities and to suffer the same
pains and sorrows again and again. But there is also no
reason to assume absolute recurrence. Experience does not show us
it, and existence, the more our knowledge deepens
in it, stands as inexhaustible; ever it reveals greater
multiplicity and more distant possibilities. A combative view of life
--- and every view of life that wants to face reality
must be a combative one --- must give a chief place
to the thought of the new and unknown, which sticks up
now as a danger that threatens, now as a task that spurs, now
as a hope that beckons. There is nothing that speaks for existence
being so narrow that a limited number of existences
should succeed one another through all eternity. Our faith in life
therefore need not pass through the purgatory through which Nietzsche in
Dionysian enthusiasm would have it go.

As he has otherwise let the mood of disharmony and contempt
be raised to the utmost, so in his concluding
thoughts he has let the reaction against all pessimism reach the
highest power. They are expansions of feeling which at every point
dominate his thought. Not through scientific treatment of problems,
but through the passionate manner and the often brilliant pathos
with which mutually conflicting points of view come to expression
in him and thereby are set forth clearly and sharply, he belongs
in the history of philosophy. His significance is symptomatic.
He has let us become witnesses to an inner drama in a
soul that felt the tendencies of the time and of life deeply and energetically.

% --- p. 148 ---
\backmatter
\chapter*{NOTES.}
\addcontentsline{toc}{chapter}{Notes}
\markboth{Notes}{Notes}
\opage{148}1) P. 8. This fear of materialism can lead to questionable consequences,
since Wundt, in order to maintain the independence of psychology, quite naturally
seeks to limit ``the elements'' to the least possible, because it is
only in the elements, not in their connection, that a physiological
correlate is supposed to be necessary. Thus he has found it questionable
that I have called the new property which a sensation acquires when it is recognized
a quality (``the quality of familiarity''), since he thinks that I
am thereby led to put a physiological explanation in the place of a psychological one.
And he finds in this a sign of the times: ``Dieser Forscher ist durchaus
nicht den principiellen Vertretern der materialistischen Psychologie zuzurechnen.
Gerade dies aber scheint mir bezeichnend für das Ueberwuchern
hypothetischer und psychologisch unfruchtbarer physiologischer Erklärungsversuche
in der heutigen Literatur, dan selbst so ... unbefangene Psychologen
wie Høffding sich nicht ganz von denselben freihalten können .''
(Philos. Studien. X. p. 61). To this I have to reply that in my view
one must in the case of all psychical processes ask whether it is
possible to find a physiological explanation; --- this holds for connections
and comparisons, as well as for the single elements and qualities.
And according to the hypothesis of identity, which it is indeed also Wundt's opinion
to maintain, a physiological explanation does not exclude the psychological one;
every phenomenon ought, as far as possible, to be illuminated from both sides. I conceive
immediate recognition psychologically in analogy with association by similarity
and comparison, as a borderline case in relation to them,
namely as the most elementary form in which the relation of similarity makes
itself felt in consciousness. The psychological explanation Wundt gives
effaces the difference between the more elementary and the more composite
psychical processes, in that he assumes that free ideas cooperate even
in the simplest cases of recognition. --- A similar position to that taken toward
immediate recognition Wundt takes in his discussion of Weber's
% --- p. 149 ---
\opage{149}Law, in that he distinguishes between the sensations themselves and their ``comparison.''
\textit{(Physiologische Psychologie.} 5th ed. I. p. 541). It is a great
question whether this distinction is tenable as regards the most elementary
processes of sensation. --- Of the literature on Wundt I would especially
refer to \emph{Edmund König}: \textit{W. Wundt.} Seine Philosophie und Psychologie.
Stuttgart 1901. (Frommanns Klassiker der Philosophie). \emph{Allan}
\emph{Vannérus}: \textit{Vid studiet af Wundts psykologi.} (Stockholm 1896), I have
also derived benefit from an unpublished seminar paper on Wundt as a psychologist
by my young friend Mag. art. \emph{Aug. Bjarnason}.

2) P. 9. Cf. e.g. \emph{Wundt}: \textit{System der Philosophie.} (1st ed.) p.
533: ``Since the growth of spiritual energy consists only in the perfecting
of the organic forms in respect of \textit{quality}, the quantitative determination
of the physico-chemical energies remains wholly untouched by it .... In all
progress in organic development the quantity of the forces of nature is increased just as
little as the quantity of matter. But in value the forces of nature and their
substrates have grown immeasurably through the development of organic life. Indeed, determinations of value
have in general only become possible, and at the same time necessary,
through the arising of purposive acts of will and of the ideas
and feelings connected with them. Cf. \textit{Physiol. Psychol.} 5 III. p. 781.

3) P. 9. The actual concept of the soul has been asserted especially by
Fichte, but already lies in Kant's concept of consciousness as synthesis.
Among us it has been maintained by F. C. Sibbern and by me. (See especially my \textit{Psykologiske
Undersøgelser.} p. 67---81).

4) P. 12. Wundt thinks that I, in my \textit{Psykologiske Undersøgelser} (p.
60), have misunderstood his statements about apperception in the third edition
of ``Physiologische Psychologie,'' in that I have concluded from the claim that the difference between association
and apperception is to rest on an abstraction that
he does not assume them to be different processes; all that is supposed to lie in it is
``dass sie thatsächlich \textit{immer zusammen} vorkommen''. (Philos. Studien. VII
p. 229 f.), But in the later published fourth edition of ``Physiologische
Psychologie'' it says (II. p. 447) of the simplest connections of ideas:
``Diese Vorgänge unterscheiden sich durch die völlige Einflusslosigkeit
des Willens auf die Art ihres Eintritts auf das bestimmteste von
den nachher zu erörternden apperceptiven Verbindungen der Vorstellungen'',
--- and further (p. 479): ``Die Association bringt die Vorstellungen nur
in diejenigen Verbindungen, in die sie vermöge ihrer eigenen Beschaffenheit,
unbeeinflusst von jeder inneren Willensthätigkeit, sich ordnen.'' And
when Wundt still (also in the fifth edition of his main psychological work)
assumes a special center of apperception in the brain, this too seems to
imply that apperception and association need not always occur together.
--- The activity that expresses itself in apperception we become, according to
Wundt, immediately conscious of by a ``Thätigkeitsgefühl,'' which is common
to all apperception. (Physiol. Psychol.\textsuperscript{4} II. p. 266, 270, 279). In other places
(e.g. Philos. Stud. X. p. 109) he presents apperception, or
% --- p. 150 ---
\opage{150}the activity of will in general, as an object of immediate intuition.
(Cf. also Physiol. Psychol.\textsuperscript{4} II. p. 560: ``Der Begriff der Thätigkeit
stammt zunächst nur aus unseren eigenen Willenshandlungen und wurde
erst von ihnen auf äussere bewegte Gegenstände übertragen''). In his earlier
presentations Wundt placed the consciousness of activity in close connection
with the immediate sensation of innervation, which he later
rejected. --- Of the difficulties that are bound up with every
kind of consciousness of activity Wundt strangely does not speak. Cf. on
this question my Psykologi VII B. 4 and Psykologiske Undersøgelser
p. 68---81.

5) P. 13. Cf. Wundt's essay \textit{Ueber die Definition der Psychologie}
(Philos. Studien. XII) p. 517. \emph{Friedrich Paulsen}: \textit{Einleitung
in die Philosophie.} Berlin 1892. p. 116---132 (Intellektualistische und voluntaristische
Psychologie).

6) P. 14. In my presentation of psychology I already in the first
edition (1882) set out to demonstrate a psychical activity in sense-perception, association of ideas,
pleasure and displeasure as well as in what is specially called
will, and I found in the will the most proper and innermost nature of mental life
(IV, 7 e). But if I could now rework it from the ground up, I would
also in the form of presentation make it a psychology of will.

7) P. 19. Characteristic statements of \emph{Wundt} on metaphysics are found
in  \textit{Philos. Studien} XIII. p. 80; 428. \textit{Einleitung in die Philosophie.} p. 348---352.

8) P. 20. Wundt's development of this point is found in his \textit{System der
Philosophie} IV, 1, 3. (1st ed. p. 361---367). I do not find his line of thought
quite clear, but have presented it according to my interpretation of it.

9) P. 23. Cf. on this whole question my \textit{Psykologi.} p. 66.
\textit{Religionsfilosofi.} p. 66–86. \textit{Filosofiske Problemer} (University program 1902).
p. 62---65. --- In his excellent memorial address on Fechner
Wundt rightly characterizes Fechner's speculation as built on analogy, in contrast
to the scientific method, which requires induction. The relation
between analogy and induction is according to Wundt that the former limits
itself to a small number of marks, whereas the latter requires as many
as possible. \textit{(Gustav Theodor Fechner.} Rede zur Feier seines hundertjahrigen
Geburtstages. Leipzig 1901. p. 70---73). But what in the case of analogy
must first and foremost be taken into account is the starting point (what
Wundt calls ``Angriffspunkt'') or the point of similarity. Analogy is a similarity of relations,
and the decisive thing will be the degree to which this similarity of relations
can be regarded as a sign of a more deeply lying identity. In metaphysics
and religion existence as a whole is interpreted by the idea of a
single one of its parts or sides, and analogy must here always remain
imperfect.

10) P. 24. Cf. my \textit{Religionsfilosofi.} p. 217–246; 334–339.

11) P. 25. By ``the heterogony of ends'' Wundt understands the same thing
% --- p. 151 ---
\opage{151}that I call objective shifting of value (see my \textit{Etik} XIII, 4). He seems
to assume that the shifting of motive (the subjective shifting of value) always
presupposes that there is more in the effect than in the intention. This is, however,
not always necessary, since partly the entering into the result of the action, partly
other accompanying circumstances can make possible the arising of new motives. --- In
a purely psychological form the shifting of motive was already described by Spinoza,
Hartley, and James Mill. In the philosophy of history Hegel laid great
weight on ``the metamorphosis of ends.''

12) P. 26. Cf. my critical remarks on Wundt's ethics in
\textit{Etiske Undersøgelser.} p. 28 (where I have shown that Wundt in his polemic
against the morality of welfare forgets his own doctrine of ``the metamorphosis of ends'')
and in \textit{Etik} VIII, 4; XIII, 4.

13) P. 31. On the history of Italian philosophy before Ardigò I refer
to \emph{Louis Ferri}: \textit{Histoire de la philosophie en Italie an 19e siècle.}
(Paris 1869). \emph{Alfred Espinas}: \textit{La philosophie experimentale en Italie.}
(Paris 1880). In Ueberweg's Handbook there is a survey written by \emph{Luigi Credaro}.
With regard to the condemnation of criticism cf.
\emph{Rudolf Eucken}: \textit{Thomas von Aqvino und Kant,} ein Kamp zweier Welten.
Berlin 1901.

14) P. 36. These last remarks on the equivalence occur in
the work \textit{L'unità della coscienza} (p. 410---413), which is far later than the essay
\textit{La formazione naturale.} They do not seem to me quite accurate,
for when it is admitted that besides the inner dispositions outer conditions must
come in, in order that a ``distinzione'' may occur, then the
special forms do not really have their equivalents in the Indistinto
alone, but in its dispositions + the outer conditions. This now lies
really also in Ardigò's doctrine, when he declares every Indistinto to be
limited and the difference between indistinto and distinto to be relative.

15) P. 36. Quando it pensiero perdi di vista l'infinito, fissandosi nel
distinto finito, esso infinito lo assiste inosservato, e costituisce la stessa
forza della logica del suo discorso. \textit{Form. nat.} p. 136. (Opere filos. II).

16) P. 37. Ardigò is clear that validity does not rest on the thought's
likeness to the object that is thought, but on its own inner nature.
\textit{(La psicologia come scienza positiva.} p. 228) (Opere filos. I). But thereby
validity is not explained. In another place (see \textit{L'unità della coscienza.}
p. 441) (Opere filos. VII) he calls truth a property of thought
(una qualità del peusiero). But when thought is a phenomenon alongside
other phenomena, what has that property of thought to do
with these? --- A certain dogmatism expresses itself here in Ardigò.

17) P. 38. In a later edition of \textit{La psicologia} (1882) Ardigò adds
the following note to the discussion of ``the psychophysical reality'': Questa espressione,
sostanza psicofisica, non è altro che l'Indistinto naturale precedente
e sottostante ai due fenomeni distinti del mondo della psiche e di quella
% --- p. 152 ---
\opage{152}della materia. (p. 387). --- The concept Indistinto was introduced only in the essay
La formazione naturale (1877).

18) P. 42. Ardigò polemicizes here especially against Mamiani. See the essay
\textit{La religione di T. Mamiani} in Opere philos. Vol. II. ---

19) P. 43. \emph{William James} has indeed recently in warm words declared
his adherence to the old English school and declared that he builds
on its ground. (Cf. on this my essay \textit{Om nogle religionsfilosofiske
Arbejder fra den nyeste Tid.} University program 1903. p. 71). But
this adherence concerns only certain general principles, which are now
recognized by most who philosophize, and which, moreover, appear in James himself
with modifications that could not be derived from the principles of the older
English school itself.

20) P. 46. Besides the writings cited Bradley has also published a
whole series of essays in the journal \textit{Mind.} In the following presentation
some of these essays have also been used.

21) P. 53. Bradley stands in decided opposition to voluntarism. It
follows from his whole standpoint that he cannot find any solution of his riddles
in an appeal to will. ``What we know as will presupposes
a relation and a process, and presupposes at the same time an unreconciled disharmony
of elements. The same holds of energy or activity, or
anything similar.'' \textit{(Appearance and Reality.} p. 483). Bradley proceeds here
from the assumption that the concept on which we build our world-view must itself
be concluded and not lead to new questions. Like Goethe, he makes
the ``primal phenomenon'' into something that requires no explanation. (Cf. my
\textit{Filosofiske Problemer.} p. 56 f.). That will is a concept that sets new
questions in motion, in this Bradley is undoubtedly right. Schopenhauer,
the progenitor of voluntarism in the 19th century, overlooked it. But
does not the same hold of all concepts that one will try to set up as
concluding? And does it not hold first and foremost of the concept of experience,
which Bradley wants to take as foundation? --- I disregard here the purely psychological
side of the matter. Bradley has not yet published his psychology.
Yet it is seen from his essays in ``Mind'' \textit{(On active attention.}
1902. --- \textit{The definition of will.} 1902---1903) that his psychology has a
predominantly intellectualistic character.

In face of James Ward, who had criticized Bradley from the standpoint of voluntarism,
A. E. \emph{Taylor}, a thinker who stands close to Bradley, remarks
that it must nonetheless be clear that activity cannot be an expression of the
Highest or the Absolute, since activity always presupposes a resistance.
(Mind. 1900. p. 258). (Cf. my Religionsfilosofi p. 64; 81). In face of
this objection there is surely no other way out than the one that philosophical
theists have often seized, to assume the Deity to be limited. Thus
\emph{Rashdall} in the essay \textit{Personality, human and divine} (in ``Personal
Idealism''). The problem of evil has especially led several thinkers to this
result. Then one sees indeed with decided clarity how the
% --- p. 153 ---
\opage{153}problems keep posing themselves anew. ``The Absolute'' is according to Rashdall
not God, but ``a society that comprises God and all spirits.'' --- It
evidently does no good to carry the treatment of problems over to the mythological
domain.

22) P. 59. Certain traits for the characterization of Taine I have taken from
\emph{Boutmy} \textit{(Taine, Renan, Schérer.} Paris 1900) and from \emph{Albert Sorel's}
memorial address on Taine at his reception into the Akadémie française (1895).

23) P. 63. \emph{Charles Renouvier}: \textit{Esquisse d'une classification systematique
des doctrines philosophiques.} Paris 1886. II. p. 395. --- Cf.
moreover on Renan G. \emph{Séailles}: Ernest Renan. Essai de biographie
psychologique. Paris 1894. Mme \emph{Darmstetter}: \textit{La vie d'Ernest Renan.}
Paris  1898. \emph{Ed. Platzhoff}: \textit{Ernest Renan.} Ein Lebensbild. Leipzig
1900.

24) P. 65. Cf. my \textit{Psykologiske Undersøgelser} (1889), 76. ---
\textit{Psykologi} VII B, 4.

25) P. 71. Renouvier got this thought from his early deceased
friend \emph{Jules Lequier}, to whom he often refers. In the essay \textit{Doute
ow croyance} (L'année philosophique 1896) p. 44---51 he quotes a longer
statement by him. --- The whole conception recalls \emph{Fichte's} philosophy.
The conflict between dogmatism and idealism must indeed according to Fichte be decided
by a choice, and the choice will depend on what kind of person one
is. On this choice see especially Fichte's writing \textit{Die Bestimmung des Mensches.}
(Berlin 1800). p. 186; 193.

26) P. 76. See more closely on this my \textit{Psykologi} V D, and \textit{Filosofiske Problemer.}
p. 60; 62. --- After the text had gone to the printer, I receive news
that Renouvier has passed away.

27) P. 83. See the Introduction to Maxwell's essay on Faraday's
lines of force (German translation by Boltzmann in Ostwald's Klassiker der
Naturwissenschaft Nr. 69 p. 4). The idea that the application of number theory to
natural phenomena rests on analogy was subsequently carried out by \emph{Helmholtz} in his
essay in the Festschrift for Zeller (1887) \textit{Zählen und Messen, erkenntnisstheoretisch
betrachtet,} and by \emph{Ernst Mach} in his work \textit{Die Principien der
Wärmelehre, historisch-critisch erläutert} (1896). --- How occupied Maxwell
was with the concept of analogy is seen from the fact that from the same time in which
he wrote the essay on Faraday's theories of force there is a draft
from his hand, which discusses the question whether there are real analogies
in nature. (It is printed in \emph{Campbell and Garnett}: \textit{The
Life of James Clerk Maxwell.} London 1882. p. 235---244). It does not seem
to me to be quite clear what the concluding result concerns. It
is really only a sketch that hints at various analogies that could
acquire significance for knowledge. A testimony to the influence of the critical philosophy
on Maxwell's line of thought is his remark on the analogy
between ground and cause: Reasons, when spoken of with relation to
% --- p. 154 ---
\opage{154}objects, get the name of causes, which are reasons, analogically referred
to objects instead of thoughts. (p. 238).

28) P. 86. At a meeting of the University College Christian Association
in May 1903 Lord Kelvin declared that, as regards the origin of life,
one could not say that science either asserted or denied a
creative force; no, science \textit{positively asserted} a creative and guiding
force, since the purely physical forces could not explain life. This statement
occasioned a lively discussion, in which Kelvin's assertion was
decidedly opposed by Thiselton-Dyer, Karl Pearson, Ray-Lankester, and others.
\textit{(Times weekly.} 8. 15. and 22. May). It can only be the great physicist's
authority in his own domain that has aroused so great attention
for this uncritical statement.

29) P. 88. Kirchhoff's conception is presented by \emph{Boltzmann} \textit{(Gustav Robert
Kirchhoff.} Festrede 1888. p. 25) thus: ``The point is not to
form bold hypotheses about the essence of matter and from the motions of the molecules
to guess the motions of bodies, but to form equations
that without hypotheses as completely as possible and with quantitative
exactness correspond to the world of experience, unconcerned about the essence of things
and forces. Indeed, in his book on mechanics Kirchhoff wants even to banish all
metaphysical concepts, such as force or cause of motion; he seeks merely
the equations that with as great exactness as possible correspond to the observed
motions.''

30) P. 95. Cf. \textit{Filosofiske Problemer.} p. 46.

31) P. 95. At the philosophical congress in Paris in 1900 it was stated, during
the discussion of a paper that adhered to Mach's and Kirchhoff's
standpoint, that when in science one puts acceleration in
the place of force, it is because one wants to put visual symbols in the place
of muscular ones, which again has its ground in the fact that the former are exactly
measurable.

32) P. 95. Cf. on this dynamic-symbolic concept of truth \textit{Filosofiske
Problemer.} p. 37 f.

33) P. 102. Cf. \textit{Filosofiske Problemer} p. 15; 18. --- Avenarius'
disciple Fr. \emph{Carstanjen} has given a summary of the first volume of ``Kritik
der reinen Erfahrung'' in his writing \textit{Richard Avenarius' biomechanische Grundlegung
der neuen allgemeinen Erkenntnisstheorie} (München 1894). On several
points the presentation of the main work becomes clearer through Carstanjen's interpretation
--- but the clarity is attained only by the use of psychological
examples, which are taken from the second part of the work, which really ought not
to have been used at all, since the ``independent'' series ought after all to illuminate
the ``dependent,'' not the reverse. Besides, Carstanjev himself adds several ``symptoms''
which I have used in the text. --- The oral
remark of Avenarius cited in the text is found in the essay \textit{Richard Avenarius. Ein
Nachruf} by Fr. Carstanjen (Vierteljahrschrift für wissenschaftliche Philosophie
XX). --- It is unjustified when \emph{Wundt} in his otherwise excellent
% --- p. 155 ---
\opage{155}critique of Avenarius \textit{(Ueber naiven und kritischen Realismus.} Philos. Studien.
XIII. p. 331) calls his doctrine materialism.

34) P. 106. Cf. my \textit{Religionsfilosofi.} p. 240.

35) P. 106. Whereas Avenarius has given a natural history of problems,
I set myself in my essay \textit{Filosofiske Problemer} the task of finding
what is characteristic of the content of the various main problems. My investigation
confirms Avenarius' insofar as it turns out to be the relation
between discontinuity and continuity that recurs within
every so-called problem. An essential difference is that I stress the inexhaustibility
of experience (of existence), whereas Avenarius operates with
the idea of a final state. I therefore stress more strongly the
abiding irrationality in our intellectual relation to existence.

36) P. 113. From a standpoint similar to the present one Guyau's theory
has been criticized in the acute essay by \emph{Ch. Christophle}: \textit{Le
principe de la vie comme mobile moral selon J. M. Guyau.} Étude critique.
(1901). (Offprint from Revue de Metaphysique et de Morale).

37) P. 118. Cf. on this my \textit{Religionsfilosofi.} p. 103. --- \emph{Fouillée}
(\textit{Le mouvemeut idealiste.} p. XXII) calls Guyau's ``Irreligion'' a religion of the future,
although Guyau himself would not have had it so called. F. \emph{Buisson}
\textit{(L'èducation morale et l'èducation religieuse,} in the collection ``Questione de Morale.''
Paris 1900. p. 329) thinks that it above all has a claim to the name
religion. --- The name naturally does not matter much. --- A good
characterization of Guyau and a survey of his ideas has his stepfather
Fouillée given in the writing \textit{La morale, l'art et la religion d'après Guyau.}
Paris 1889. The American philosopher of religion \emph{Josiah Royce} has also
given a very sympathetic characterization of Guyau (in an essay in
his \textit{Studies in Good and Evil.} New York 1898. p. 379---384).

38) P. 122. In his book \textit{Richard Wagner, poète et peuseur} (3. éd. p.
301; 429---432) \emph{Henri Lichtenberger} has sought to show that Nietzsche
does Wagner an injustice when he thinks that the latter has ``converted.'' There is
--- despite various nuances --- a continuity in Wagner's view of life.
That he later places Christianity above Buddhism, whereas
he earlier placed them side by side, is essentially an expression
of the fact that he --- just as in his own way Nietzsche himself --- reacted
against pessimism.
He never became confessional in his religious conception.

39) P. 122. Nietzsche's sister, Frau \emph{Elisabeth Förster-Nietzsche}
has begun the publication of a biography, of which so far two volumes (\textit{Das
Leben Friedrich Nietzsches.} I and II, 1. Leipzig 1895---1897) have appeared.
They go to 1880. In addition come Nietzsche's \textit{Briefe} (2. Bde. 1900---1902). In
\emph{Crusius'} biography of Erwin Rohde many contributions are found toward the elucidation
and understanding of Nietzsche; likewise in Paul Deussen's beautiful
book \textit{Erinnerungen an Nietzsche} (Leipzig 1901). Nietzsche's motherly friend
\emph{Malwida von Meysenbug} has in her book \textit{Individualitäten} (1901)
p. 1---41 given an interesting depiction of Nietzsche's personality.

% --- p. 156 ---
\opage{156}40) P. 133. Against Socrates Nietzsche is in general severe and unjust.
He cannot find abuse enough against him. In ``Geburt der Tragödie''
it is, besides intellectualism, the optimistic faith in life
that is reproached to him. Later \textit{(Fröliche Wissenschaft.} Aphor. 340) he must
be spoken ill of for the pessimism that for Nietzsche lies in his last words
in Plato, the request to sacrifice a cock to Asclepius. He has regarded
life as a disease! --- Nietzsche follows here the traditional, Neoplatonic
interpretation of the well-known remark. Professor \emph{Heiberg} has
shown that this can be interpreted in a far simpler, less mystical way,
when one looks at the whole context in which it occurs. (See
Oversigten over det danske Videnskabernes Selskabs Forhandlinger 1902.
\emph{J. L. Heiberg}: \textit{Sokrates' sidste Ord).}

41) P. 134. In his last writing \textit{Wille zur Macht} Nietzsche
himself designates his view as aristocratism (Werke XV. p. 427). In the passage
cited Nietzsche seems to polemicize against Guyau and Fouillée, who
lay so great weight on the sociological consideration and its significance
for the view of life.

42) P. 136. A change has here occurred in Nietzsche's conception,
for in \textit{Unzeitgemeine Betrachtungen} (Drittes Stück § 5) it is the task of
\textit{everyone} to work toward the genesis of the great human being. Here the dualism
is not yet so sharp, and the problem not so intricate.

43) P. 137. Nietzsche's copies of Guyau's books with his marginal notes
are found in the Nietzsche Archive in Weimar, whence \emph{Fouillée}
has obtained the most important of these in transcript. He has used them in an interesting
chapter in his book \textit{Nietzsche et l'Immoralisme} (1902). Livre III:
Les jugements de Nietzsche sur Guyau d'après des documents inédits.

44) P. 138. Cf. my \textit{Etik.} X, 3 (toward the end) on the significance of
models --- in analogy with the significance of variations for biology.

45) P. 143. I do not go into the hints in an epistemological
and metaphysical direction that Nietzsche has given, especially in the third book of
\textit{Wille zur Macht.} In his conception of the principles of knowledge he recalls
the economic-biological theory of knowledge, and even more decidedly he recalls
in his metaphysics of will Schopenhauer's philosophy (only that he speaks
of ``the will to power'' instead of ``the will to life'').

46) P. 146. The thought of the recurrence of all things is not new. Nietzsche
has probably first come to know it through his classical studies.
It is an ancient thought, and it is closely connected with the idea
of the world as a limited whole. When elements and forces at the same time
are assumed to be imperishable, a rhythmic recurrence is necessary. With the Stoics
such a recurrence came after every conflagration of the world, both for
gods and human beings. Everything then proceeds in the same way; Socrates
marries Xanthippe again, etc. It was for the Stoics (see especially \emph{Marcus}
\emph{Aurelius}. Comm. XI, 1) a support for their peace of mind that nothing new
came; neither will our descendants get to see anything essentially new, nor
% --- p. 157 ---
\opage{157}did our forefathers see another world than we. --- It has in more recent
times more often been expressed as an alleged consequence of the scientific
conception. Thus by Blanqui, Le Bon, Nàgeli, Guyau. (Cf. \emph{Fouillée}:
Nietzsche et l'Immoralisme. IV, 3: Le retour éternel). Also in
Dostojevski the thought comes forth --- as a temptation from the Devil himself!
(Cf. \emph{Merejkowsky}: \textit{Tolstoi et Dostojewsky.} Paris 1903. p. 300).
--- Recurrence is naturally necessary only when one presupposes the limitation of the world
and at the same time has a dogmatic trust in the absolute validity of the scientific
principles. If, on the other hand, it is impossible for us
to think any absolute limit for the world in time, in space, in energy, and
if those principles are only the guiding hypotheses for our inquiry,
the possibility of the arising of new elements and combinations is not excluded.
The law of relation leads us naturally in this direction and places us
thereby before the possibility of the irrationality of existence. Cf. on this
my \textit{Filosofiske Problemer.} Chap. II and III. Thereafter Goethe will most likely
be right when he says: ``Neu ist immer die Erfahrung.'' --- In Nietzsche
the doctrine of the necessity of recurrence is not consistent. For he
lays in his theory of knowledge great weight on the fact that the principle of identity and the
other principles are postulates of our will, expressing our wish
to bring nature into our power. It is we who put identity into nature.
``Man soll die Nöthigung, Begriffe, Gattungen, Formen, Zwecke,
Gesetze zu bilden (`eine Welt der identischen Fälle'), nicht so verstehen,
als ob wir damit die wahre Welt zu fixiren im Stande wären; sondern
als Nöthigung, uns eine Welt zurecht zu machen, bei der unsre Existenz
ermöglicht wird: --- wir schaffen damit eine Welt, die berechenbar, vereinfacht,
verständlich u. s. w. für uns ist.'' \textit{(Der Wille zur Macht.} Aphor.
279). When Nietzsche assumes that there are not in reality identical
cases, he must consistently give up the belief in absolute recurrence. Experience
indeed also shows only an approximation to absolute identity and recurrence.
Nietzsche's metaphysics and his theory of knowledge do not agree.
He moreover polemicizes sometimes against natural science, because
it introduces equality before the law into nature!

\end{document}
